% Money and banking — Economics Upper Sixth — Cours signé OnBuch+
% Official MINESEC syllabus. Compiler avec : tectonic ECO03-money-and-banking.tex
\newcommand{\DOCMATIERE}{Economics}
\newcommand{\DOCNIVEAU}{Upper Sixth}
\newcommand{\DOCMODULE}{Upper Sixth — Economics}
\newcommand{\DOCLECON}{3}
\newcommand{\DOCTITRE}{Money and banking}
% OnBuch+ — préambule « cours signés » ENRICHI (compilé avec tectonic / XeLaTeX)
% Système visuel « Pop » d'OnBuch+ : fond crème (jamais blanc), bordures encre
% épaisses, ombres « dures » décalées, titres Archivo Black en capitales.
% Version MATHÉMATIQUES : graphes (pgfplots/tikz), boîtes pédagogiques
% (définition, propriété, méthode, attention, le savais-tu, exercice résolu,
% exercice, TP, bilan), page de garde et signature OnBuch+ ancrée en bas.
%
% Macros attendues dans main.tex AVANT \input{preamble} :
%   \DOCMATIERE  \DOCNIVEAU  \DOCMODULE  \DOCLECON (numéro)  \DOCTITRE
\documentclass[11pt]{article}

\usepackage{fontspec}
\usepackage[english]{babel}
\usepackage[a4paper,margin=2cm,top=3.4cm,bottom=2.8cm,headheight=1.4cm,footskip=1.3cm]{geometry}
\usepackage{amsmath,amssymb,amsfonts}
\usepackage{mathtools} % \xmapsto, \coloneqq... souvent utilisés par les modèles
\usepackage{cancel} % simplifications barrées (bilans de forces)
% \abs{z} (module, valeur absolue) et \norm{u} (norme), utilisés par les modèles
\providecommand{\abs}{}\let\abs\relax\DeclarePairedDelimiter{\abs}{\lvert}{\rvert}
\providecommand{\norm}{}\let\norm\relax\DeclarePairedDelimiter{\norm}{\lVert}{\rVert}
\usepackage[table]{xcolor} % [table] : fournit aussi \rowcolors (alternance de lignes)
\usepackage{pagecolor}
\usepackage{tikz}
\usetikzlibrary{arrows.meta,positioning,calc,shapes.geometric,shapes.misc,shapes.arrows,decorations.pathmorphing,decorations.pathreplacing,decorations.markings,patterns,shadows,mindmap,trees,fit,backgrounds,matrix,intersections,angles,quotes}
\usepackage{pgfplots}
\usepackage[version=4]{mhchem} % équations nucléaires \ce{^{235}_{92}U}
\usepackage[european,siunitx]{circuitikz} % schémas électriques
\pgfplotsset{compat=1.18}
\usepgfplotslibrary{fillbetween,groupplots} % aires entre courbes (calcul intégral)
\usepackage{enumitem}
\usepackage{fancyhdr}
% [most] charge toutes les bibliothèques tcolorbox (listings, minted, raster,
% documentation...) et gonfle inutilement la mémoire TeX sur un document long
% (~300 boîtes) : on ne charge que ce qui est réellement utilisé.
\usepackage[skins,breakable]{tcolorbox}
\usepackage{graphicx}
\usepackage{colortbl}
\usepackage{booktabs}
\usepackage{tabularx}
\usepackage{diagbox} % case d'en-tête diagonale des tableaux à double entrée
% Colonnes extensibles centrée / gauche / droite (C, L, R), souvent utilisées par les modèles
\newcolumntype{C}{>{\centering\arraybackslash}X}
\newcolumntype{L}{>{\raggedright\arraybackslash}X}
\newcolumntype{R}{>{\raggedleft\arraybackslash}X}
\usepackage{array}
\usepackage{multicol}
\usepackage{siunitx}
% parse-numbers=false : accepte aussi \SI{2,5 \times 10^{-3}}{...}
\sisetup{locale=UK,output-decimal-marker={.},group-minimum-digits=4,parse-numbers=false}
\DeclareSIUnit\ohm{\ensuremath{\symup{\Omega}}} % U+2126 absent de la police
\DeclareSIUnit\division{div} % réglages d'oscilloscope (ms/div, V/div)
\DeclareSIUnit\tr{tr} % tours (vitesse de rotation, ex. \SI{1500}{\tr\per\minute})
\usepackage{qrcode}
% \degree : commande fréquemment utilisée par les modèles pour un angle ou une
% température en dehors de \SI{}{} (non fournie par les paquets chargés).
\providecommand{\degree}{^{\circ}}
% \qed (fin de démonstration), utilisé par les modèles sans amsthm
\providecommand{\qed}{\hfill$\square$}
% \barre : double barre des valeurs interdites dans un tableau de variations (tkz-tab)
\providecommand{\barre}{\ensuremath{\|}}
% environnement proof (amsthm non chargé) utilisé par les modèles
\ifdefined\proof\else\newenvironment{proof}[1][Proof]{\par\noindent\textit{#1.}\ }{\qed\par}\fi
\usepackage{lastpage}
\usepackage{hyperref}
\hypersetup{hidelinks}

% ============================================================
% Palette « Pop » OnBuch+ — mêmes hex que la classe P (plus_ui.dart)
% ============================================================
\definecolor{poporange}{HTML}{F59321}
\definecolor{popdark}{HTML}{E07A0C}
\definecolor{popcream}{HTML}{FFF6E8}
\definecolor{popink}{HTML}{1C1714}
\definecolor{poppurple}{HTML}{7A5AE0}
\definecolor{popgreen}{HTML}{2E9E6A}
\definecolor{poppink}{HTML}{E0457B}
\definecolor{popblue}{HTML}{2F7DE1}
\definecolor{popgold}{HTML}{C98A12}
\definecolor{popmuted}{HTML}{8A7F76}
\definecolor{popline}{HTML}{EADFD2}
\definecolor{popgray}{HTML}{8A7F76}
\colorlet{popgrey}{popgray}
% Alias tolérants (le modèle nomme parfois les couleurs différemment)
\colorlet{popwhite}{white}
\colorlet{popblack}{popink}
\colorlet{popred}{poppink}
\colorlet{popyellow}{popgold}
% Teintes claires pour fonds de tableaux / zones
\colorlet{popblueL}{popblue!10!white}
\colorlet{poporangeL}{poporange!14!white}
\colorlet{popgreenL}{popgreen!12!white}
\colorlet{poppurpleL}{poppurple!12!white}
\colorlet{poppinkL}{poppink!12!white}
\colorlet{popgoldL}{popgold!14!white}

\pagecolor{popcream}

% ============================================================
% Typographie
% ============================================================
\defaultfontfeatures{Ligatures=TeX}
\setmainfont{PlusJakartaSans-Medium.ttf}[Path=../../../fonts/,BoldFont=PlusJakartaSans-Bold.ttf]
% Maths en sans-serif (Fira Math) avec chiffres et lettres droites de Plus
% Jakarta, pour que les formules se fondent dans le texte.
\usepackage[math-style=ISO,bold-style=ISO]{unicode-math}
\setmathfont{FiraMath-Regular.otf}[Path=../../../fonts/]
\setmathfont{PlusJakartaSans-Medium.ttf}[Path=../../../fonts/,range={up/num,up/{Latin,latin},\mathup}]
% (icomma retiré : décimales avec point en anglais)
% Caractères mathématiques Unicode absents des polices de texte (Plus Jakarta,
% Archivo Black) : rendus par leur équivalent mathématique, en texte comme en
% formule — sinon ils s'affichent en carré vide (ex. « Arithmétique dans ℤ »).
\usepackage{newunicodechar}
\newunicodechar{ℝ}{\ensuremath{\mathbb{R}}}
\newunicodechar{ℤ}{\ensuremath{\mathbb{Z}}}
\newunicodechar{ℂ}{\ensuremath{\mathbb{C}}}
\newunicodechar{ℕ}{\ensuremath{\mathbb{N}}}
\newunicodechar{ℚ}{\ensuremath{\mathbb{Q}}}
\newunicodechar{𝔻}{\ensuremath{\mathbb{D}}}
\newunicodechar{α}{\ensuremath{\alpha}}
\newunicodechar{β}{\ensuremath{\beta}}
\newunicodechar{γ}{\ensuremath{\gamma}}
\newunicodechar{δ}{\ensuremath{\delta}}
\newunicodechar{ε}{\ensuremath{\varepsilon}}
\newunicodechar{θ}{\ensuremath{\theta}}
\newunicodechar{λ}{\ensuremath{\lambda}}
\newunicodechar{μ}{\ensuremath{\mu}}
\newunicodechar{σ}{\ensuremath{\sigma}}
\newunicodechar{τ}{\ensuremath{\tau}}
\newunicodechar{φ}{\ensuremath{\varphi}}
\newunicodechar{ψ}{\ensuremath{\psi}}
\newunicodechar{ω}{\ensuremath{\omega}}
\newunicodechar{Σ}{\ensuremath{\Sigma}}
% majuscules produites par \MakeUppercase dans les titres (α-aminés -> Α)
\newunicodechar{Α}{\ensuremath{\alpha}}
\newunicodechar{Β}{\ensuremath{\beta}}
\newunicodechar{Γ}{\ensuremath{\Gamma}}
\newunicodechar{Δ}{\ensuremath{\Delta}}
\newunicodechar{Ε}{\ensuremath{\varepsilon}}
\newunicodechar{Θ}{\ensuremath{\Theta}}
\newunicodechar{Λ}{\ensuremath{\Lambda}}
\newunicodechar{Μ}{\ensuremath{\mu}}
\newunicodechar{Τ}{\ensuremath{\tau}}
\newunicodechar{Φ}{\ensuremath{\Phi}}
\newunicodechar{Ψ}{\ensuremath{\Psi}}
\newunicodechar{Ω}{\ensuremath{\Omega}}
\newunicodechar{↦}{\ensuremath{\mapsto}}
\newunicodechar{∈}{\ensuremath{\in}}
\newunicodechar{∉}{\ensuremath{\notin}}
\newunicodechar{∧}{\ensuremath{\wedge}}
\newunicodechar{⇔}{\ensuremath{\Leftrightarrow}}
\newunicodechar{⟺}{\ensuremath{\Longleftrightarrow}}
\newunicodechar{⇒}{\ensuremath{\Rightarrow}}
\newunicodechar{⟹}{\ensuremath{\Longrightarrow}}
\newunicodechar{∘}{\ensuremath{\circ}}
\newunicodechar{∪}{\ensuremath{\cup}}
\newunicodechar{∩}{\ensuremath{\cap}}
\newunicodechar{≡}{\ensuremath{\equiv}}
\newunicodechar{⊂}{\ensuremath{\subset}}
\newunicodechar{⊥}{\ensuremath{\perp}}
\newunicodechar{∃}{\ensuremath{\exists}}
\newunicodechar{∀}{\ensuremath{\forall}}
\newunicodechar{∛}{\ensuremath{\sqrt[3]{\,}}}
\newunicodechar{∜}{\ensuremath{\sqrt[4]{\,}}}
\newunicodechar{⃗}{\ensuremath{{}^{\to}}}
\newunicodechar{ⁿ}{\ifmmode^{n}\else\textsuperscript{\ensuremath{n}}\fi}
\newunicodechar{ˣ}{\ifmmode^{x}\else\textsuperscript{\ensuremath{x}}\fi}
\newunicodechar{ᵇ}{\ifmmode^{b}\else\textsuperscript{\ensuremath{b}}\fi}
\newunicodechar{ʳ}{\ifmmode^{r}\else\textsuperscript{\ensuremath{r}}\fi}
\newunicodechar{ᵘ}{\ifmmode^{u}\else\textsuperscript{\ensuremath{u}}\fi}
\newunicodechar{ʸ}{\ifmmode^{y}\else\textsuperscript{\ensuremath{y}}\fi}
\newunicodechar{ᵃ}{\ifmmode^{a}\else\textsuperscript{\ensuremath{a}}\fi}
\newunicodechar{ᵉ}{\ifmmode^{e}\else\textsuperscript{\ensuremath{e}}\fi}
\newunicodechar{ᵏ}{\ifmmode^{k}\else\textsuperscript{\ensuremath{k}}\fi}
\newunicodechar{ᵖ}{\ifmmode^{p}\else\textsuperscript{\ensuremath{p}}\fi}
\newunicodechar{ᴺ}{\ifmmode^{N}\else\textsuperscript{\ensuremath{N}}\fi}
\newunicodechar{ᶜ}{\ifmmode^{c}\else\textsuperscript{\ensuremath{c}}\fi}
\newunicodechar{ʰ}{\ifmmode^{h}\else\textsuperscript{\ensuremath{h}}\fi}
\newunicodechar{ᵠ}{\ifmmode^{q}\else\textsuperscript{\ensuremath{q}}\fi}
\newunicodechar{ₙ}{\ifmmode_{n}\else\textsubscript{\ensuremath{n}}\fi}
\newunicodechar{ₐ}{\ifmmode_{a}\else\textsubscript{\ensuremath{a}}\fi}
\newunicodechar{ᵢ}{\ifmmode_{i}\else\textsubscript{\ensuremath{i}}\fi}
\newunicodechar{ᵣ}{\ifmmode_{r}\else\textsubscript{\ensuremath{r}}\fi}
\newunicodechar{ₖ}{\ifmmode_{k}\else\textsubscript{\ensuremath{k}}\fi}
\newunicodechar{ᵦ}{\ifmmode_{\beta}\else\textsubscript{\ensuremath{\beta}}\fi}
\newunicodechar{ₓ}{\ifmmode_{x}\else\textsubscript{\ensuremath{x}}\fi}
\newfontfamily\popdisplay{ArchivoBlack-Regular.ttf}[Path=../../../fonts/]
\newfontfamily\poplogo{SpaceGrotesk-Bold.ttf}[Path=../../../fonts/]
\newfontfamily\popsemi{PlusJakartaSans-SemiBold.ttf}[Path=../../../fonts/]
\newfontfamily\popxbold{PlusJakartaSans-ExtraBold.ttf}[Path=../../../fonts/]
\newfontfamily\popxboldls{PlusJakartaSans-ExtraBold.ttf}[Path=../../../fonts/,LetterSpace=3.0]

\setlist[enumerate]{leftmargin=1.7em,itemsep=5pt,topsep=4pt}
\setlist[itemize]{leftmargin=1.7em,itemsep=4pt,topsep=4pt,label={\color{poporange}\textbullet}}
\setlength{\parindent}{0pt}
\setlength{\parskip}{5pt}
\renewcommand{\arraystretch}{1.25}

% pgfplots : style Pop par défaut
\pgfplotsset{
  every axis/.append style={axis line style={popink,line width=1pt},tick style={popink},
    grid style={popline,line width=0.5pt},label style={font=\small},tick label style={font=\footnotesize}},
  popcurve/.style={poporange,line width=1.8pt,no markers,smooth},
}
% Même style disponible dans un \draw TikZ simple (hors pgfplots)
\tikzset{popcurve/.style={poporange,line width=1.8pt}}
\tikzset{popgrid/.style={popline,very thin}}
\colorlet{popcurve}{poporange} % le nom du style est parfois utilisé comme couleur

% ============================================================
% Pill / squiggle
% ============================================================
\newcommand{\poppill}[3]{%
  \tikz[baseline=-0.55ex]\node[draw=popink,line width=1.2pt,rounded corners=2.6pt,%
    fill=#2,inner xsep=6.5pt,inner ysep=3pt]{%
    \popxboldls\fontsize{7}{7}\selectfont\color{#3}\MakeUppercase{#1}};%
}
\newcommand{\popsquiggle}[1]{%
  \tikz[baseline=-0.4ex]{
    \draw[#1,line width=2.6pt,line cap=round]
      plot [smooth] coordinates {(0,0.09) (0.34,0.01) (0.68,0.21) (1.02,0.01) (1.36,0.10)};
  }%
}
% Mot-clé surligné (marqueur orange sous le texte)
\newcommand{\cle}[1]{{\popxbold\color{popdark}#1}}

% ============================================================
% En-tête / pied
% ============================================================
\pagestyle{fancy}
\fancyhf{}
\renewcommand{\headrulewidth}{0pt}
\renewcommand{\footrulewidth}{0pt}
\lhead{\raisebox{-0.15ex}{\poplogo\fontsize{13}{13}\selectfont\color{popink}OnBuch\color{poporange}+}}
\rhead{\poppill{\DOCMATIERE\ \textbullet\ \DOCNIVEAU\ \textbullet\ Chapter \DOCLECON}{white}{popink}}
\lfoot{\popxbold\fontsize{7.6}{7.6}\selectfont\color{popmuted}\MakeUppercase{Signed course OnBuch+}\ \textbullet\ {\popsemi\DOCTITRE}}
\rfoot{\tikz[baseline=-0.6ex]\node[rounded corners=8.5pt,fill=popink,minimum height=17pt,minimum width=17pt,inner xsep=5pt,inner sep=0]{\popxbold\fontsize{8}{8}\selectfont\color{popcream}\thepage\,/\,\pageref*{LastPage}};}

% ============================================================
% Titres
% ============================================================
\newcommand{\chaptitre}[2]{%
  \begin{tcolorbox}[enhanced,colback=poporange,colframe=popink,boxrule=2.6pt,arc=11pt,
      drop shadow=popink,left=15pt,right=15pt,top=12pt,bottom=11pt,before skip=3pt,after skip=18pt]
    \poppill{#2}{popink}{popcream}\par\vspace{9pt}
    {\raggedright\hyphenpenalty=10000\exhyphenpenalty=10000\popdisplay\fontsize{22}{24}\selectfont\color{popcream}\MakeUppercase{#1}\par}
  \end{tcolorbox}
}
% Section numérotée : pastille orange avec le numéro + titre Archivo
\newcounter{popsec}
\newcommand{\coursec}[1]{%
  \stepcounter{popsec}\setcounter{popsub}{0}%
  \par\vspace{16pt}\noindent
  \begin{minipage}{\linewidth}
  \tikz[baseline=-0.7ex]\node[draw=popink,line width=1.6pt,fill=poporange,rounded corners=4pt,minimum size=22pt,inner sep=0]{\popdisplay\fontsize{12}{12}\selectfont\color{popcream}\thepopsec};%
  \hspace{8pt}{\popdisplay\fontsize{15}{17}\selectfont\color{popink}\MakeUppercase{#1}}\par
  \vspace{4pt}\noindent{\color{poporange}\rule{36pt}{3.6pt}}
  \end{minipage}\par\nopagebreak\vspace{8pt}
}
\newcounter{popsub}
\newcommand{\courssub}[1]{%
  \stepcounter{popsub}%
  \par\vspace{9pt}\noindent{\popxbold\fontsize{11.5}{13}\selectfont\color{popink}{\color{poporange}\thepopsec.\thepopsub}\hspace{6pt}#1}\par\nopagebreak\vspace{4pt}
}

% ============================================================
% Boîtes pédagogiques — même recette : fond blanc, bordure épaisse colorée,
% ombre dure encre, badge plein. Titre optionnel [#1].
% ============================================================
\tcbset{popbase/.style={breakable,enhanced,colback=white,boxrule=2.3pt,arc=9pt,
  shadow={3pt}{-3pt}{0pt}{popink},left=14pt,right=14pt,top=11pt,bottom=11pt,before skip=11pt,after skip=15pt,
  fontupper=\popsemi\fontsize{10.6}{14.5}\selectfont\color{popink},
  fontlower=\popsemi\fontsize{10.6}{14.5}\selectfont\color{popink}}}
% Coche dessinée (le glyphe ✓ n'existe pas dans les polices du document)
\newcommand{\popcheck}[1]{\tikz[baseline=-0.5ex]\draw[#1,line width=1.6pt,line cap=round,line join=round](-0.13,0.0)--(-0.04,-0.09)--(0.13,0.1);}
\providecommand{\coche}{\popcheck{popgreen}}
\providecommand{\resultat}[1]{\par\smallskip\noindent\fcolorbox{popgreen}{popgreenL}{\parbox{\dimexpr\linewidth-2\fboxsep-2\fboxrule\relax}{\textbf{Result.} #1}}\par\smallskip}
\newcommand{\popboxhead}[3]{\poppill{#1}{#2}{white}\ifx\relax#3\relax\else\hspace{6pt}{\popxbold\fontsize{10.6}{12}\selectfont\color{popink}#3}\fi\par\vspace{7pt}}

\newtcolorbox{aretenir}[1][]{popbase,colframe=popgreen,before upper={\popboxhead{Key points}{popgreen}{#1}}}
\newtcolorbox{exemplebox}[1][]{popbase,colframe=poporange,before upper={\popboxhead{Exemple}{poporange}{#1}}}
\newtcolorbox{definition}[1][]{popbase,colframe=popblue,before upper={\popboxhead{Definition}{popblue}{#1}}}
\newtcolorbox{propriete}[1][]{popbase,colframe=poppurple,before upper={\popboxhead{Property}{poppurple}{#1}}}
\newtcolorbox{methode}[1][]{popbase,colframe=popgold,before upper={\popboxhead{Method}{popgold}{#1}}}
\newtcolorbox{attention}[1][]{popbase,colframe=poppink,before upper={\popboxhead{Warning}{poppink}{#1}}}
\newtcolorbox{experience}[1][]{popbase,colframe=popblue,colback=popblueL,before upper={\popboxhead{Experiment}{popblue}{#1}}}
% « Le savais-tu ? » : carte violette pleine (contexte, culture, Cameroun)
\newtcolorbox{savaistu}[1][]{popbase,colback=poppurple,colframe=popink,
  fontupper=\popsemi\fontsize{10.4}{14.2}\selectfont\color{white},
  before upper={\poppill{Did you know?}{white}{poppurple}\ifx\relax#1\relax\else\hspace{6pt}{\popxbold\color{white}#1}\fi\par\vspace{7pt}}}
% Exercice résolu : énoncé en haut, \tcblower puis la solution sur fond crème
\newcounter{popexo}\newcounter{popexr}
\newtcolorbox{exoresolu}[1][]{popbase,colframe=poporange,colbacklower=popcream,
  segmentation style={popink,line width=1.2pt,dashed},
  before upper={\stepcounter{popexr}\popboxhead{Worked example \thepopexr}{poporange}{#1}},
  before lower={\poppill{Solution}{popink}{popcream}\par\vspace{6pt}}}
% Exercice (non résolu en place) — la correction est dans la partie Corrigés
\newtcolorbox{exercice}[1][]{popbase,colframe=popink,
  before upper={\stepcounter{popexo}\popboxhead{Exercise \thepopexo}{popink}{#1}}}
% Corrigé
\newtcolorbox{corrige}[1][]{popbase,colframe=popgreen,colback=popgreenL,
  before upper={\popboxhead{Solution}{popgreen}{#1}}}
% Objectifs / prérequis / bilan
\newtcolorbox{objectifs}{popbase,colframe=popink,colback=poporangeL,
  before upper={\popboxhead{Chapter objectives}{popink}{}}}
\newtcolorbox{prerequis}{popbase,colframe=popink,colback=popblueL,
  before upper={\popboxhead{Prerequisites}{popblue}{}}}
\newtcolorbox{bilan}{popbase,colframe=popink,colback=white,boxrule=2.8pt,
  before upper={\popboxhead{Fiche bilan}{poporange}{L'essentiel à connaître pour l'examen}}}
% Figure encadrée avec légende
\newtcolorbox{popfigure}[1][]{enhanced,breakable=false,colback=white,colframe=popline,boxrule=1.4pt,arc=8pt,
  left=10pt,right=10pt,top=10pt,bottom=8pt,before skip=10pt,after skip=12pt,halign=center,
  lower separated=false,halign lower=center,
  fontlower=\popsemi\fontsize{9}{11.5}\selectfont\color{popmuted},
  #1}
\newcommand{\legende}[1]{\par\vspace{6pt}{\popsemi\fontsize{9}{11.5}\selectfont\color{popmuted}#1}}

% ============================================================
% Signature OnBuch+
% ============================================================
\newcommand{\poplogoinline}[1]{{\poplogo\fontsize{#1}{#1}\selectfont\color{popink}OnBuch\color{poporange}+}}
\newcommand{\signatureonbuch}{%
  \begin{tcolorbox}[enhanced,colback=popink,colframe=popink,boxrule=2.2pt,arc=10pt,
      drop shadow=poporange,left=14pt,right=14pt,top=10pt,bottom=10pt,before skip=0pt,after skip=0pt]
    \tikz[baseline=-0.7ex]\node[circle,fill=popgreen,draw=popcream,line width=1.3pt,minimum size=22pt,inner sep=0]{\popcheck{white}};%
    \hspace{9pt}\begin{minipage}[c]{0.62\linewidth}
      {\popxbold\fontsize{12}{13}\selectfont\color{popcream}Signed\ \textbullet\ {\poplogo OnBuch\color{poporange}+}}\par\vspace{2pt}
      {\raggedright\popsemi\fontsize{8}{10.5}\selectfont\color{popline}\MakeUppercase{OnBuch+ teaching team}\\\MakeUppercase{Follows the official MINESEC syllabus}\par}
    \end{minipage}\hfill
    {\poplogo\fontsize{22}{22}\selectfont\color{popcream}OnBuch\color{poporange}+}
  \end{tcolorbox}%
}

% ============================================================
% Page de garde
% #1 titre · #2 matière · #3 niveau · #4 module · #5 n° leçon · #6 durée ·
% #7 description / objectifs (texte ou itemize)
% ============================================================
\newcommand{\pagegarde}[7]{%
  \thispagestyle{empty}
  \newgeometry{a4paper,margin=1.7cm,top=1.6cm,bottom=1.4cm}%
  \begin{tikzpicture}[remember picture,overlay]
    \fill[poporange!18] ([xshift=-3.2cm,yshift=-3.6cm]current page.north east) circle (4.6cm);
    \fill[poppurple!14] ([xshift=2.4cm,yshift=7.6cm]current page.south west) circle (3.8cm);
  \end{tikzpicture}%
  {\poplogo\fontsize{30}{30}\selectfont\color{popink}OnBuch\color{poporange}+}\hfill
  \raisebox{0.6ex}{\poppill{Signed course \textbullet\ Premium}{popink}{popcream}}\par
  \vspace{1pt}\popsquiggle{poppurple}\par
  \vspace{8pt}
  \begin{tcolorbox}[enhanced,colback=poporange,colframe=popink,boxrule=2.8pt,arc=13pt,
      drop shadow=popink,left=18pt,right=18pt,top=12pt,bottom=13pt,after skip=10pt]
    \poppill{#2 \textbullet\ #3}{popink}{popcream}\hspace{5pt}\poppill{#4}{white}{popink}\par\vspace{8pt}
    {\popdisplay\fontsize{14}{14}\selectfont\color{popink}CHAPTER #5}\par\vspace{4pt}
    {\raggedright\hyphenpenalty=10000\exhyphenpenalty=10000\popdisplay\fontsize{25}{27}\selectfont\color{popcream}\MakeUppercase{#1}\par}
    \vspace{8pt}
    {\popxbold\fontsize{9.5}{9.5}\selectfont\color{popink}\MakeUppercase{Suggested time: #6}}
  \end{tcolorbox}
  \begin{tcolorbox}[enhanced,colback=white,colframe=popink,boxrule=2.2pt,arc=10pt,
      drop shadow=popink,left=16pt,right=16pt,top=10pt,bottom=10pt,after skip=8pt]
    \poppill{In this chapter}{poppurple}{white}\par\vspace{5pt}
    \popsemi\fontsize{9.3}{12.6}\selectfont\color{popink}#7
  \end{tcolorbox}
  \noindent\begin{minipage}[t]{0.487\linewidth}
    \begin{tcolorbox}[enhanced,colback=popgreen,colframe=popink,boxrule=2.2pt,arc=10pt,
        drop shadow=popink,left=11pt,right=11pt,top=10pt,bottom=10pt,height=3.1cm,valign=center]
      \tcbox[colback=white,colframe=popink,boxrule=1.3pt,arc=5pt,boxsep=0pt,nobeforeafter,
        left=3pt,right=3pt,top=3pt,bottom=3pt,valign=center]{\qrcode[height=1.9cm]{https://play.google.com/store/apps/details?id=com.luvvix.onbuch&hl=en}}%
      \hspace{8pt}\begin{minipage}[c]{0.5\linewidth}
        {\popdisplay\fontsize{11}{12.5}\selectfont\color{white}\MakeUppercase{Download OnBuch}}\par\vspace{4pt}
        {\raggedright\popsemi\fontsize{8.4}{11}\selectfont\color{white}Free \textbullet\ Google Play\\AI tutor, past papers, results\par}
      \end{minipage}
    \end{tcolorbox}
  \end{minipage}\hfill
  \begin{minipage}[t]{0.487\linewidth}
    \begin{tcolorbox}[enhanced,colback=poppurple,colframe=popink,boxrule=2.2pt,arc=10pt,
        drop shadow=popink,left=11pt,right=11pt,top=10pt,bottom=10pt,height=3.1cm,valign=center]
      \tikz[baseline=-0.7ex]\node[circle,fill=white,draw=popink,line width=1.3pt,minimum size=1.6cm,inner sep=0]{\popdisplay\fontsize{24}{24}\selectfont\color{poppurple}+};%
      \hspace{8pt}\begin{minipage}[c]{0.6\linewidth}
        {\popdisplay\fontsize{11}{12.5}\selectfont\color{white}\MakeUppercase{Go OnBuch+}}\par\vspace{4pt}
        {\raggedright\popsemi\fontsize{8.4}{11}\selectfont\color{white}All signed courses, unlimited AI tutor, PDF sheets — OnBuch+ tab in the app\par}
      \end{minipage}
    \end{tcolorbox}
  \end{minipage}
  \vspace{8pt}
  \popcontenu
  \vfill
  \signatureonbuch
  \restoregeometry
  \newpage
}

% Rangée « Ce cours contient » (page de garde)
\newcommand{\poptile}[3]{% #1 couleur #2 gros texte #3 légende
  \begin{tcolorbox}[enhanced,colback=#1,colframe=popink,boxrule=2pt,arc=9pt,shadow={2.5pt}{-2.5pt}{0pt}{popink},
      left=6pt,right=6pt,top=9pt,bottom=9pt,halign=center,valign=center,height=1.75cm,width=\linewidth,nobeforeafter]
    {\popdisplay\fontsize{12.5}{13}\selectfont\color{white}\MakeUppercase{#2}}\par\vspace{3pt}
    {\centering\popsemi\fontsize{7.8}{9.5}\selectfont\color{white}#3\par}
  \end{tcolorbox}}
\newcommand{\popcontenu}{%
  \par\noindent{\popxboldls\fontsize{7.5}{7.5}\selectfont\color{popmuted}THIS SIGNED COURSE CONTAINS}\par\vspace{7pt}
  \noindent\begin{minipage}[t]{0.235\linewidth}\poptile{popblue}{Course}{complete, illustrated, step by step}\end{minipage}\hfill
  \begin{minipage}[t]{0.235\linewidth}\poptile{poporange}{Methods}{and worked examples}\end{minipage}\hfill
  \begin{minipage}[t]{0.235\linewidth}\poptile{poppink}{Exercises}{exam style + solutions}\end{minipage}\hfill
  \begin{minipage}[t]{0.235\linewidth}\poptile{popgreen}{Summary sheet}{the essentials to remember}\end{minipage}\par}

% Sommaire
\newtcolorbox{sommaire}{enhanced,colback=white,colframe=popink,boxrule=2.2pt,arc=10pt,
  drop shadow=popink,left=18pt,right=18pt,top=13pt,bottom=13pt,before skip=6pt,after skip=20pt,
  before upper={\poppill{Contents}{poppurple}{white}\par\vspace{9pt}\popsemi\fontsize{10.6}{14.5}\selectfont\color{popink}}}

% Signature de fin de cours — poussée tout en bas de la dernière page
\newcommand{\findecours}{%
  \par\vfill\nopagebreak
  \begin{center}\popsquiggle{poporange}\end{center}\vspace{4pt}
  \signatureonbuch
}
\AtBeginDocument{\def\llbracket{\mathopen{[\mkern-3.5mu[}}\def\rrbracket{\mathclose{]\mkern-3.5mu]}}} % intervalles d'entiers (glyphe ⟦ absent de la police)
% Glyphes absents de Fira Math : redéfinis à partir de glyphes disponibles
\AtBeginDocument{%
  \def\setminus{\mathbin{\backslash}}\let\smallsetminus\setminus
  \def\vdots{\mathord{\vbox{\baselineskip3.2pt\lineskiplimit0pt\kern4pt\hbox{.}\hbox{.}\hbox{.}}}}%
  \def\ddots{\mathinner{\mkern1mu\raise6.5pt\hbox{.}\mkern2mu\raise3.6pt\hbox{.}\mkern2mu\raise0.7pt\hbox{.}\mkern1mu}}%
  \def\triangle{\mathord{\scalebox{0.9}{$\symup{\Delta}$}}}%
  \def\nabla{\mathord{\raisebox{\depth}{\scalebox{1}[-1]{$\symup{\Delta}$}}}}%
  \def\checkmark{\popcheck{popdark}}%
  \def\star{\mathbin{\ast}}\def\bigstar{\mathord{\ast}}%
  \def\diamond{\mathbin{\scalebox{0.6}{$\Diamond$}}}%
  \def\bigcup{\mathop{\mathchoice{\raisebox{-0.3em}{\scalebox{1.7}{$\cup$}}}{\scalebox{1.25}{$\cup$}}{\cup}{\cup}}\displaylimits}%
  \def\bigcap{\mathop{\mathchoice{\raisebox{-0.3em}{\scalebox{1.7}{$\cap$}}}{\scalebox{1.25}{$\cap$}}{\cap}{\cap}}\displaylimits}%
  \def\lesseqqgtr{\mathrel{\lessgtr}}\def\bigoplus{\mathop{\oplus}\displaylimits}%
}
\providecommand{\ssi}{if and only if} % abréviation fréquente
\AtBeginDocument{\providecommand{\wideparen}{\overparen}} % arc de cercle

% Fonctions usuelles que les modèles utilisent sans que amsmath les fournisse
\providecommand{\arsinh}{\operatorname{arsinh}}\providecommand{\arcosh}{\operatorname{arcosh}}
\providecommand{\artanh}{\operatorname{artanh}}\providecommand{\arsech}{\operatorname{arsech}}
\providecommand{\arcsch}{\operatorname{arcsch}}\providecommand{\arcoth}{\operatorname{arcoth}}
\providecommand{\arcsinh}{\operatorname{arsinh}}\providecommand{\arccosh}{\operatorname{arcosh}}
\providecommand{\arctanh}{\operatorname{artanh}}\providecommand{\sech}{\operatorname{sech}}
\providecommand{\csch}{\operatorname{csch}}\providecommand{\arcsec}{\operatorname{arcsec}}
\providecommand{\arccsc}{\operatorname{arccsc}}\providecommand{\arccot}{\operatorname{arccot}}
\providecommand{\sgn}{\operatorname{sgn}}\providecommand{\lcm}{\operatorname{lcm}}
\providecommand{\Var}{\operatorname{Var}}\providecommand{\Cov}{\operatorname{Cov}}

%% FIGFIX-BEGIN v4 — correctifs globaux des figures (docs/onbuch-generation/tools/figfix)
\usepackage{adjustbox}\usepackage{etoolbox}
% toute figure TikZ/pgfplots plus large que la ligne est réduite à la largeur disponible
\BeforeBeginEnvironment{tikzpicture}{\begin{adjustbox}{max width=\linewidth}}
\AfterEndEnvironment{tikzpicture}{\end{adjustbox}}
% légendes pgfplots lisibles par-dessus les courbes
\pgfplotsset{every axis/.append style={legend style={fill=white,fill opacity=0.92,text opacity=1,draw=black!15,inner sep=3pt}}}
% étiquettes d'arêtes (syntaxe edge["..."]) avec fond blanc
\tikzset{every edge quotes/.append style={fill=white,inner sep=1.2pt}}
%% FIGFIX-END
\begin{document}

\pagegarde{Money and banking}{Economics}{Upper Sixth}{Upper Sixth — Economics}{3}{7 periods}{This chapter explores the nature and functions of money, the forces determining its demand and supply, and how its value is measured through price indices and explained by the quantity theory of money. It then examines the structure and operations of commercial banks, including credit creation and the interpretation of a bank's balance sheet, with applications to the Cameroonian banking system.
\vspace{4pt}
\begin{multicols}{2}\small
\begin{itemize}
\item By the end of this chapter I can define money, trace its historical evolution and explain its functions and characteristics.
\item By the end of this chapter I can distinguish the motives for demanding money and identify the components of money supply in Cameroon (M1, M2).
\item By the end of this chapter I can explain how the value of money changes and relate it to the general price level.
\item By the end of this chapter I can construct and interpret a retail price index (Laspeyres) and compute rates of inflation.
\item By the end of this chapter I can state the Fisher equation of exchange (MV = PT) and use it to explain and calculate changes in the price level.
\item By the end of this chapter I can describe the types, functions and role of commercial banks in the Cameroonian economy.
\end{itemize}\end{multicols}}

\begin{sommaire}
\begin{multicols}{2}
\begin{enumerate}
\item Money: Nature, Evolution and Functions
\item The Demand for and Supply of Money
\item The Value of Money, Price Indices and the Quantity Theory
\item Commercial Banks: Types, Functions and Credit Creation
\item The Balance Sheet of a Commercial Bank
\item Integration activity
\item Methods and worked examples
\item Exercises
\item Solutions to the exercises
\item Summary sheet
\end{enumerate}
\end{multicols}
\end{sommaire}

\begin{objectifs}
\begin{itemize}
\item By the end of this chapter I can define money, trace its historical evolution and explain its functions and characteristics.
\item By the end of this chapter I can distinguish the motives for demanding money and identify the components of money supply in Cameroon (M1, M2).
\item By the end of this chapter I can explain how the value of money changes and relate it to the general price level.
\item By the end of this chapter I can construct and interpret a retail price index (Laspeyres) and compute rates of inflation.
\item By the end of this chapter I can state the Fisher equation of exchange (MV = PT) and use it to explain and calculate changes in the price level.
\item By the end of this chapter I can describe the types, functions and role of commercial banks in the Cameroonian economy.
\item By the end of this chapter I can explain the process of credit creation and calculate the credit multiplier.
\item By the end of this chapter I can read, construct and interpret the balance sheet of a commercial bank.
\end{itemize}
\end{objectifs}

\begin{prerequis}
\begin{itemize}
\item Basic national income and circular flow concepts from Lower Sixth
\item Elementary demand and supply analysis and price determination
\item Percentage and index-number calculations (mathematics)
\item The role of the State and economic agents in the economy
\item Basic knowledge of financial institutions in Cameroon (BEAC, commercial banks)
\end{itemize}
\end{prerequis}

\newpage

\coursec{Money: Nature, Evolution and Functions}

In 1994, the CFA franc was devalued by 50\% overnight: a trader in Mokolo market who had saved 100,000 FCFA suddenly found her money buying only half the goods it used to. Yet today, the same 100,000 FCFA can move from Yaoundé to Bamenda in seconds through MTN Mobile Money, without a single banknote changing hands. What exactly is money, why does its value rise and fall, and who creates it? This chapter takes you inside the world of money, prices and banks --- from the barter stalls of the past to the balance sheets of Afriland and SGC --- and shows how the quantity of money in circulation shapes the prices you pay every day.

Before we can analyse banks or inflation, we must answer a deceptively simple question: \emph{what is money?} You will discover that money is not defined by what it looks like, but by what people \emph{accept} it for.

\courssub{What is Money? From Observation to Definition}

\textbf{A market observation.}\par
Walk through Mokolo market in Yaoundé. A woman sells plantains and receives banknotes. A taxi driver accepts a 1,000 FCFA note for a ride he could have refused to price in yams or in litres of petrol. A shopkeeper in Bamenda accepts a mobile money transfer on a cracked phone screen and immediately trusts it. None of these people \emph{wants} the money for itself --- the trader cannot eat a banknote. They accept it because they are confident that \emph{someone else} will accept it in turn. This observation leads directly to the definition.

\begin{definition}[Money]
\cle{Money} is anything that is \textbf{generally acceptable} as a \textbf{medium of exchange} for goods and services and in the settlement of debts within a community. An object is money not because of its physical form (metal, paper, digits on a phone) but because of the \emph{general acceptability} it commands.
\end{definition}

Notice the two key ideas: \emph{anything} (the form can change over history) and \emph{generally acceptable} (the social convention that gives money its power). A 10,000 FCFA note is just printed paper; it is money because every trader, transporter and tax collector in Cameroon will take it.

\begin{attention}[Money is not wealth, and not income]
Candidates frequently confuse three distinct concepts:
\begin{itemize}
\item \cle{Money}: the stock of notes, coins and bank deposits you hold \emph{at a point in time} (e.g.\ 25,000 FCFA in your wallet and mobile account).
\item \cle{Wealth}: the total value of \emph{all} assets owned --- land, a house in Buea, cattle, shares, \emph{plus} money. A wealthy chief may hold very little money; a poor student may hold only money.
\item \cle{Income}: a \emph{flow} received over a period (e.g.\ a salary of 150,000 FCFA \emph{per month}).
\end{itemize}
Money is only one component of wealth, and it is a \emph{stock}, whereas income is a \emph{flow}. In the examination, never write ``money is wealth''.
\end{attention}

\courssub{The Barter Economy and Its Limitations}

\textbf{Imagine a world without money.}\par
Long ago, exchange took place by \cle{barter}: the direct exchange of goods for goods. A fisherman in Limbe who wants cassava must find a cassava farmer who wants fish. At first sight this seems natural, but try to organise a whole economy this way and it collapses. Barter suffers from four fundamental limitations.

\begin{enumerate}
\item \textbf{The double coincidence of wants.} Exchange requires that each party wants \emph{exactly} what the other offers, at the same time. If the cassava farmer does not want fish but wants a hoe, the fisherman must first find a blacksmith who wants fish, trade fish for a hoe, then trade the hoe for cassava. Searching for partners wastes enormous time and effort --- economists call these \emph{transaction costs}.
\item \textbf{Indivisibility of goods.} Some goods cannot be divided without losing their value. A cattle owner in the North-West who wants only a bag of maize cannot cut off a fraction of a live cow. Barter makes small, everyday purchases of large-value assets impossible.
\item \textbf{No common measure of value.} Under barter, every good must be priced in terms of every other good: how many fish for a goat? How many goats for a bicycle? With $n$ goods there are $\frac{n(n-1)}{2}$ exchange ratios to remember (with 100 goods, that is 4,950 prices!). There is no single ``price'' to compare values.
\item \textbf{Difficulty of storing value and of deferred payment.} Fish rots, plantains ripen. Saving under barter means storing goods that perish or lose value. Lending is equally awkward: if you borrow three baskets of garri today, in what exactly must you repay next season, and how do you write that contract?
\end{enumerate}

\begin{popfigure}
\begin{center}
\begin{tikzpicture}[
box/.style={draw=popink, rounded corners=2pt, fill=popblueL, text width=2.6cm, align=center, minimum height=1.1cm, font=\small},
bad/.style={-{Stealth[length=2.5mm]}, thick, poporange},
good/.style={-{Stealth[length=2.5mm]}, thick, popgreen},
lbl/.style={font=\scriptsize\itshape, align=center}]
% Barter side
\node[box, align=center] (fish) at (0,2.2){Fisherman\\ has fish, wants cassava};
\node[box, align=center] (farm) at (5.4,2.2){Farmer\\ has cassava, wants a hoe};
\draw[bad] (fish) -- node[lbl, above]{wants to trade} (farm);
\draw[bad] (farm) -- node[lbl, below]{does NOT want fish} (fish);
\node[lbl, poporange, align=center] at (2.7,0.9){No double coincidence:\\ exchange FAILS};
% Money side
\node[box, align=center] (fish2) at (0,-2.2){Fisherman\\ sells fish};
\node[box, fill=popgreenL, align=center] (money) at (5.4,-2.2){Money\\ (FCFA)};
\node[box, align=center] (farm2) at (10.8,-2.2){Farmer\\ sells cassava};
\draw[good] (fish2) -- node[lbl, above]{fish} (money);
\draw[good] (money) -- node[lbl, above]{payment} (fish2);
\draw[good] (money) -- node[lbl, above]{payment} (farm2);
\draw[good] (farm2) -- node[lbl, above]{cassava} (money);
\node[lbl, popgreen, align=center] at (5.4,-3.7){Money breaks the exchange into two simple sales:\\ exchange SUCCEEDS};
\node[font=\small\bfseries] at (2.7,3.6) {BARTER};
\node[font=\small\bfseries] at (5.4,-0.9) {MONETARY EXCHANGE};
\end{tikzpicture}
\legende{Barter requires a double coincidence of wants; money removes it by splitting exchange into a sale and a purchase.}
\end{center}
\end{popfigure}

Money solves all four problems at once: it is wanted by everyone (no coincidence needed), it is divisible, it provides a single unit in which all prices are expressed, and it stores value across time.

\courssub{The Historical Evolution of Money}

Money was not invented once; it evolved through stages, each correcting the weaknesses of the previous one.

\begin{enumerate}
\item \textbf{Commodity money.} The earliest money was a useful, desired commodity: cowrie shells (used for centuries in West and Central Africa), salt, cattle, cocoa pods, cloth. These had \emph{intrinsic value} --- they were useful in themselves --- which made them acceptable. But cowries could be counterfeited, cattle could die, and salt was bulky.
\item \textbf{Metallic money.} Precious metals (gold, silver, bronze) were scarce, durable, divisible and homogeneous. Coins stamped by an authority guaranteed weight and purity. Their drawback: heavy to carry in large amounts and tempting to clip or debase.
\item \textbf{Paper money (banknotes).} Originally, notes were \emph{claims} on gold held by a bank or the state --- the note itself had almost no intrinsic value, but it was convertible. Modern notes are \emph{fiduciary money}: they circulate purely on \emph{trust} (\emph{fides} = confidence in Latin) in the issuing authority. Paper is light and cheap to produce, but wears out and can be forged.
\item \textbf{Bank money (deposit money).} The largest part of money today is not printed at all: it is the \emph{book-entry} balances held in current accounts at commercial banks (Afriland, SGC, BICEC), transferred by cheque, card or bank transfer. A deposit is money because it is generally acceptable in payment.
\item \textbf{Electronic and mobile money.} The newest stage: value stored electronically and transferred by phone or computer. In Cameroon, \emph{MTN Mobile Money} and \emph{Orange Money} allow a mother in Yaoundé to pay school fees in Bamenda in seconds, and a farmer in Kupe Muanenguba to receive payment for cocoa without travelling to a bank branch.
\end{enumerate}

\begin{popfigure}
\begin{center}
\begin{tikzpicture}[
stage/.style={draw=popink, rounded corners=2pt, fill=popblueL, text width=2.5cm, align=center, minimum height=1.3cm, font=\small},
arr/.style={-{Stealth[length=2.5mm]}, very thick, poppurple}]
\node[stage, align=center] (s1) at (0,0){Commodity money\\ \scriptsize cowries, salt, cattle};
\node[stage, align=center] (s2) at (3.4,0){Metallic money\\ \scriptsize gold \& silver coins};
\node[stage, align=center] (s3) at (6.8,0){Paper money\\ \scriptsize banknotes};
\node[stage, align=center] (s4) at (10.2,0){Bank money\\ \scriptsize deposits, cheques};
\node[stage, fill=popgreenL, align=center] (s5) at (13.6,0){E-money\\ \scriptsize MoMo, Orange};
\draw[arr] (s1) -- (s2);
\draw[arr] (s2) -- (s3);
\draw[arr] (s3) -- (s4);
\draw[arr] (s4) -- (s5);
\node[font=\scriptsize\itshape, text width=13.5cm, align=center] at (6.8,-1.4) {Each stage improves portability, divisibility and convenience --- and relies less on intrinsic value, more on trust.};
\end{tikzpicture}
\legende{The evolution of money: from commodities with intrinsic value to pure trust-based electronic money.}
\end{center}
\end{popfigure}

\begin{savaistu}[Mobile money in Cameroon]
Cameroon is one of the leaders of mobile money in the CEMAC zone. Millions of Cameroonians --- including market traders, moto-taxi drivers and students --- hold value in MTN MoMo or Orange Money wallets, many of them without any traditional bank account. Economists debate whether mobile money is a \emph{new} money or simply a new \emph{way of moving} bank money; what is certain is that it performs the functions of money for its users every day.
\end{savaistu}

\courssub{The Four Functions of Money}

Why is money so powerful? Because it performs four functions that no barter good could perform simultaneously.

\begin{propriete}[The functions of money]
Money serves four functions (learn them with the mnemonic \emph{``M-U-S-S''}):
\begin{enumerate}
\item \cle{Medium of exchange}: money is accepted in payment for all goods and services, eliminating the double coincidence of wants. \emph{Example:} you pay your Douala--Yaoundé bus fare in FCFA, not in plantains.
\item \cle{Unit of account} (measure of value): money is the common standard in which prices, debts and accounts are expressed. \emph{Example:} a phone is priced at 85,000 FCFA and a bag of rice at 28,000 FCFA, so we instantly know the phone ``costs'' about three bags of rice --- no need for 4,950 barter ratios.
\item \cle{Store of value}: money allows purchasing power to be kept and used later, linking the present and the future. \emph{Example:} a trader saves 100,000 FCFA in September to pay school fees in January. (This function weakens when inflation is high, as the 1994 devaluation showed.)
\item \cle{Standard of deferred payment}: money is the unit in which future payments --- loans, rents, hire-purchase instalments --- are fixed. \emph{Example:} a shopkeeper sells a freezer on credit: ``20,000 FCFA per month for six months.''
\end{enumerate}
This classification is examinable; be ready to \emph{illustrate} each function, not merely list it.
\end{propriete}

\begin{attention}[The functions are linked to confidence]
The store-of-value and deferred-payment functions depend on the \emph{stability} of money's value. When prices rise sharply (inflation) or the currency is devalued, people rush to hold goods, foreign currency or land instead of money --- and money begins to fail at its job. This is why the 1994 devaluation hurt savers so much.
\end{attention}

\courssub{The Qualities of Good Money}

Not every object can serve as money. Experience shows that good money must possess the following characteristics:

\begin{itemize}
\item \textbf{General acceptability}: everyone must be willing to take it in payment --- the supreme quality, from which all others follow.
\item \textbf{Durability}: it must not perish or wear out quickly (why cattle and fish failed; why coins succeeded).
\item \textbf{Divisibility}: it must split into small units for small purchases (the FCFA divides into 1, 5, 10, 25, 50, 100, 500 coins and larger notes).
\item \textbf{Portability}: it must be easy to carry --- high value in small weight and volume (a phone wallet beats a bag of cowries).
\item \textbf{Homogeneity} (uniformity): every unit must be identical to every other unit of the same denomination, so no one argues about \emph{which} 1,000 FCFA note is offered.
\item \textbf{Scarcity} (limited supply): it must not be so abundant that it loses value --- this is why the power to issue money is reserved to a central bank.
\item \textbf{Stability of value}: its purchasing power should be reasonably constant over time, so people dare to save and to write contracts in it.
\end{itemize}

\begin{methode}[Testing whether an asset qualifies as money]
To decide whether an asset (a cowrie, a deposit, a MoMo balance, a treasury bill) is money, apply the two-step test:
\begin{enumerate}
\item \textbf{Acceptability test}: is it \emph{generally acceptable} in final payment for goods, services and debts in the community?
\item \textbf{Functions test}: does it perform the four functions --- medium of exchange, unit of account, store of value, standard of deferred payment --- especially the first?
\end{enumerate}
If both answers are yes, it is money. If it must first be \emph{converted} into something else before it can be spent (like a savings deposit or a share), it is \emph{near-money}, not money proper.
\end{methode}

\courssub{Legal Tender, the CFA Franc and the CEMAC Zone}

\begin{definition}[Legal tender]
\cle{Legal tender} is money that, by law, \emph{must} be accepted in settlement of a debt within a country; a creditor cannot legally refuse it. In Cameroon, banknotes and coins issued by the BEAC are legal tender. A cheque is \emph{not} legal tender --- a seller may legally refuse it --- even though bank deposits are money.
\end{definition}

The currency of Cameroon is the \cle{CFA franc} (FCFA), whose full name in our region is the \emph{franc de la Coopération Financière en Afrique centrale}. It is issued not by Cameroon alone but by the \cle{BEAC} (\emph{Banque des États de l'Afrique Centrale}), the common central bank of the \cle{CEMAC} zone (Cameroon, Chad, Central African Republic, Republic of Congo, Equatorial Guinea and Gabon).

\begin{aretenir}[Key facts on the CFA franc]
\begin{itemize}
\item The CFA franc used in Cameroon is issued by the \textbf{BEAC} for the \textbf{CEMAC} monetary zone.
\item It is \textbf{pegged to the euro} at a fixed parity, guaranteed by the French Treasury under cooperation agreements.
\item It is \textbf{legal tender} throughout the CEMAC zone.
\item Its value was halved against the French franc in the \textbf{1994 devaluation} --- a reminder that even legal tender can lose purchasing power.
\end{itemize}
\end{aretenir}

\courssub{Money and Near-Money: The Question of Liquidity}

\textbf{Observation.}\par
A student holds 15,000 FCFA in her wallet, 40,000 FCFA in a current account at SGC, and 100,000 FCFA in a blocked savings account that requires one month's notice to withdraw. Which of these is money?

The answer introduces the idea of \cle{liquidity}: the ease and speed with which an asset can be converted into a means of payment \emph{without loss of value}.

\begin{definition}[Near-money]
\cle{Near-money} consists of assets that are highly liquid and act as good stores of value but are \emph{not directly usable} as a medium of exchange: they must first be converted into cash or a transferable deposit. Examples: savings and time deposits, treasury bills, shares and bonds.
\end{definition}

Notes, coins, current-account (demand) deposits and mobile-money balances pass the acceptability test --- they are \emph{money proper}. The savings deposit fails it: you cannot pay a taxi driver with a savings booklet; you must first withdraw. It is therefore near-money. The distinction matters because, as we shall see in the next lesson, economists define the \emph{money supply} by deciding exactly where to draw this line.

\begin{exoresolu}[Is it money? Applying the test]
Classify the following as \emph{money} or \emph{near-money} in Cameroon today, justifying each answer: (a) a 2,000 FCFA BEAC note; (b) a balance of 30,000 FCFA on an MTN MoMo wallet; (c) a 6-month treasury bill; (d) a cow in Adamawa; (e) a demand deposit at Afriland First Bank.
\tcblower
Apply the two-step test (acceptability + functions):
\begin{itemize}
\item \textbf{(a) The banknote}: generally acceptable, legal tender, performs all four functions $\Rightarrow$ \emph{money}.
\item \textbf{(b) MoMo balance}: accepted by millions of users and merchants for payment and transfers, stores value, denominated in FCFA $\Rightarrow$ functions as \emph{money} (electronic money) for its users.
\item \textbf{(c) Treasury bill}: an excellent store of value, but it must be sold or held to maturity before it can be spent; a trader will not accept it for bread $\Rightarrow$ \emph{near-money}.
\item \textbf{(d) The cow}: not divisible, not portable, not homogeneous, and not generally acceptable in Yaoundé shops $\Rightarrow$ \emph{neither} money nor near-money; it is simply wealth (a real asset).
\item \textbf{(e) Demand deposit}: transferable instantly by cheque, card or transfer and acceptable in payment $\Rightarrow$ \emph{money} (bank money).
\end{itemize}
\end{exoresolu}

\begin{exercice}[Check your understanding]
\begin{enumerate}
\item Define money and explain why general acceptability is its essential feature.
\item A yam farmer wants a pair of shoes. Using this example, explain the problem of double coincidence of wants and how money solves it.
\item Identify which function of money is illustrated in each case: (i) a landlord fixes rent at 60,000 FCFA per month, payable over one year; (ii) a trader keeps her daily takings in a cash box; (iii) prices in a supermarket are all labelled in FCFA.
\item Explain why a savings account at a microfinance institution in Bafoussam is near-money rather than money.
\end{enumerate}
\end{exercice}

\begin{corrige}[Exercise 1]
\begin{enumerate}
\item Money is anything generally acceptable as a medium of exchange and in settlement of debts. Acceptability is essential because money has little or no intrinsic use: people hold it only because they trust others will accept it in turn; without general acceptability no object can perform the functions of money.
\item The farmer must find a shoemaker who simultaneously wants yams --- a double coincidence. If the shoemaker wants fish instead, the farmer must make a chain of intermediate exchanges, wasting time. With money, the farmer simply sells yams for FCFA and buys shoes with FCFA: money splits the exchange into two independent sales.
\item (i) Standard of deferred payment; (ii) store of value; (iii) unit of account.
\item The savings balance cannot be used directly to pay for goods; it must first be withdrawn as cash or transferred to a current account. It is a liquid store of value but not a medium of exchange, hence near-money.
\end{enumerate}
\end{corrige}

\begin{aretenir}[Summary of Section 1]
\begin{itemize}
\item Money is \emph{anything generally acceptable} as a medium of exchange; it is a stock, distinct from wealth and income.
\item Barter fails because of the double coincidence of wants, indivisibility, the absence of a common measure of value, and the difficulty of storing value and deferring payment.
\item Money evolved: commodity $\rightarrow$ metallic $\rightarrow$ paper $\rightarrow$ bank (deposit) $\rightarrow$ electronic/mobile money.
\item Four functions: medium of exchange, unit of account, store of value, standard of deferred payment.
\item Good money is acceptable, durable, divisible, portable, homogeneous, scarce and stable in value.
\item The CFA franc is legal tender in Cameroon, issued by the BEAC for the CEMAC zone and pegged to the euro.
\item Near-money (savings deposits, treasury bills) is liquid but must be converted before it can be spent.
\end{itemize}
\end{aretenir}

\coursec{The Demand for and Supply of Money}

In the previous section we established what money \emph{is}: anything generally accepted as a medium of exchange, unit of account, store of value and standard of deferred payment. We now ask the natural follow-up questions. \emph{Why} do people choose to hold part of their wealth as money, even though money earns little or no interest? \emph{Who} decides how much money exists in the economy? And how do these two forces --- the \cle{demand for money} and the \cle{supply of money} --- interact to determine the \cle{rate of interest}? These questions take us to the heart of the \cle{money market}, and the answers were given their classic form by the British economist John Maynard Keynes.

\courssub{The Demand for Money}

\textbf{Meaning of the demand for money}\par
Observe a trader at Mokolo market in Yaound\'e at the end of a good week. She has earned 200\,000 FCFA. She must decide in what \emph{form} to hold this wealth: as cash in her drawer, as a deposit at her bank, as a sack of cocoa beans, as a plot of land, or as a government bond. Holding wealth as \textbf{money} (cash or demand deposits) means holding it in the most \cle{liquid} form: it can be spent immediately, but it earns no interest (or very little). Holding it as a bond earns interest, but the bond must be sold before it can be spent, and its market price can fall.

\begin{definition}[Demand for money]
The \cle{demand for money} is the desire of households and firms to hold part of their wealth in liquid form --- as notes, coins and demand (current account) deposits --- rather than in interest-earning assets such as bonds, savings accounts or physical property. It is a demand for \emph{money to hold} (a stock concept, measured at a point in time), not a demand for money to spend (a flow concept, measured over a period).
\end{definition}

\begin{attention}[A frequent confusion]
The demand for money is \textbf{not} the desire to be rich. Everyone wants more wealth. The demand for money is the desire to hold a given wealth in \emph{liquid} form. A millionaire who keeps all his wealth in land and shares has a \emph{low} demand for money; a modest salaried worker who keeps everything in a current account has a \emph{high} demand for money relative to his wealth.
\end{attention}

\courssub{Keynes's Motives for Holding Money}

Keynes, in his \emph{General Theory of Employment, Interest and Money} (1936), identified three reasons --- three \cle{motives} --- why rational economic agents sacrifice interest income in order to hold money.

\begin{definition}[The transactions motive]
The \cle{transactions motive} is the desire to hold money in order to finance everyday, planned expenditure: buying food, paying transport fares, paying school fees. Because income is received at discrete intervals (monthly salary, seasonal cocoa sales) while spending occurs continuously, people must hold money to bridge the gap. The amount held depends mainly on the \textbf{level of income}: the higher the income of the household or firm, the larger its transactions balances.
\end{definition}

\begin{definition}[The precautionary motive]
The \cle{precautionary motive} is the desire to hold money as a reserve against \emph{unforeseen} events: sudden illness, a broken motorcycle, an unexpected funeral contribution, or an unexpected business opportunity requiring immediate cash. It is money held ``just in case''. Like the transactions motive, it rises with income, but it also rises with the degree of uncertainty about the future.
\end{definition}

\begin{definition}[The speculative motive]
The \cle{speculative motive} is the desire to hold money instead of interest-earning assets (bonds) in order to avoid capital losses and to be ready to profit from future changes in asset prices. A speculator who expects bond prices to \emph{fall} (that is, expects the interest rate to \emph{rise}) prefers to hold money now and buy bonds later, when they are cheaper.
\end{definition}

\begin{propriete}[Interest rate and speculative demand for money]
There is an \textbf{inverse relationship} between the rate of interest and the speculative demand for money. When the interest rate is \emph{high}: (i) the opportunity cost of holding money (foregone interest) is high, and (ii) most people expect the rate to fall (bond prices to rise) in the future, so they hold bonds, not money: speculative demand for money is \emph{low}. When the interest rate is \emph{low}: the reward for holding bonds is small and most people expect the rate to rise (bond prices to fall), so they prefer to hold money: speculative demand for money is \emph{high}. (The economic reasoning, not a formal mathematical proof, is examinable.)
\end{propriete}

\begin{exemplebox}[Bond prices and the interest rate]
Consider a government bond that pays a fixed coupon of 5\,000 FCFA per year. If the prevailing market interest rate is 10\%, the bond's market price is $5\,000 / 0.10 = 50\,000$ FCFA. If the market rate falls to 5\%, the same bond's price rises to $5\,000 / 0.05 = 100\,000$ FCFA. Hence a \emph{fall} in the interest rate means a \emph{rise} in bond prices --- which is why speculators expecting lower future rates buy bonds today (to sell later at a capital gain), while those expecting higher future rates hold money.
\end{exemplebox}

\courssub{Determinants of the Demand for Money}

\begin{itemize}
\item \textbf{Level of income (Y).} Higher national income raises transactions and precautionary balances: money demand rises. This is the primary determinant of the \emph{transactions} and \emph{precautionary} demand.
\item \textbf{General price level (P).} If prices double, twice as much nominal money is needed to carry out the same real volume of transactions: nominal money demand rises proportionally. Real money demand ($M_d/P$) depends on real income.
\item \textbf{Rate of interest (r).} It is the opportunity cost of holding money; a higher rate reduces money demand (especially the speculative component).
\item \textbf{Frequency of payments.} A worker paid weekly needs smaller average cash balances than one paid monthly for the same annual income. The more frequent the income receipts, the lower the demand for money.
\item \textbf{Financial innovation and payment habits.} The spread of mobile money services (such as MTN Mobile Money and Orange Money, widely used in Cameroon), ATMs, debit cards and electronic transfers reduces the need to hold physical cash, shifting the demand for narrow money downward.
\end{itemize}

\courssub{The Liquidity Preference Curve and the Liquidity Trap}

Keynes called the demand for money \cle{liquidity preference}. Plotting the quantity of money demanded ($M_d$) on the horizontal axis against the rate of interest ($r$) on the vertical axis gives the \cle{liquidity preference curve}: it slopes downward because of the inverse link between the interest rate and speculative balances. The transactions and precautionary demands are largely interest-inelastic (vertical), so the slope comes entirely from the speculative motive.

\begin{popfigure}
\begin{center}
\begin{tikzpicture}
\begin{axis}[
width=11cm, height=7.5cm,
axis lines=left,
xlabel={Quantity of money ($M$)},
ylabel={Rate of interest ($r$)},
xmin=0, xmax=10, ymin=0, ymax=10,
xtick=\empty, ytick=\empty,
every axis x label/.style={at={(ticklabel* cs:1.02)}, anchor=west},
every axis y label/.style={at={(ticklabel* cs:1.02)}, anchor=south},
clip=false
]
\addplot[very thick, popblue, domain=0.6:6.5, samples=100] {8*exp(-0.55*x)+1.2};
\addplot[very thick, popblue, domain=6.5:9.5, samples=20] {1.2};
\draw[dashed, popmuted] (axis cs:6.5,0) -- (axis cs:6.5,1.2);
\draw[dashed, popmuted] (axis cs:0,1.2) -- (axis cs:9.5,1.2);
\node[popblue, anchor=west] at (axis cs:1.2,6.2) {LP ($M_d$)};
\node[poporange, anchor=south west, align=center] at (axis cs:6.9,1.35) {Liquidity trap};
\node[popmuted, anchor=east] at (axis cs:0,1.2) {$r_{\min}$};
\end{axis}
\end{tikzpicture}
\end{center}
\legende{The liquidity preference (money demand) curve: at very low interest rates it becomes perfectly elastic --- the liquidity trap.}
\end{popfigure}

\begin{definition}[The liquidity trap]
The \cle{liquidity trap} is the situation in which the interest rate has fallen so low ($r_{\min}$) that almost everyone expects it to rise in the future (bond prices to fall). No one is willing to buy bonds; any additional money supplied is simply held as idle balances. The demand for money becomes \emph{infinitely interest-elastic} (horizontal) at the floor rate $r_{\min}$.
\end{definition}

\begin{savaistu}[Why the trap matters for policy]
In a liquidity trap, printing more money cannot push the interest rate any lower, so conventional monetary policy (open market operations) loses its power to stimulate investment. Keynes used this idea to explain why economies can remain stuck in deep depression with idle resources --- a lesson forcefully revisited after the 2008 global financial crisis and during the COVID-19 pandemic, when many central banks hit the ``zero lower bound''.
\end{savaistu}

\courssub{The Supply of Money}

\textbf{Meaning and measurement}\par
The \cle{supply of money} is the total quantity of money in circulation in an economy at a given point in time. Because ``money'' exists in several forms of differing liquidity, statisticians measure it using \cle{monetary aggregates}.

\begin{definition}[Monetary aggregates M1 and M2]
\textbf{M1 (Narrow money)} = Notes and coins in circulation \emph{plus} demand deposits (current account balances usable on demand by cheque, card, transfer or mobile money). \textbf{M2 (Broad money)} = M1 \emph{plus} savings and time deposits (deposits that require notice before withdrawal or have a fixed maturity). M2 therefore includes money proper plus \emph{near-money} assets that are highly liquid but not perfectly so.
\end{definition}

\begin{exemplebox}[Computing M1 and M2]
Suppose in an economy: notes and coins = 400 billion FCFA; demand deposits = 600 billion FCFA; savings and time deposits = 500 billion FCFA. Then:
\begin{align*}
M1 &= 400 + 600 = 1\,000 \text{ billion FCFA}\\
M2 &= M1 + 500 = 1\,500 \text{ billion FCFA}
\end{align*}
\end{exemplebox}

\textbf{Determinants of the money supply}\par
\begin{itemize}
\item \textbf{Central bank policy.} The central bank --- for Cameroon and the CEMAC zone, the \cle{BEAC} (Banque des \'Etats de l'Afrique Centrale), which issues the CFA franc --- controls the \emph{monetary base} (currency in circulation + commercial banks' reserves). It expands or contracts the base through its key instruments: the refinancing rate (taux de refinancement), reserve requirements (coefficients de r\'eserves obligatoires), and open market operations (op\'erations d'open market).
\item \textbf{Commercial bank credit creation.} Every loan granted by a commercial bank creates a new deposit; the lending behaviour of banks, constrained by reserve requirements and their own liquidity preferences, is the largest source of money supply growth (analysed in detail in Section 4 of this chapter).
\item \textbf{Government financing.} When government borrows directly from the central bank (monetisation of the deficit, colloquially ``printing money''), the monetary base and money supply rise. Borrowing from the public by selling bonds merely transfers existing money and does not directly increase $M_s$.
\item \textbf{Balance of payments.} A surplus on the overall balance of payments brings foreign currency into the country; when exporters convert it into CFA francs at the BEAC, the domestic monetary base and money supply increase. A deficit has the opposite effect.
\end{itemize}

\begin{attention}[The vertical money supply curve]
In the \emph{simple Keynesian model} used at this level, the money supply is treated as \textbf{exogenously fixed by the monetary authorities}, independent of the rate of interest. It is therefore drawn as a \emph{vertical} line. Do \textbf{not} draw it upward-sloping: a common exam mistake is to assume that a higher interest rate automatically increases the money supply because banks lend more. (In reality, the money supply is endogenous to some degree, but that refinement belongs to the credit-creation analysis in the next section, not to the basic money-market diagram.)
\end{attention}

\courssub{Money Market Equilibrium}

The \cle{money market} is in equilibrium when the quantity of money demanded equals the quantity supplied ($M_d = M_s$). The \emph{price} that clears this market is the \textbf{rate of interest}. If the interest rate is above equilibrium, the supply of money exceeds the demand: the public uses its surplus money balances to buy bonds, bond prices rise, and the interest rate falls back. If the rate is below equilibrium, money demand exceeds supply: the public sells bonds to obtain liquidity, bond prices fall, and the interest rate rises.

\begin{popfigure}
\begin{center}
\begin{tikzpicture}
\begin{axis}[
width=11cm, height=7.5cm,
axis lines=left,
xlabel={Quantity of money ($M$)},
ylabel={Rate of interest ($r$)},
xmin=0, xmax=10, ymin=0, ymax=10,
xtick=\empty, ytick=\empty,
every axis x label/.style={at={(ticklabel* cs:1.02)}, anchor=west},
every axis y label/.style={at={(ticklabel* cs:1.02)}, anchor=south},
clip=false
]
\addplot[very thick, popblue, domain=0.7:9.3, samples=100] {8*exp(-0.5*x)+0.8};
\node[popblue, anchor=west] at (axis cs:1.3,6.0) {LP ($M_d$)};
\addplot[very thick, poporange] coordinates {(4.5,0) (4.5,9.5)};
\addplot[very thick, popgreen] coordinates {(6.8,0) (6.8,9.5)};
\node[poporange, anchor=south] at (axis cs:4.5,9.5) {$M_{s1}$};
\node[popgreen, anchor=south] at (axis cs:6.8,9.5) {$M_{s2}$};
\draw[dashed, popmuted] (axis cs:0,3.6) -- (axis cs:4.5,3.6);
\draw[dashed, popmuted] (axis cs:0,2.1) -- (axis cs:6.8,2.1);
\node[popmuted, anchor=east] at (axis cs:0,3.6) {$r_1$};
\node[popmuted, anchor=east] at (axis cs:0,2.1) {$r_2$};
\addplot[only marks, mark=*, popdark] coordinates {(4.5,3.6) (6.8,2.1)};
\draw[-stealth, very thick, poppurple] (axis cs:4.9,8.6) -- (axis cs:6.4,8.6);
\node[poppurple, anchor=south] at (axis cs:5.6,8.8) {$\Delta M_s > 0$};
\end{axis}
\end{tikzpicture}
\end{center}
\legende{Money market equilibrium: an increase in money supply from $M_{s1}$ to $M_{s2}$ lowers the equilibrium interest rate from $r_1$ to $r_2$.}
\end{popfigure}

\begin{propriete}[Effect of a change in money supply]
Ceteris paribus, an \textbf{increase in the money supply} shifts the vertical supply curve to the \emph{right} and \textbf{lowers the equilibrium rate of interest}; a decrease in money supply shifts it left and raises the rate. The fall in the interest rate is the primary channel through which expansionary monetary policy stimulates investment and aggregate demand. \emph{Exception}: if the economy is in the liquidity trap (horizontal part of LP), extra money is hoarded and the interest rate does not fall.
\end{propriete}

\begin{methode}[Analysing a shift in money supply]
\begin{enumerate}
\item Draw the downward-sloping liquidity preference curve $M_d$ and the initial vertical supply curve $M_{s1}$; mark the initial equilibrium point $(M_1, r_1)$.
\item Identify the direction of the policy change: expansionary policy (BEAC lowers refinancing rate, buys securities, lowers reserve requirements, or banks lend more) shifts $M_s$ \emph{right}; contractionary policy shifts it \emph{left}.
\item Draw the new vertical curve $M_{s2}$ and read off the new intersection with $M_d$.
\item Conclude: rightward shift $\Rightarrow r$ falls (investment and aggregate demand stimulated); leftward shift $\Rightarrow r$ rises.
\item Check whether the economy is in the liquidity trap (horizontal part of LP): if so, state explicitly that the interest rate remains unchanged at $r_{\min}$.
\end{enumerate}
\end{methode}

\begin{exoresolu}[Worked exercise: money market equilibrium]
In the economy of ``Kamerland'', the demand for money is given by $M_d = 200 - 10r$ (billions of FCFA, $r$ in \%) and the money supply fixed by the central bank is $M_s = 120$ billion FCFA. (a) Find the equilibrium interest rate. (b) The central bank increases the money supply to 150 billion FCFA. Find the new equilibrium rate and comment on the transmission mechanism.
\tcblower
\textbf{(a)} Equilibrium requires $M_d = M_s$:
\begin{align*}
200 - 10r &= 120\\
10r &= 80\\
r_1 &= 8\%
\end{align*}

\textbf{(b)} With the new money supply $M_s = 150$:
\begin{align*}
200 - 10r &= 150\\
10r &= 50\\
r_2 &= 5\%
\end{align*}

\textbf{Comment.} The increase in money supply of 30 billion FCFA lowers the equilibrium interest rate from 8\% to 5\%. At the old rate of 8\%, the public now holds more money (150) than it wishes to hold (120 demanded); the excess money balances are used to purchase bonds, pushing bond prices up and the interest rate down until a new equilibrium is reached where the public willingly holds the larger money stock. This illustrates how expansionary monetary policy by the central bank cheapens the cost of borrowing and encourages investment spending.
\end{exoresolu}

\begin{exercice}[Practice]
Using the same demand function $M_d = 200 - 10r$ (billions of FCFA, $r$ in \%):
\begin{enumerate}
\item Find the equilibrium interest rate if the central bank \emph{reduces} the money supply to 100 billion FCFA.
\item Explain, in three or four sentences, the mechanism that moves the market from the old equilibrium (120 billion) to this new equilibrium.
\item State one circumstance in which a change in money supply would have \emph{no} effect on the interest rate.
\end{enumerate}
\end{exercice}

\begin{corrige}[Exercise 1]
\begin{enumerate}
\item $200 - 10r = 100 \Rightarrow 10r = 100 \Rightarrow r = 10\%$.
\item At the initial equilibrium ($r=8\%$, $M=120$), the supply falls to 100. At $r=8\%$, money demanded (120) now exceeds money supplied (100). The public sells bonds to rebuild cash balances; bond prices fall, and since bond prices and interest rates move inversely, the interest rate rises. The rise in $r$ reduces the speculative demand for money until it falls to 100 at the new equilibrium $r=10\%$.
\item In the \cle{liquidity trap}, where the liquidity preference curve is horizontal at $r_{\min}$: any change in money supply is absorbed by an equal change in idle balances, leaving the interest rate unchanged.
\end{enumerate}
\end{corrige}

\begin{aretenir}[Key points]
\begin{itemize}
\item The demand for money is the desire to hold wealth in \emph{liquid} form; Keynes identified the \textbf{transactions}, \textbf{precautionary} and \textbf{speculative} motives.
\item Transactions and precautionary demands depend mainly on income; speculative demand varies \emph{inversely} with the interest rate. At very low rates the economy may fall into a \textbf{liquidity trap} (horizontal LP curve).
\item Money supply is measured by the aggregates \textbf{M1} (notes, coins, demand deposits) and \textbf{M2} (M1 + savings and time deposits). Its determinants are central bank policy (BEAC), commercial bank lending, government financing mode, and the balance of payments.
\item In the simple model, money supply is a \emph{vertical} curve; equilibrium in the money market determines the \textbf{rate of interest}.
\item An increase in money supply shifts $M_s$ rightward and \emph{lowers} the equilibrium interest rate --- the key transmission channel of monetary policy (ineffective in a liquidity trap).
\end{itemize}
\end{aretenir}

\coursec{The Value of Money, Price Indices and the Quantity Theory}

In the previous section we studied how much money people wish to hold (the demand for money) and how the banking system creates and supplies it. But a natural question follows: what is all this money actually \emph{worth}? A 10\,000 FCFA note has no intrinsic value --- its value lies in what it can buy. When prices in the markets of Douala or Bamenda rise, the same note buys less; when prices fall, it buys more. This section therefore examines the \cle{value of money}, how economists measure its changes through \cle{price indices}, and the famous \cle{Quantity Theory of Money}, which links the supply of money to the general price level.

\courssub{The Meaning of the Value of Money}

\begin{definition}[Value of money]
The \cle{value of money} is the \textbf{purchasing power of money} --- the quantity of goods and services that a unit of money (for example one CFA franc) can buy at a given time. It is the reciprocal of the general price level.
\end{definition}

Money is desired not for its own sake but for what it commands in exchange. If a loaf of bread costs 250 FCFA, then 1\,000 FCFA ``is worth'' four loaves. If the price of bread rises to 500 FCFA, the same 1\,000 FCFA is worth only two loaves: the value of money has \emph{fallen}, even though the number of francs is unchanged.

\begin{propriete}[Inverse relationship between the value of money and the price level]
The value of money varies \textbf{inversely} with the general price level. If $P$ denotes the price index (base year = 100), then:
\begin{itemize}
\item The \textbf{purchasing power of one franc} (as a fraction of its base-year value) is $\dfrac{100}{P}$.
\item The \textbf{value of money index} (base year = 100) is $\dfrac{100}{P}\times 100$.
\end{itemize}
When prices double ($P$ goes from 100 to 200), the purchasing power is halved (from 1 to 0.5) and the value-of-money index falls from 100 to 50. When prices fall by half, the value of money doubles.
\end{propriete}

\begin{exemplebox}[The shrinking franc]
Suppose a representative basket of goods cost 20\,000 FCFA in 2015 (base year) and 30\,000 FCFA in 2024. The price index has risen from 100 to 150. The purchasing power of one franc has therefore fallen from $\frac{100}{100}=1$ to $\frac{100}{150}\approx 0.67$: each franc in 2024 buys only about two-thirds of what it bought in 2015. The value-of-money index falls from 100 to $\frac{100}{150}\times 100 \approx 66.7$.
\end{exemplebox}

\begin{popfigure}
\begin{center}
\begin{tikzpicture}
\begin{axis}[
  width=11cm, height=7cm,
  xlabel={General price level $P$ (index, base 100)},
  ylabel={Purchasing power $= 100/P$},
  xmin=40, xmax=210, ymin=0.3, ymax=2.6,
  grid=both, grid style={popline!40},
  axis line style={popink},
  label style={popink},
]
\addplot[domain=45:205, samples=100, very thick, popblue] {100/x};
\addplot[mark=*, only marks, poporange] coordinates {(50,2) (100,1) (200,0.5)};
\node[popdark, anchor=west, font=\small] at (axis cs:105,2.05) {prices halve $\Rightarrow$ value doubles};
\node[popdark, anchor=west, font=\small] at (axis cs:105,0.55) {prices double $\Rightarrow$ value halves};
\end{axis}
\end{tikzpicture}
\end{center}
\legende{The inverse (hyperbolic) relationship between the general price level and the purchasing power of money.}
\end{popfigure}

\courssub{Inflation, Deflation and Disinflation}

Because the value of money moves inversely to prices, any sustained movement in the general price level changes the value of money.

\begin{definition}[Inflation, deflation, disinflation]
\begin{itemize}
\item \cle{Inflation} is a \textbf{sustained rise in the general price level} over time. It causes a \textbf{fall in the value of money}.
\item \cle{Deflation} is a \textbf{sustained fall in the general price level}. It causes a \textbf{rise in the value of money}, but discourages spending and investment because consumers delay purchases expecting lower prices, and it raises the real burden of debt.
\item \cle{Disinflation} is a \textbf{fall in the rate of inflation}: prices are still rising, but more slowly (for example, inflation falling from 6\% to 3\%).
\end{itemize}
\end{definition}

\begin{attention}[One expensive good is not inflation]
A very common examination mistake is to call any price rise ``inflation''. A rise in the price of \emph{one} good (say, petrol or plantain) is a \textbf{relative price change}, not inflation. Inflation refers to a rise in the \textbf{general} price level --- the weighted average of prices across the whole economy. Similarly, a one-off jump in prices (e.g., due to a tax increase) that does not continue is not, strictly speaking, inflation, which must be \emph{sustained}.
\end{attention}

Inflation redistributes income: debtors gain (they repay in ``cheaper'' money) while creditors, pensioners on fixed incomes and cash holders lose. Deflation raises the real burden of debt and can deepen recessions. This is why the Bank of Central African States (BEAC) aims at low, stable inflation (typically around 3\% per year) within the CEMAC zone.

\courssub{Measuring Changes in the Value of Money: The Retail Price Index}

To know how the value of money changes, statisticians track the cost of a representative \cle{basket of goods and services} bought by a typical household.

\begin{definition}[Retail Price Index (RPI) / Consumer Price Index (CPI)]
The \cle{Retail Price Index} (RPI), also called the \cle{Consumer Price Index} (CPI), is a weighted average index number that measures the change over time in the cost of a fixed basket of goods and services purchased by a typical household. In Cameroon, the National Institute of Statistics (INS) computes the CPI using Yaoundé, Douala and other urban centres as reference markets. The base period is periodically updated (e.g., 2018 = 100).
\end{definition}

\begin{methode}[Constructing a Laspeyres price index (base-weighted)]
\begin{enumerate}
\item \textbf{Choose a base year} (year 0) whose index is set at 100.
\item \textbf{Select the basket}: a representative list of goods and services (food, transport, housing, clothing, education, health...).
\item \textbf{Fix the quantities} $q_0$ consumed of each item in the base year --- these quantities become the \textbf{weights}.
\item \textbf{Collect prices} in the base year ($p_0$) and in the current year ($p_1$).
\item \textbf{Compute the cost of the base-year basket} at current prices, $\sum p_1 q_0$, and at base-year prices, $\sum p_0 q_0$.
\item \textbf{Divide and multiply by 100} to obtain the index.
\end{enumerate}
\end{methode}

\begin{propriete}[Laspeyres index, inflation rate and value of money --- formulas to know]
The \cle{Laspeyres price index} uses base-year quantities as weights:
\[
\text{Laspeyres index } (P_1) = \frac{\sum p_1 q_0}{\sum p_0 q_0}\times 100
\]
The \cle{rate of inflation} between two periods (with indices $P_0$ and $P_1$) is:
\[
\text{Inflation rate} = \frac{P_1 - P_0}{P_0}\times 100
\]
The \cle{value of money} (purchasing power as a fraction of base year) and its \textbf{index} (base 100) are:
\[
\text{Purchasing power} = \frac{100}{P_1} \qquad\text{and}\qquad \text{Value-of-money index} = \frac{100}{P_1}\times 100
\]
These calculations are examinable in the GCE practical papers.
\end{propriete}

\begin{exoresolu}[Constructing a price index and measuring inflation]
A simplified household basket contains three items. Base-year (2020) quantities and prices, and 2024 prices, are:

\begin{center}
\begin{tabular}{|l|c|c|c|}
\hline
\rowcolor{popblueL} Item & $q_0$ (2020 qty) & $p_0$ (2020 price, FCFA) & $p_1$ (2024 price, FCFA) \\
\hline
Rice (kg) & 50 & 400 & 550 \\
\hline
Bread (loaves) & 100 & 250 & 350 \\
\hline
Transport (trips) & 40 & 500 & 600 \\
\hline
\end{tabular}
\end{center}

Calculate (a) the Laspeyres index for 2024 (base 2020), (b) the inflation rate between 2020 and 2024, (c) the value of money index in 2024.
\tcblower
\textbf{Step 1: Cost of the basket at base-year prices.}
\[
\sum p_0 q_0 = (400\times 50) + (250\times 100) + (500\times 40) = 20\,000 + 25\,000 + 20\,000 = 65\,000
\]
\textbf{Step 2: Cost of the same basket at 2024 prices.}
\[
\sum p_1 q_0 = (550\times 50) + (350\times 100) + (600\times 40) = 27\,500 + 35\,000 + 24\,000 = 86\,500
\]
\textbf{Step 3: The Laspeyres price index for 2024 (base 2020 = 100).}
\[
P_{2024} = \frac{86\,500}{65\,000}\times 100 \approx 133.1
\]
\textbf{Step 4: Inflation rate between 2020 and 2024.}
\[
\frac{133.1 - 100}{100}\times 100 = 33.1\%
\]
\textbf{Step 5: Value of money in 2024.}
\begin{itemize}
\item Purchasing power (fraction of 2020 value) $= \dfrac{100}{133.1} \approx 0.751$.
\item Value-of-money index (base 2020 = 100) $= \dfrac{100}{133.1}\times 100 \approx 75.1$.
\end{itemize}
Interpretation: prices rose by about 33\% between 2020 and 2024, so one franc in 2024 buys only about 75\% of what it bought in 2020.
\end{exoresolu}

\begin{popfigure}
\begin{center}
\begin{tikzpicture}
\begin{axis}[
  width=11cm, height=7cm,
  ybar, bar width=14pt,
  xlabel={Year}, ylabel={Consumer price index (base 2018 = 100)},
  symbolic x coords={2018,2019,2020,2021,2022,2023},
  xtick=data,
  ymin=0, ymax=140,
  grid=major, grid style={popline!40},
  axis line style={popink}, label style={popink},
  nodes near coords, every node near coord/.append style={font=\scriptsize, popdark},
]
\addplot[fill=popblue, draw=popdark] coordinates {(2018,100) (2019,103) (2020,106) (2021,110) (2022,118) (2023,127)};
\end{axis}
\end{tikzpicture}
\end{center}
\legende{A hypothetical consumer price index for Cameroon (base 2018 = 100). The steady rise of the index illustrates persistent inflation, i.e.\ a falling value of money.}
\end{popfigure}

\begin{attention}[Interpreting index numbers]
An index of 127 does \emph{not} mean ``prices are 127 FCFA''. It means prices are 27\% higher than in the base year. Also, when computing inflation \emph{between two non-base years}, divide by the \emph{first} year's index, not by 100: e.g.\ from 110 to 127, inflation $= \frac{127-110}{110}\times 100 \approx 15.5\%$, not 17\%.
\end{attention}

\courssub{The Quantity Theory of Money}

Why does the general price level rise? The classical answer, associated with the American economist Irving Fisher, is that prices depend on the quantity of money in circulation.

\begin{propriete}[Fisher's equation of exchange]
\[
MV = PT
\]
where:
\begin{itemize}
\item $M$ = the \textbf{money supply} (stock of money in circulation);
\item $V$ = the \textbf{velocity of circulation} --- the average number of times a unit of money changes hands in a period;
\item $P$ = the \textbf{general price level} of transactions;
\item $T$ = the \textbf{volume of transactions} (real output) in the economy.
\end{itemize}
The equation is an \textbf{identity}: it is true by definition, because the money spent ($MV$) must equal the money value of what is bought ($PT$).
\end{propriete}

\begin{exemplebox}[A numerical illustration of $MV = PT$]
Suppose an economy has a money supply $M = 500$ billion FCFA, each franc is spent on average $V = 4$ times a year, and the economy carries out $T = 2\,000$ billion units of transactions. Then:
\[
P = \frac{MV}{T} = \frac{500 \times 4}{2\,000} = 1.0
\]
If the central bank doubles the money supply to $M = 1\,000$ billion while $V$ and $T$ stay constant, then $P = \frac{1\,000 \times 4}{2\,000} = 2.0$: the price level doubles.
\end{exemplebox}

The equation becomes a \emph{theory} of inflation only when combined with two assumptions:

\begin{propriete}[Assumptions and conclusion of the Quantity Theory]
\textbf{Assumptions:}
\begin{itemize}
\item $V$ is \textbf{constant} in the short run, being determined by institutional habits (how often wages are paid, banking practices, payment technologies);
\item $T$ is \textbf{constant} in the short run, because the economy is assumed to be at or near full employment, so output cannot expand quickly.
\end{itemize}
\textbf{Conclusion:} with $V$ and $T$ fixed, $P = \dfrac{MV}{T}$ implies that any increase in the money supply $M$ causes a \textbf{proportionate increase in the price level}. Doubling $M$ doubles prices and halves the value of money. ``Inflation is always and everywhere a monetary phenomenon'' (monetarist view).
\end{propriete}

\begin{attention}[$V$ and $T$ are assumed constant]
A frequent error is to apply the theory mechanically. If $V$ falls (people hold money longer, as in a crisis) or $T$ rises (real output grows), an increase in $M$ need \emph{not} raise prices proportionately. Always state the assumptions before concluding --- examiners reward this.
\end{attention}

\textbf{Criticisms and limitations of the Quantity Theory.}
\begin{itemize}
\item \textbf{$V$ is not stable}: velocity varies with confidence, interest rates and financial innovation (mobile money in Cameroon, for instance, can speed up circulation).
\item \textbf{$T$ is not fixed}: in an economy with unemployed resources, extra money can raise \emph{output} rather than prices (the Keynesian criticism).
\item \textbf{Causality may run backwards}: rising costs (imported fuel, food) can push prices up and force the money supply to accommodate, rather than money causing prices.
\item \textbf{It ignores interest rates and expectations}, which influence how money affects spending.
\item It is a \textbf{long-run} relationship; in the short run the link between money and prices is loose.
\end{itemize}

\begin{savaistu}[The Cambridge cash-balance approach]
Cambridge economists (Marshall, Pigou) refined Fisher by asking how much money people \emph{choose to hold}:
\[
M = kPY
\]
where $Y$ is real national income, $P$ the price level, and $k$ the fraction of money income people wish to hold as cash balances. Since $k$ is the inverse of velocity ($k = 1/V$), this is mathematically close to Fisher, but it shifts attention to the \textbf{demand for money} --- linking this section neatly back to the previous one. If $k$ and $Y$ are stable, the same conclusion follows: more money means higher prices.
\end{savaistu}

\begin{aretenir}[Key points of the section]
\begin{itemize}
\item The value of money is its purchasing power; it moves \textbf{inversely} to the general price level: purchasing power $= 100/P$, value-of-money index $= (100/P)\times 100$.
\item Inflation (sustained general price rise) erodes the value of money; deflation raises it; disinflation is a slowing of inflation.
\item The RPI/CPI measures the cost of a fixed basket; the Laspeyres index is $\frac{\sum p_1 q_0}{\sum p_0 q_0}\times 100$; inflation $= \frac{P_1 - P_0}{P_0}\times 100$.
\item Fisher's equation $MV = PT$ is an identity; with $V$ and $T$ constant, money supply growth causes proportional inflation.
\item The theory is limited by unstable velocity, unemployed resources, and reverse causality; the Cambridge equation $M = kPY$ restates it in terms of money demand.
\end{itemize}
\end{aretenir}

\begin{exercice}[Practice: index numbers and the quantity theory]
\begin{enumerate}
\item The RPI of a country rises from 125 in 2022 to 140 in 2023. Calculate (a) the inflation rate for 2023; (b) the value of money index in each year (base 100).
\item In an economy, $M = 800$ billion FCFA, $V = 5$, $T = 2\,500$ billion units. (a) Find $P$. (b) If $M$ rises by 25\% while $V$ and $T$ are unchanged, find the new price level and comment.
\item Explain briefly why a 40\% rise in the price of petrol alone does not necessarily mean the economy is experiencing 40\% inflation.
\end{enumerate}
\end{exercice}

\begin{corrige}[Exercise 1]
\begin{enumerate}
\item (a) Inflation $= \frac{140-125}{125}\times 100 = 12\%$. (b) Value of money index: 2022: $\frac{100}{125}\times 100 = 80$; 2023: $\frac{100}{140}\times 100 \approx 71.4$. The franc's purchasing power fell from 80 to about 71.4 (base 100).
\item (a) $P = \frac{MV}{T} = \frac{800\times 5}{2\,500} = 1.6$. (b) New $M = 1\,000$ billion; $P = \frac{1\,000\times 5}{2\,500} = 2.0$. Prices rise by exactly 25\% (from 1.6 to 2.0), confirming the proportionality conclusion of the quantity theory \emph{given constant $V$ and $T$}.
\item Inflation is a rise in the \emph{general} price level measured by a weighted index of the whole basket. Petrol is only one item with a limited weight; its price rise raises the RPI only in proportion to its weight, and other prices may be stable or falling.
\end{enumerate}
\end{corrige}

Having understood what determines the value of money, we turn in the next section to the institutions that create and distribute most of the money in a modern economy: the \textbf{commercial banks}.

\coursec{Commercial Banks: Types, Functions and Credit Creation}

Having examined how the value of money is measured through price indices and the quantity theory, we now turn to the institutions that actually create and manage the bulk of the money supply in a modern economy: commercial banks. In Cameroon, as in all CEMAC countries, these banks operate under the supervision of the \cle{BEAC} (Banque des États de l'Afrique Centrale) and play a pivotal role in mobilising savings, financing investment and facilitating everyday transactions.

\courssub{Definition and Types of Commercial Banks}

\begin{definition}[Commercial Bank]
A \cle{commercial bank} is a financial institution licensed by the monetary authority to accept deposits from the public, make loans, and provide payment services. It creates \cle{bank deposits} (the bulk of the broad money supply M2) when it extends credit. In Cameroon, leading examples include \textbf{Afriland First Bank}, \textbf{Société Générale Cameroun (SGC)}, \textbf{BICEC} (Banque Internationale du Cameroun pour l'Épargne et le Crédit) and \textbf{Ecobank Cameroun}.
\end{definition}

\begin{definition}[Bank Deposit]
A \cle{bank deposit} is a liability of the bank owed to the depositor. It exists in three main forms:
\begin{itemize}
\item \cle{Current (demand) deposits} – withdrawable on demand without notice, typically pay no or very low interest, used for payments (cheques, transfers).
\item \cle{Savings deposits} – interest-bearing, subject to withdrawal limits or notice periods (e.g., \emph{livret d'épargne}).
\item \cle{Fixed (time) deposits} – fixed maturity (30 days to several years), higher interest, penalties for early withdrawal.
\end{itemize}
\end{definition}

\begin{aretenir}[Types of Banking Institutions in Cameroon (CEMAC Framework)]
\begin{itemize}
\item \textbf{Commercial banks (Universal banks)} – offer the full range of deposit, credit and payment services to households and firms. Licensed by COBAC.
\item \textbf{Merchant (investment) banks} – specialise in corporate finance, underwriting, long-term lending and advisory; they do not usually accept demand deposits from the general public.
\item \textbf{Central bank (BEAC)} – the monetary authority for the CEMAC zone; issues currency (CFA franc), sets the \cle{policy rate} (taux directeur), supervises banks via COBAC, acts as lender of last resort, and manages foreign reserves.
\item \textbf{Microfinance institutions (MFIs)} – serve low-income clients and small enterprises; regulated by COBAC under a specific framework (lighter capital requirements, limited product range). Examples: CAMCCUL, ACEP Cameroun, UCEC.
\item \textbf{Specialised financial institutions} – e.g., \textbf{Crédit Foncier du Cameroun} (housing finance), \textbf{Société Nationale d'Investissement (SNI)} (development finance).
\end{itemize}
\end{aretenir}

\courssub{Functions of Commercial Banks}

Commercial banks perform three core functions, supplemented by a growing array of auxiliary services.

\begin{propriete}[Core Functions of Commercial Banks]
\begin{enumerate}
\item \textbf{Accepting deposits} – providing safe custody for funds and paying interest on savings and time deposits. This mobilises idle balances into the formal financial system.
\item \textbf{Lending (Credit granting)} – granting \cle{loans} (medium/long term, e.g., equipment, housing), \cle{overdrafts} (short-term current account advances up to an agreed limit) and \cle{discounting bills of exchange} (providing immediate cash against future payment promises, deducting interest).
\item \textbf{Money transmission (Payment services)} – operating the payments system through \cle{cheques}, \cle{electronic transfers} (virements), \cle{standing orders} (ordres permanents), \cle{direct debits} (prélèvements) and, increasingly, \cle{mobile banking} platforms (e.g., Orange Money, MTN MoMo partnerships with banks).
\end{enumerate}
\end{propriete}

\begin{aretenir}[Other Services]
\begin{itemize}
\item Safe custody of valuables and documents (safe deposit boxes).
\item Foreign exchange dealing (buying/selling foreign currency, trade finance: letters of credit, documentary collections).
\item Advisory services (investment advice, project appraisal, feasibility studies).
\item Automated Teller Machines (ATMs), POS terminals, and digital banking apps (internet banking, USSD).
\item Issuance of bank guarantees and performance bonds.
\end{itemize}
\end{aretenir}

\courssub{The Process of Credit Creation: From Theory to Practice}

When a commercial bank grants a loan, it does not merely hand out existing deposits; it \emph{creates a new deposit} in the borrower's account. This new deposit becomes part of the money supply (M1/M2). The process repeats across the banking system, generating multiple rounds of deposit creation. This is the \cle{credit creation} (or \cle{deposit multiplication}) process.

\begin{propriete}[Cash Ratio (Reserve Requirement)]
The \cle{cash ratio} (or \cle{reserve ratio}, \cle{liquidity ratio} in some textbooks) $c$ is the fraction of total deposits that a bank must hold as liquid reserves (cash in vault + balances at the central bank). In the CEMAC zone, the regulatory minimum reserve requirement is set by BEAC/COBAC and differs by deposit category (e.g., sight deposits vs. term deposits). For pedagogical simplicity, we often assume a single constant ratio $c$ applied to all deposits.
\end{propriete}

\begin{propriete}[Credit (Deposit) Multiplier — Simple Model]
If the banking system receives an initial fresh deposit $\Delta D_0$ (new reserves) and all banks maintain the same cash ratio $c$, with \textbf{no leakages} (public holds no currency, banks hold no excess reserves), the total increase in deposits $\Delta D$ is given by the \cle{credit multiplier}:
\[
\Delta D = \frac{1}{c} \times \Delta D_0.
\]
The total new credit (loans) created equals $\Delta D - \Delta D_0$.
\end{propriete}

\begin{methode}[Computing Total Credit Creation (Simple Multiplier)]
\begin{enumerate}
\item Identify the initial injection of reserves (or fresh deposit) $\Delta D_0$ and the cash ratio $c$ (expressed as a decimal, e.g., $10\% = 0.10$).
\item Compute the multiplier $m = 1/c$.
\item Total new deposits $\Delta D = m \times \Delta D_0$.
\item Total new credit (loans) created $= \Delta D - \Delta D_0$.
\end{enumerate}
\end{methode}

\begin{exoresolu}[Multiple Deposit Expansion — Round-by-Round Analysis]
\textbf{Statement:} Suppose the banking system in Cameroon receives a fresh deposit of \SI{100}{million} CFA francs (e.g., from cocoa export proceeds deposited by a marketing board) and the regulatory cash ratio is $20\%$. Assuming no leakages, show the first three rounds of credit creation and calculate (a) the credit multiplier, (b) the total deposits eventually created, and (c) the total new loans granted.

\tcblower
\textbf{Detailed solution:}

\textbf{Round-by-round expansion (c = 0.20):}

\begin{center}
\begin{tabular}{|c|c|c|c|c|}
\hline
\textbf{Round} & \textbf{New Deposit Received} & \textbf{Reserves Kept (c $\times$ Deposit)} & \textbf{New Loans Made (1-c) $\times$ Deposit} & \textbf{Cumulative Deposits} \\
\hline
0 (Initial) & 100.00 & 20.00 & 80.00 & 100.00 \\
1 & 80.00 & 16.00 & 64.00 & 180.00 \\
2 & 64.00 & 12.80 & 51.20 & 244.00 \\
3 & 51.20 & 10.24 & 40.96 & 295.20 \\
$\vdots$ & $\vdots$ & $\vdots$ & $\vdots$ & $\vdots$ \\
\hline
\textbf{Total ($\infty$)} & \textbf{500.00} & \textbf{100.00} & \textbf{400.00} & \textbf{500.00} \\
\hline
\end{tabular}
\end{center}

\begin{enumerate}
\item Cash ratio $c = 20\% = 0.20$. Multiplier $m = 1/0.20 = 5$.
\item Total deposits $\Delta D = 5 \times 100 = \SI{500}{million}$ CFA francs.
\item Total new loans $= \Delta D - \Delta D_0 = 500 - 100 = \SI{400}{million}$ CFA francs.
\end{enumerate}

\textbf{Interpretation:} The initial \SI{100}{million} CFA francs of reserves supports a five-fold expansion of the money supply, with \SI{400}{million} CFA francs of new credit injected into the economy. The final reserves in the system exactly equal the initial injection (100 million), all held as required reserves against the 500 million total deposits.
\end{exoresolu}

\begin{popfigure}
\begin{tikzpicture}[node distance=1.6cm, >=stealth, font=\small]
  % Nodes for rounds with numbers
  \node[draw, popblue, thick, rounded corners, text width=2.8cm, align=center] (R0) {Round 0\\Initial deposit\\$\Delta D_0 = 100$\\Reserves kept: 20\\Loans: 80};
  \node[draw, poporange, thick, rounded corners, text width=2.8cm, align=center, right of=R0] (R1) {Round 1\\New deposit: 80\\Reserves kept: 16\\Loans: 64\\Cumul. deposits: 180};
  \node[draw, popgreen, thick, rounded corners, text width=2.8cm, align=center, right of=R1] (R2) {Round 2\\New deposit: 64\\Reserves kept: 12.8\\Loans: 51.2\\Cumul. deposits: 244};
  \node[draw, poppurple, thick, rounded corners, text width=2.8cm, align=center, right of=R2] (R3) {Round $n \to \infty$\\Total deposits: 500\\Total reserves: 100\\Total loans: 400};

  \draw[->, thick] (R0) -- node[above, font=\tiny] {Loan of 80 becomes new deposit} (R1);
  \draw[->, thick] (R1) -- node[above, font=\tiny] {Loan of 64 becomes new deposit} (R2);
  \draw[->, thick] (R2) -- node[above, font=\tiny] {Process continues...} (R3);
\end{tikzpicture}
\legende{Credit creation through successive rounds of lending and redepositing (cash ratio $c=0.20$, initial deposit 100 million CFAF).}
\end{popfigure}

\courssub{Credit Creation and the Bank's Balance Sheet}

Understanding credit creation requires looking at the bank's balance sheet. A balance sheet shows \textbf{Assets} (what the bank owns or is owed) = \textbf{Liabilities} (what the bank owes) + \textbf{Capital} (owners' equity).

\begin{experience}[T-Account Analysis of a Single Bank Making a Loan]
Consider Bank A, which starts with reserves of 100 and no loans. Cash ratio $c=20\%$.

\textbf{Initial Balance Sheet (Bank A):}
\begin{center}
\begin{tabular}{|l|r|l|r|}
\hline
\textbf{Assets} & \textbf{(million CFAF)} & \textbf{Liabilities + Capital} & \textbf{(million CFAF)} \\
\hline
Reserves (at BEAC + vault) & 100 & Deposits (initial) & 100 \\
Loans & 0 & Capital & 0 \\
\hline
\textbf{Total} & \textbf{100} & \textbf{Total} & \textbf{100} \\
\hline
\end{tabular}
\end{center}

Bank A can lend up to 80 (keeping 20 as required reserves). It grants a loan of 80 to a customer, crediting the customer's current account.

\textbf{After Loan of 80 (Bank A):}
\begin{center}
\begin{tabular}{|l|r|l|r|}
\hline
\textbf{Assets} & & \textbf{Liabilities + Capital} & \\
\hline
Reserves & 100 & Deposits (initial 100 + new 80) & 180 \\
Loans & 80 & Capital & 0 \\
\hline
\textbf{Total} & \textbf{180} & \textbf{Total} & \textbf{180} \\
\hline
\end{tabular}
\end{center}

\textbf{Observation:} The loan (asset) and the new deposit (liability) are created \emph{simultaneously}. The bank's reserves have not changed; only their composition (required vs. excess) has changed. Required reserves are now $0.20 \times 180 = 36$; excess reserves are $100 - 36 = 64$. The borrower spends the 80, the recipient deposits it in Bank B, and the process continues.
\end{experience}

\courssub{Limitations to Credit Creation (Leakages and Constraints)}

The theoretical multiplier $1/c$ assumes ideal conditions. In practice, several \cle{leakages} and constraints reduce the actual expansion.

\begin{aretenir}[Main Limitations to Credit Creation]
\begin{itemize}
\item \textbf{Currency leakage (Cash drain)} – the public holds part of the new money as currency (notes and coins) rather than redepositing it. If the currency-deposit ratio is $k$, this reduces the multiplier.
\item \textbf{Excess reserves} – banks may voluntarily hold reserves above the minimum (precautionary motive, uncertainty, lack of creditworthy borrowers). Let $e$ be the excess reserve ratio.
\item \textbf{Demand for loans} – credit creation requires creditworthy borrowers willing to borrow at the prevailing interest rate. In a recession, loan demand may be weak (\emph{pushing on a string}).
\item \textbf{Central bank regulations} – BEAC/COBAC may impose higher reserve requirements, capital adequacy ratios (Basel II/III standards), concentration limits (large exposure limits), or sectoral credit guidelines.
\item \textbf{Liquidity requirements} – banks must maintain a liquidity ratio (liquid assets / short-term liabilities) which can bind before the cash ratio.
\item \textbf{Non-bank financial intermediaries} – funds diverted to MFIs, money market funds, or capital markets do not feed the banking multiplier.
\end{itemize}
\end{aretenir}

\begin{propriete}[Money Multiplier with Leakages (Extended Model)]
When the public holds currency (currency-deposit ratio $k$) and banks hold excess reserves (excess reserve ratio $e$), the \cle{money multiplier} becomes:
\[
m = \frac{1 + k}{c + e + k}.
\]
If $k = 0$ and $e = 0$, this reduces to the simple credit multiplier $1/c$. The total money supply $M = m \times H$, where $H$ (high-powered money / monetary base) = Currency + Reserves.
\end{propriete}

\begin{attention}[Common Mistake]
Many students believe that \emph{``banks lend out the deposits they receive''}. In reality, \textbf{banks create deposits when they lend}. The loan appears as an asset, and a matching deposit (liability) is created simultaneously. The cash ratio only limits the \emph{multiple} of this creation, not the act itself. Also, do not confuse the \emph{cash ratio} (reserves/deposits) with the \emph{capital adequacy ratio} (capital/risk-weighted assets) — the latter constrains lending based on equity, not reserves.
\end{attention}

\begin{exoresolu}[Credit Multiplier with Currency Leakage]
\textbf{Statement:} Assume an initial deposit of \SI{200}{million} CFA francs, a cash ratio of $15\%$, and a currency leakage of $10\%$ of each new deposit (the public withdraws this fraction as cash). Banks hold no excess reserves. Compute the total deposit creation, total money supply change, and total new loans.

\tcblower
\textbf{Detailed solution:}
Here, the initial injection is a \emph{deposit}, so the monetary base $H$ does not increase by the full 200 million (part of it is just a shift from currency to deposits). However, if we treat the 200 million as \textbf{new reserves} injected into the banking system (e.g., BEAC open market purchase), and the public maintains a currency-deposit ratio $k = 0.10$, with $c = 0.15$, $e = 0$:

Effective denominator $= c + e + k = 0.15 + 0 + 0.10 = 0.25$.
Money multiplier $m = \frac{1 + k}{c + e + k} = \frac{1.10}{0.25} = 4.4$.

Total money supply increase $\Delta M = m \times \Delta H = 4.4 \times 200 = \SI{880}{million}$ CFA francs.
Of this, currency held by public $\Delta C = k \times \Delta D$ and deposits $\Delta D = \frac{1}{1+k} \Delta M = \frac{1}{1.10} \times 880 = \SI{800}{million}$ CFA francs.
Total new loans $\approx \Delta D - \Delta D_0$ (if initial deposit was part of $\Delta H$) or simply the increase in bank credit.
\textbf{Simpler approach (as per hint):} If the 200 million is an initial \emph{deposit} and leakage applies to each round, effective leakage-adjusted ratio $= c + k = 0.15 + 0.10 = 0.25$. Deposit multiplier $= 1/0.25 = 4$. Total deposits $= 4 \times 200 = \SI{800}{million}$ CFA francs. Total new loans $= 800 - 200 = \SI{600}{million}$ CFA francs. Currency in circulation increases by $k \times 800 = \SI{80}{million}$.
\end{exoresolu}

\courssub{Role of Commercial Banks in Cameroon's Economic Development}

Commercial banks are the engine of financial intermediation in Cameroon. Their contribution can be summarised under three headings.

\begin{aretenir}[Key Contributions to Cameroon's Economy]
\begin{itemize}
\item \textbf{Mobilising savings} – by offering safe, interest-bearing accounts (savings books, term deposits), banks channel household and corporate savings into the formal financial system, reducing hoarding and enabling capital formation.
\item \textbf{Financing investment} – loans fund agriculture (cocoa, coffee, cotton, banana), industry (SONARA refinery upgrades, aluminium smelting), infrastructure (roads, energy: Nachtigal, Lom Pangar), SMEs and housing. The \cle{credit multiplier} amplifies the impact of each unit of reserves.
\item \textbf{Facilitating trade} – payment instruments (cheques, electronic transfers, letters of credit, documentary collections) lower transaction costs for domestic and cross-border trade within CEMAC and under the AfCFTA.
\item \textbf{Monetary policy transmission} – banks are the channel through which BEAC's policy rates (taux d'appel d'offres, taux de pénalité) influence market interest rates and credit conditions.
\item \textbf{Financial inclusion} – expanding branch networks, agency banking, and mobile money partnerships bring unbanked populations (rural areas, informal sector) into the formal system.
\end{itemize}
\end{aretenir}

\courssub{Structure of the Cameroonian Banking System}

The following diagram shows the hierarchical and regulatory structure, with the BEAC at the apex.

\begin{popfigure}
\begin{tikzpicture}[node distance=1.4cm, >=stealth, font=\small]
  % Central Bank
  \node[draw, popdark, thick, rounded corners, fill=popblueL, text width=3.2cm, align=center] (BEAC) {BEAC\\ (Banque des États de l'Afrique Centrale)\\ Monetary policy, currency issue,\\ lender of last resort, foreign reserves};

  % Regulatory body
  \node[draw, popdark, thick, rounded corners, fill=popgreenL, text width=3.2cm, align=center, below of=BEAC] (COBAC) {COBAC\\ (Commission Bancaire de l'Afrique Centrale)\\ Prudential regulation, licensing,\\ supervision, resolution};

  % Commercial banks
  \node[draw, poporange, thick, rounded corners, fill=poporange!20, text width=3.2cm, align=center, left of=COBAC, xshift=-2.8cm] (Comm) {Commercial Banks (Universal)\\ Afriland First Bank, SGC, BICEC,\\ Ecobank, UBA, Standard Chartered,\\ SCB, CBC, NFC, etc.};

  % Microfinance
  \node[draw, poppurple, thick, rounded corners, fill=poppurple!20, text width=3.2cm, align=center, right of=COBAC, xshift=2.8cm] (MFI) {Microfinance Institutions (MFIs)\\ CAMCCUL, ACEP Cameroun, UCEC,\\ MC2, IFCE, etc.};

  % Specialised
  \node[draw, poppink, thick, rounded corners, fill=poppink!20, text width=3.2cm, align=center, below of=COBAC, yshift=-1.2cm] (Spec) {Specialised Institutions\\ Crédit Foncier (housing), SNI\\ (development finance)};

  % Public / Non-bank public
  \node[draw, popblue, thick, rounded corners, fill=popblue!10, text width=3.2cm, align=center, below of=Spec, yshift=-1.2cm] (Public) {Public / Non-bank Sector\\ Households, Firms, Government,\\ Public enterprises};

  % Arrows
  \draw[->, thick] (BEAC) -- node[left, font=\tiny] {sets policy rate, refinancing} (COBAC);
  \draw[->, thick] (COBAC) -- node[above left, font=\tiny] {licenses, supervises} (Comm);
  \draw[->, thick] (COBAC) -- node[above right, font=\tiny] {licenses, supervises} (MFI);
  \draw[->, thick] (COBAC) -- node[left, font=\tiny] {supervises} (Spec);
  \draw[<->, thick] (Comm) -- node[above, font=\tiny] {deposits / loans / payments} (Public);
  \draw[<->, thick] (MFI) -- node[above, font=\tiny] {micro-savings / micro-credit} (Public);
  \draw[<->, thick] (Spec) -- node[right, font=\tiny] {long-term finance} (Public);
  \draw[->, dashed, thick] (BEAC) -- node[left, font=\tiny] {reserves, SYSTAC settlement} (Comm);
  \draw[->, dashed, thick] (BEAC) -- node[right, font=\tiny] {refinancing window} (MFI);
  \draw[->, dashed, thick] (BEAC) -- node[left, font=\tiny] {refinancing} (Spec);
\end{tikzpicture}
\legende{Structure of the Cameroonian/CEMAC banking system under BEAC/COBAC supervision.}
\end{popfigure}

\begin{savaistu}[Key Fact: CEMAC Monetary Union]
Commercial banks in Cameroon operate within the \cle{CEMAC} (Communauté Économique et Monétaire de l'Afrique Centrale) monetary union, comprising 6 countries: Cameroon, Central African Republic, Chad, Congo, Equatorial Guinea, Gabon. They hold their reserves at the BEAC, settle interbank payments through the \cle{SYSTAC} (Système de Télécompensation de l'Afrique Centrale) and are subject to \cle{COBAC} prudential rules (capital adequacy $\ge 9.5\%$, liquidity ratio $\ge 100\%$, concentration limits). The CFA franc's fixed parity with the euro (1 EUR = 655.957 XAF) provides exchange-rate stability but ties monetary policy to the European Central Bank; BEAC's policy rate typically shadows the ECB's rate.
\end{savaistu}

\begin{exercice}[Practice: Credit Multiplier with Excess Reserves]
Assume the banking system receives a fresh reserve injection of \SI{500}{million} CFA francs. The required cash ratio is $10\%$, banks choose to hold excess reserves of $5\%$ of deposits, and the public's currency-deposit ratio is $20\%$. 
\begin{enumerate}
\item Calculate the money multiplier.
\item Calculate the total increase in money supply (M2).
\item Calculate the total increase in bank deposits.
\item Calculate the total increase in currency in circulation.
\end{enumerate}
\end{exercice}

\begin{corrige}[Exercise 1]
Given: $c = 0.10$, $e = 0.05$, $k = 0.20$, $\Delta H = 500$ million.
\begin{enumerate}
\item Money multiplier $m = \frac{1 + k}{c + e + k} = \frac{1.20}{0.10 + 0.05 + 0.20} = \frac{1.20}{0.35} \approx 3.4286$.
\item $\Delta M = m \times \Delta H = 3.4286 \times 500 \approx \SI{1714.3}{million}$ CFA francs.
\item $\Delta D = \frac{1}{1+k} \Delta M = \frac{1}{1.20} \times 1714.3 \approx \SI{1428.6}{million}$ CFA francs.
\item $\Delta C = k \times \Delta D = 0.20 \times 1428.6 \approx \SI{285.7}{million}$ CFA francs.
Check: $\Delta C + \Delta D = 285.7 + 1428.6 = 1714.3 = \Delta M$. Reserves increase = $\Delta H - \Delta C = 500 - 285.7 = 214.3$ million (required reserves $= c \Delta D = 142.9$, excess reserves $= e \Delta D = 71.4$, total $= 214.3$). Consistent.
\end{enumerate}
\end{corrige}

\courssub{Summary}

\begin{aretenir}[Chapter Essentials: Commercial Banks and Credit Creation]
\begin{itemize}
\item A \cle{commercial bank} is a deposit-taking, loan-making financial intermediary licensed by COBAC. It creates \cle{bank deposits} (money) when it grants credit.
\item The \cle{credit multiplier} (simple model) $= 1/c$, where $c$ is the cash/reserve ratio. With leakages (currency drain $k$, excess reserves $e$), the \cle{money multiplier} $= \frac{1+k}{c+e+k}$.
\item Credit creation is constrained by: currency leakage, excess reserves, loan demand, capital adequacy, liquidity rules, and BEAC/COBAC regulations.
\item In Cameroon/CEMAC, commercial banks (Afriland, SGC, BICEC, Ecobank, etc.) are the main conduit for savings mobilisation, investment finance, trade facilitation, and monetary policy transmission under BEAC oversight.
\item The balance sheet identity (Assets = Liabilities + Capital) shows that every loan creates a matching deposit; reserves are not "lent out" but serve as the basis for multiple deposit expansion.
\end{itemize}
\end{aretenir}

\coursec{The Balance Sheet of a Commercial Bank}

In the previous section we saw how commercial banks create credit by lending out a multiple of their cash base. But how can we actually \emph{see} what a bank owns and what it owes at a given moment? The answer is the bank's \cle{balance sheet} — a financial photograph of the bank taken at a particular date. Every commercial bank in Cameroon, from Afriland First Bank to SGBC or BICEC, must publish such a statement under COBAC (Commission Bancaire de l'Afrique Centrale) regulations. Understanding it is essential for the GCE, where you may be asked to construct one from a list of items, interpret an extract, or compute key ratios.

\courssub{Meaning and Purpose of a Bank Balance Sheet}

Imagine you run a small provision store in Mokolo market, Yaoundé. At the end of the month you sit down and write two lists: on one side, everything you \emph{own} (your stock of goods, the cash in your drawer, the money customers owe you); on the other side, everything you \emph{owe} (the loan from your brother, the unpaid supplier's bill). A bank does exactly the same thing, but on a much larger scale and with stricter rules governed by the CEMAC banking charter.

\begin{definition}[Balance Sheet, Assets and Liabilities]
A \cle{balance sheet} is a financial statement showing the \cle{assets} and \cle{liabilities} of a firm at a particular point in time (usually the end of the financial year).
\begin{itemize}
\item \textbf{Assets} are everything the bank \emph{owns} or is \emph{owed}: cash, loans granted to customers, investments, buildings.
\item \textbf{Liabilities} are everything the bank \emph{owes} to others: shareholders' capital, reserves, and above all customers' deposits.
\end{itemize}
By construction, \textbf{Total Assets = Total Liabilities}.
\end{definition}

The purposes of the balance sheet are:
\begin{itemize}
\item to show the \textbf{financial position} of the bank at a given date;
\item to reveal how the bank has used the funds entrusted to it (the structure of its assets);
\item to allow shareholders, depositors, the central bank (BEAC in the CEMAC zone) and the public to judge the bank's \textbf{solvency} (ability to meet long-term obligations) and \textbf{liquidity} (ability to meet short-term cash demands);
\item to permit comparison of a bank's performance over time or with other banks.
\end{itemize}

\begin{attention}[Deposits are Liabilities, not Assets]
The most common GCE mistake is to classify \textbf{customers' deposits as assets}. Remember: a deposit is money that \emph{belongs to the customer} and that the bank must repay on demand or at an agreed date. From the \emph{bank's} point of view, a deposit is a \textbf{debt}, hence a \textbf{liability}. Conversely, a \textbf{loan granted to a customer} is money the customer owes the bank — it is an \textbf{asset} of the bank. Always reason from the bank's standpoint, never from the customer's.
\end{attention}

\courssub{The Liabilities Side: Sources of Funds}

The liabilities side shows where the bank obtained its funds. The main items, listed in order of permanence (most permanent first), are:

\begin{itemize}
\item \textbf{Share capital}: the money contributed by the shareholders when the bank was created or when capital was increased. It is a liability because it ultimately belongs to the shareholders.
\item \textbf{Reserves}: accumulated profits which have not been distributed as dividends but kept in the bank to strengthen it (statutory reserve, general reserve, retained earnings). Under COBAC rules, banks must transfer a percentage of annual profits to a statutory reserve.
\item \textbf{Deposit liabilities}: the largest item for most banks. They are usually subdivided into:
\begin{itemize}
\item \emph{Current (demand) deposits}: repayable on demand, by cheque, transfer, or ATM withdrawal. Usually non-interest-bearing or very low interest.
\item \emph{Savings deposits}: withdrawable subject to short notice (or with penalties for frequent withdrawal), earning modest interest. In Cameroon, the \emph{livret d'épargne} is a classic example.
\item \emph{Fixed (time) deposits}: blocked for an agreed period (30 days, 90 days, 1 year, etc.), earning higher interest. Early withdrawal incurs a penalty.
\end{itemize}
\item \textbf{Borrowings}: funds the bank itself has borrowed, for example from the central bank (refinancing window, advances at the policy rate) or from other banks on the interbank market (money market).
\item \textbf{Other liabilities}: accrued expenses, provisions for bad debts, taxes payable, etc. (often grouped as ``Other liabilities'' in simplified GCE questions).
\end{itemize}

\courssub{The Assets Side: Uses of Funds — Order of Liquidity and Profitability}

The assets side shows how the bank has \emph{used} its funds. Assets are conventionally listed in a very specific order, which examiners expect you to know and reproduce exactly.

\begin{propriete}[Ordering of Assets on the Balance Sheet]
On a commercial bank's balance sheet, assets are arranged in \textbf{descending order of liquidity} (from the most liquid to the least liquid) and, equivalently, in \textbf{ascending order of profitability} (from the least profitable to the most profitable). This ordering is a convention you must reproduce exactly in the examination.
\end{propriete}

The typical order, with Cameroonian/CEMAC context, is:

\begin{enumerate}
\item \textbf{Cash in hand and at the central bank}: notes and coins in the vault (\emph{caisse}) plus the balance held at BEAC (compte de réserve obligatoire and compte courant). Perfectly liquid, but earns \emph{nothing} (or very little on the reserve account). This item satisfies the \cle{legal reserve requirement} (currently a percentage of deposit liabilities set by BEAC).
\item \textbf{Money at call and short notice}: very short-term loans to other banks or financial institutions, recallable within hours or days (typically overnight to 7 days). Highly liquid, very low return (interbank rate).
\item \textbf{Bills discounted}: treasury bills (Bons du Trésor) and eligible trade bills (effets de commerce) bought at a discount before maturity. Liquid (they can be rediscounted at the central bank or sold on the secondary market), modest return.
\item \textbf{Investments (Securities)}: medium- and long-term securities such as government bonds (Obligations d'État), corporate bonds, and equity holdings. Less liquid than bills, but they pay regular interest/coupons. In CEMAC, banks hold significant amounts of sovereign securities.
\item \textbf{Loans and advances}: overdrafts (découverts), short-term loans (crédits à court terme), medium- and long-term loans (crédits d'équipement) granted to customers. The \emph{most profitable} asset class but illiquid — the bank cannot demand instant repayment of a five-year equipment loan. This is where credit creation materialises.
\item \textbf{Fixed assets (Premises, equipment, intangible assets)}: the bank's buildings, branches, IT systems, vehicles, software licences. Least liquid of all; they earn no direct financial return but are necessary for operations.
\end{enumerate}

\begin{popfigure}
\begin{center}
\begin{tabular}{|l|r|l|r|}
\hline
\rowcolor{popblueL} \textbf{Liabilities} & \textbf{FCFA m} & \textbf{Assets} & \textbf{FCFA m} \\
\hline
Share capital & 5\,000 & Cash in hand and at central bank & 4\,000 \\
Reserves & 2\,000 & Money at call and short notice & 3\,000 \\
Current deposits & 20\,000 & Bills discounted & 5\,000 \\
Savings deposits & 8\,000 & Investments (securities) & 8\,000 \\
Fixed deposits & 6\,000 & Loans and advances & 19\,000 \\
Borrowings & 2\,000 & Premises and equipment & 4\,000 \\
\hline
\textbf{Total} & \textbf{43\,000} & \textbf{Total} & \textbf{43\,000} \\
\hline
\end{tabular}
\end{center}
\legende{Model balance sheet of a commercial bank. Assets are ranked from the most liquid (cash, top) to the least liquid and most profitable (loans, premises, bottom). Totals on both sides are equal.}
\end{popfigure}

\courssub{The Fundamental Conflict: Liquidity versus Profitability}

Here lies the eternal dilemma of the banker. \textbf{Liquidity} is the ease with which an asset can be converted into cash without loss of value; \textbf{profitability} is the income an asset yields per period. The two move in \emph{opposite directions} along the asset spectrum:

\begin{itemize}
\item If the bank keeps too much cash and call money, it is perfectly safe and can meet any withdrawal — but it earns almost no profit, and shareholders will complain about low return on equity.
\item If the bank lends nearly everything as long-term loans, profits are high — but a wave of withdrawals could force it into a \emph{liquidity crisis} and even insolvency, because loans cannot be recalled instantly without severe loss.
\end{itemize}

Good bank management therefore consists of \textbf{striking a balance}: holding enough liquid assets to honour withdrawals (this is also why the central bank imposes \cle{reserve requirements} and \cle{liquidity ratios}) while lending enough to remain profitable. The BEAC monitors this through prudential ratios (liquidity ratio, solvency ratio, concentration risk limits).

\begin{popfigure}
\begin{center}
\begin{tikzpicture}[scale=1.0]
\draw[->, thick, popdark] (0,0) -- (10.5,0) node[below left=2pt and -30pt, text=popdark, font=\small] {Increasing profitability $\rightarrow$};
\draw[->, thick, popdark] (0,0) -- (0,3.4) node[rotate=90, above, midway, font=\small, text=popdark] {Liquidity};
\draw[very thick, popblue] (0.6,2.9) -- (9.9,0.4);
\node[font=\small, text=popblue] at (5.2,2.15) {Liquidity falls as profitability rises};
\foreach \x/\lab in {0.9/{Cash}, 2.7/{Call money}, 4.5/{Bills}, 6.3/{Investments}, 8.1/{Loans}, 9.7/{Premises}}{
  \fill[poporange] (\x,0) circle (2pt);
  \node[below, font=\footnotesize, text=popdark] at (\x,-0.05) {\lab};
}
\node[font=\small\itshape, text=popmuted] at (8.6,2.9) {Most liquid, least profitable};
\node[font=\small\itshape, text=popmuted] at (8.6,0.9) {Least liquid, most profitable};
\end{tikzpicture}
\end{center}
\legende{The liquidity--profitability trade-off: as the bank moves along the spectrum of assets from cash towards loans and premises, liquidity falls while profitability rises.}
\end{popfigure}

\begin{attention}[The Balance Sheet Must Always Balance]
Total assets must \emph{always} equal total liabilities. If your two totals differ in an examination question, you have either misclassified an item or made an arithmetic slip — go back and check. Note also that share capital and reserves count as liabilities: they are what the bank owes to its owners. This equality is not a coincidence; every franc the bank holds (asset) was obtained from someone (liability).
\end{attention}

\courssub{Constructing a Balance Sheet: Method and Worked Example}

\begin{methode}[Checklist for Classifying Any Item]
To place any item correctly, ask yourself two questions:
\begin{enumerate}
\item \textbf{Does the bank own it, or is it owed to the bank?} If yes $\rightarrow$ \textbf{Asset} (cash, call money, bills, investments, loans granted, premises).
\item \textbf{Does the bank owe it to someone else?} If yes $\rightarrow$ \textbf{Liability} (share capital, reserves, all deposits, borrowings from the central bank or other banks).
\item Place assets in descending order of liquidity: cash $\rightarrow$ call money $\rightarrow$ bills $\rightarrow$ investments $\rightarrow$ loans $\rightarrow$ premises.
\item Add up each side and \textbf{verify that the two totals are equal}.
\end{enumerate}
\end{methode}

\begin{exoresolu}[Constructing a Balance Sheet from Given Items]
The following items were extracted from the books of Wouri Bank Ltd on 31 December (figures in millions of FCFA): loans and advances 12\,000; share capital 4\,000; cash in hand 1\,500; fixed deposits 5\,000; investments 6\,000; reserves 1\,500; bills discounted 2\,500; current deposits 10\,000; premises 3\,000; money at call 2\,000; savings deposits 4\,500; cash at the central bank 1\,000. \textbf{Required:} prepare the balance sheet of Wouri Bank Ltd.
\tcblower
\textbf{Solution.} Classify each item using the checklist: deposits, capital and reserves are liabilities; cash, call money, bills, investments, loans and premises are assets. Arrange assets in descending order of liquidity. Note that ``cash in hand'' and ``cash at the central bank'' are combined as the first asset item in the standard presentation, but may be shown separately for clarity.
\begin{center}
\begin{tabular}{|l|r|l|r|}
\hline
\rowcolor{popblueL} \textbf{Liabilities} & \textbf{FCFA m} & \textbf{Assets} & \textbf{FCFA m} \\
\hline
Share capital & 4\,000 & Cash in hand & 1\,500 \\
Reserves & 1\,500 & Cash at central bank & 1\,000 \\
Current deposits & 10\,000 & Money at call & 2\,000 \\
Savings deposits & 4\,500 & Bills discounted & 2\,500 \\
Fixed deposits & 5\,000 & Investments & 6\,000 \\
 &  & Loans and advances & 12\,000 \\
 &  & Premises & 3\,000 \\
\hline
\textbf{Total} & \textbf{25\,000} & \textbf{Total} & \textbf{25\,000} \\
\hline
\end{tabular}
\end{center}
Check: liabilities $= 4\,000 + 1\,500 + 10\,000 + 4\,500 + 5\,000 = 25\,000$; assets $= 1\,500 + 1\,000 + 2\,000 + 2\,500 + 6\,000 + 12\,000 + 3\,000 = 25\,000$. The balance sheet balances. $\blacksquare$
\end{exoresolu}

\courssub{Effects of Transactions on the Balance Sheet}

A balance sheet is not static: every banking operation changes it, but the equality Assets = Liabilities is always preserved. Understanding these mechanics is crucial for analysing credit creation and monetary policy transmission.

\begin{exemplebox}[Three Typical Transactions]
Take a simplified bank with cash of 2\,000 and deposits of 2\,000 (all figures in FCFA m). Assume no reserve requirement for simplicity.
\begin{enumerate}
\item \textbf{A new cash deposit of 500 by a customer}: cash (asset) rises to 2\,500 and deposits (liability) rise to 2\,500. Both sides increase by the same amount. The bank's liquidity improves.
\item \textbf{Granting a loan of 800 by crediting the customer's current account}: loans (asset) rise by 800 and current deposits (liability) rise by 800 — this is precisely \emph{credit creation}: the bank creates a deposit out of nothing by writing a loan. Cash is unchanged, but liquidity falls (cash/deposit ratio drops).
\item \textbf{A cash withdrawal of 300 by a depositor}: cash (asset) falls by 300 and deposits (liability) fall by 300. If withdrawals exceed available cash and quickly realisable assets (cash + call money + bills), the bank faces a liquidity crisis — the practical consequence of the liquidity--profitability conflict.
\end{enumerate}
In each case, the balance sheet still balances after the transaction.
\end{exemplebox}

\begin{savaistu}[Reading a Real Cameroonian Bank Balance Sheet]
Published balance sheets of Cameroonian banks (Afriland First Bank, BICEC, SGBC, Ecobank Cameroun, UBA Cameroun...) follow exactly the logic studied here, although under COBAC/CEMAC accounting rules (Plan Comptable Bancaire Unifié) the vocabulary is more technical: ``customer deposits'' appear as \emph{amounts due to customers} (dettes envers la clientèle), loans as \emph{amounts due from customers} (créances sur la clientèle), and balances at the central bank as \emph{amounts due from the BEAC} (créances sur la BEAC). In the GCE, an extract may be given and you may be asked to identify which items are assets or liabilities, or to compute ratios such as the \textbf{loan-to-deposit ratio} (loans and advances / total deposits) or the \textbf{liquidity ratio} (liquid assets / short-term liabilities). The classification rules you have learned are all you need.
\end{savaistu}

\begin{exercice}[Balance Sheet Construction and Verification]
From the following items of Nkam Bank Ltd (FCFA m), prepare its balance sheet at 31 December: premises 2\,500; current deposits 12\,000; cash in hand 1\,000; loans and advances 9\,000; share capital 3\,000; bills discounted 1\,500; savings deposits 3\,500; investments 4\,000; reserves 1\,000; money at call 1\,500; fixed deposits 2\,000; cash at the central bank 500. State clearly why you classified each of the deposit items as you did.
\end{exercice}

\begin{corrige}[Exercise 1]
\textbf{Step 1: Classification.}
\begin{itemize}
\item \textbf{Liabilities (sources of funds):} Share capital 3\,000; Reserves 1\,000; Current deposits 12\,000; Savings deposits 3\,500; Fixed deposits 2\,000. Total Liabilities $= 21\,500$.
\item \textbf{Assets (uses of funds), in descending order of liquidity:} Cash in hand 1\,000; Cash at central bank 500; Money at call 1\,500; Bills discounted 1\,500; Investments 4\,000; Loans and advances 9\,000; Premises 2\,500. Total Assets $= 20\,000$.
\end{itemize}

\textbf{Step 2: Verification.} The totals differ by 1\,500 million FCFA. In a real examination, this signals an error: either a figure was miscopied, an item was omitted from the list, or an item was misclassified. \emph{Do not force the balance by inventing a figure.} Instead, state: ``The totals do not match (21\,500 vs 20\,000), indicating a discrepancy in the provided data. Assuming the data is correct as given, the balance sheet does not balance.'' (In this specific exercise, the discrepancy is deliberate to test your checking reflex.)

\textbf{Step 3: Justification for deposit classification.} All deposit items (current, savings, fixed) are classified as \textbf{liabilities} because they represent sums the bank \emph{owes} to its customers. The bank has a legal obligation to repay:
\begin{itemize}
\item Current deposits: on demand (cheque, transfer, ATM).
\item Savings deposits: on demand or after short notice, per the savings contract.
\item Fixed deposits: at the agreed maturity date (or earlier with penalty).
\end{itemize}
From the bank's perspective, these are debts — hence liabilities. $\blacksquare$
\end{corrige}

\begin{aretenir}[Key Points for the GCE]
\begin{itemize}
\item A balance sheet is a statement of assets and liabilities at a point in time; \textbf{Total Assets = Total Liabilities}, always.
\item Liabilities (sources): share capital, reserves, deposits (current, savings, fixed), borrowings. \textbf{Deposits are liabilities} because the bank owes them to customers.
\item Assets (uses) are ranked in \textbf{descending order of liquidity} and \textbf{ascending order of profitability}: cash $\rightarrow$ call money $\rightarrow$ bills $\rightarrow$ investments $\rightarrow$ loans $\rightarrow$ premises.
\item The banker's dilemma is the \textbf{liquidity--profitability conflict}: the most profitable assets (loans) are the least liquid.
\item Every transaction (deposit, loan, withdrawal) changes both sides so that the balance sheet continues to balance; granting a loan by crediting an account is the mechanics of credit creation.
\item \textbf{Exam skills:} Construct a balance sheet from a jumbled list; identify assets vs liabilities in an extract; compute loan-to-deposit ratio, liquid asset ratio; explain the effect of a given transaction on the balance sheet.
\end{itemize}
\end{aretenir}

\newpage

\coursec{Integration activity}

\coursec{Integration Activity: Money, Banking and Inflation in Mbankomo}

\courssub{Aims of the Activity}

This integration activity brings together, in a single authentic case study, all the notions developed in this chapter: the functions of \cle{money}, the construction of a \cle{retail price index}, the measurement of \cle{inflation} and of changes in the \cle{value of money}, the \cle{quantity theory of money} ($MV = PT$), \cle{credit creation} by commercial banks through the credit multiplier, and the structure of a commercial bank's \cle{balance sheet}.

By the end of the activity, you should be able to:
\begin{itemize}
\item construct a base-weighted (Laspeyres) retail price index from raw price and quantity data;
\item compute an inflation rate and the corresponding percentage change in the value (purchasing power) of money;
\item apply the Fisher equation $MV = PT$ to predict a change in the price level and compare it with an observed index;
\item calculate the total deposit expansion and the total credit created by the banking system from a primary deposit, given a cash ratio;
\item classify financial items into assets and liabilities and present a balanced balance sheet;
\item use all these results to give a reasoned financial advice to an economic agent — the skill most valued in GCE essay questions.
\end{itemize}

\begin{methode}[How to Work]
Form groups of four to six students. Each group appoints a secretary and a presenter. Work through Tasks 1 to 7 in order (60--90 minutes), then prepare a five-minute oral presentation ending with your advice to the cooperative. Calculators are allowed; show all formulae and working, as required at the GCE Advanced Level.
\end{methode}

\courssub{Situation}

\textbf{Mbankomo Traders' Cooperative: Money, Prices and a Bank Loan.}\par
Mbankomo is a fast-growing town on the Yaound\'e--Mbalmayo road in the Centre Region of Cameroon. The \emph{Mbankomo Traders' Cooperative Society} (MBATRACO) groups 120 retailers who buy foodstuffs in the surrounding villages and resell them in the town market. At its annual general meeting in January 2025, the executive bureau must report on three worrying questions raised by members:
\begin{itemize}
\item ``Our money buys less and less in the market. By how much, exactly, have prices risen since 2022?''
\item ``The quantity of money circulating in our local economy has increased. Is that the cause of the price rise?''
\item ``The cooperative deposited 5,000,000 FCFA at \emph{Avenir Bank S.A.}, a commercial bank. The bank now proposes a loan to finance a storage warehouse. How much credit can banks really create, and should we save more or borrow?''
\end{itemize}

The bureau has gathered the following dossier.

\textbf{Document 1 -- Prices of the cooperative's reference basket.} The basket represents the typical monthly purchases of a member household. Quantities are those consumed in the base year 2022.

\begin{center}
\begin{tabularx}{\linewidth}{|l|c|c|c|X|}
\hline
\rowcolor{popblueL} \textbf{Item (unit)} & \textbf{Qty 2022} & \textbf{Price 2022 (FCFA)} & \textbf{Price 2024 (FCFA)} & \textbf{Unit} \\
\hline
Garri & 100 & 400 & 500 & per kg \\
\hline
Palm oil & 50 & 1\,000 & 1\,300 & per litre \\
\hline
Plantain & 200 & 250 & 300 & per kg \\
\hline
Fresh fish & 80 & 1\,500 & 1\,800 & per kg \\
\hline
\end{tabularx}
\end{center}

\textbf{Document 2 -- Money and output in the local economy.} A study by a local NGO estimates that between 2022 and 2024 the money supply circulating in the Mbankomo area rose from 200 million FCFA to 250 million FCFA, while the real volume of transactions (output of goods and services) rose by 5\%. The velocity of circulation of money is assumed constant over the period.

\textbf{Document 3 -- The cooperative's deposit.} In January 2024, MBATRACO deposited 5,000,000 FCFA in cash at Avenir Bank. All banks in the system keep a \cle{cash ratio} (cash reserve ratio) of 10\% of deposits and lend out the rest; all money lent is re-deposited in the banking system.

\textbf{Document 4 -- Financial items of Avenir Bank S.A. (31 December 2024, in millions of FCFA).}

\begin{center}
\begin{tabularx}{\linewidth}{|X|c|X|c|}
\hline
\rowcolor{popblueL} \textbf{Item} & \textbf{Amount} & \textbf{Item} & \textbf{Amount} \\
\hline
Customer deposits & 50 & Treasury bills & 7 \\
\hline
Share capital & 6 & Loans and advances to customers & 41 \\
\hline
Reserves & 2 & Investments in local firms & 4 \\
\hline
Borrowing from the central bank & 5 & Cash in vault & 2 \\
\hline
Other liabilities (creditors) & 1 & Balances at the central bank & 3 \\
\hline
Premises and equipment & 6 & Cheques in course of collection & 1 \\
\hline
\end{tabularx}
\end{center}

\begin{exercice}[Tasks]
\begin{enumerate}
\item \textbf{Cost of the basket.} Using the 2022 quantities, calculate the total cost of the reference basket (a) at 2022 prices; (b) at 2024 prices.
\item \textbf{Retail price index and inflation.} Construct the retail price index for 2024, base 100 in 2022. Deduce the rate of inflation between 2022 and 2024.
\item \textbf{Value of money.} Calculate the purchasing power of 1 FCFA of 2024 in terms of 2022 FCFA, and the percentage change in the value of money. Interpret the result for a trader holding cash.
\item \textbf{Quantity theory.} Using $MV = PT$ with $V$ constant, predict the percentage change in the price level implied by Document 2. Compare this prediction with the index obtained in Task 2 and suggest one reason for any difference.
\item \textbf{Credit creation.} From the 5,000,000 FCFA primary deposit and the 10\% cash ratio, calculate (a) the credit multiplier; (b) the total deposits the banking system can create; (c) the total new credit (loans) created. Show the first two rounds of the process.
\item \textbf{Balance sheet.} Classify the twelve items of Document 4 into assets and liabilities, present the balance sheet of Avenir Bank in proper form, and verify that it balances.
\item \textbf{Synthesis and advice.} In about ten lines, advise the cooperative: given the inflation measured, the behaviour of the value of money and the bank's capacity to lend, is it wiser for MBATRACO to keep its savings in cash, to deposit them, or to borrow now for the warehouse? Justify with your figures.
\end{enumerate}
\end{exercice}

\begin{exoresolu}[Task 1: Cost of the Basket]
\textbf{Statement:} Using the 2022 quantities as weights, compute the total cost of the basket at 2022 prices and at 2024 prices.

\tcblower
\textbf{Detailed Solution:}

The Laspeyres (base-weighted) cost of the basket is $\sum p_t q_0$ where $q_0$ are the 2022 quantities.

\begin{align*}
\text{Cost}_{2022} &= (100 \times 400) + (50 \times 1\,000) + (200 \times 250) + (80 \times 1\,500) \\
&= 40\,000 + 50\,000 + 50\,000 + 120\,000 \\
&= \mathbf{260\,000\ FCFA}. \\[6pt]
\text{Cost}_{2024} &= (100 \times 500) + (50 \times 1\,300) + (200 \times 300) + (80 \times 1\,800) \\
&= 50\,000 + 65\,000 + 60\,000 + 144\,000 \\
&= \mathbf{319\,000\ FCFA}.
\end{align*}

\begin{attention}[Examiner's Tip]
Always show the multiplication for each item. In the GCE, marks are awarded for the correct application of the formula $\sum p_1 q_0$ and $\sum p_0 q_0$, not just the final totals.
\end{attention}
\end{exoresolu}

\begin{exoresolu}[Task 2: Retail Price Index and Inflation Rate]
\textbf{Statement:} Construct the retail price index (RPI) for 2024 with base 100 in 2022, and deduce the inflation rate over the period.

\tcblower
\textbf{Detailed Solution:}

The Laspeyres price index formula is:
\[
\text{RPI}_{2024} = \frac{\sum p_{2024} q_{2022}}{\sum p_{2022} q_{2022}} \times 100.
\]

Substituting the totals from Task 1:
\[
\text{RPI}_{2024} = \frac{319\,000}{260\,000} \times 100 = 122.6923\ldots \approx \mathbf{122.7}\ \text{(base 100 in 2022)}.
\]

The inflation rate between 2022 and 2024 is the percentage increase in the index:
\[
\text{Inflation rate} = \frac{122.7 - 100}{100} \times 100\% = \mathbf{22.7\%}.
\]

\begin{aretenir}[Key Formula]
\[
\text{Inflation rate} = \frac{\text{Index}_{\text{current}} - \text{Index}_{\text{base}}}{\text{Index}_{\text{base}}} \times 100\%.
\]
Since the base index is always 100, this simplifies to $\text{Index}_{\text{current}} - 100$.
\end{aretenir}
\end{exoresolu}

\begin{exoresolu}[Task 3: Value of Money]
\textbf{Statement:} Calculate the purchasing power of 1 FCFA of 2024 in terms of 2022 FCFA, the percentage change in the value of money, and interpret for a trader holding cash.

\tcblower
\textbf{Detailed Solution:}

The purchasing power (value) of money is the reciprocal of the price index (expressed as a decimal):
\[
\text{Purchasing power of 2024 FCFA in 2022 terms} = \frac{100}{\text{RPI}_{2024}} = \frac{100}{122.7} \approx \mathbf{0.815\ FCFA}.
\]

This means 1 FCFA in 2024 buys only what 0.815 FCFA bought in 2022.

The percentage change in the value of money is:
\[
\frac{0.815 - 1}{1} \times 100\% = -18.5\% \quad \text{or} \quad \left(1 - \frac{100}{122.7}\right) \times 100\% \approx \mathbf{18.5\%\ fall}.
\]

\textbf{Interpretation:} A trader who kept 260,000 FCFA in a cash box since 2022 can now buy only about 81.5\% of the former basket. The real value of idle cash has fallen by 18.5\%. In inflationary times, holding money as cash is a losing strategy — it is a \cle{negative real return} asset.

\begin{attention}[Common Mistake]
Do \textbf{not} say ``the value of money fell by 22.7\%''. The fall in the value of money is measured by $1 - \frac{100}{\text{Index}}$, which is always smaller than the rise in prices. The two magnitudes are reciprocal, not equal:
\[
\text{\% fall in value} = \frac{\text{\% rise in prices}}{100 + \text{\% rise in prices}} \times 100\%.
\]
Here $\frac{22.7}{122.7} \times 100\% \approx 18.5\%$.
\end{attention}
\end{exoresolu}

\begin{exoresolu}[Task 4: Quantity Theory of Money]
\textbf{Statement:} Using $MV = PT$ with $V$ constant, predict the percentage change in the price level from Document 2. Compare with the measured index and suggest a reason for any difference.

\tcblower
\textbf{Detailed Solution:}

The Fisher equation is $MV = PT$. With $V$ constant, changes satisfy:
\[
\frac{\Delta M}{M} + \frac{\Delta V}{V} = \frac{\Delta P}{P} + \frac{\Delta T}{T} \quad \Rightarrow \quad \frac{\Delta P}{P} \approx \frac{\Delta M}{M} - \frac{\Delta T}{T}.
\]

From Document 2:
\begin{itemize}
\item Money supply $M$ rose from 200 million to 250 million FCFA: $\frac{\Delta M}{M} = \frac{50}{200} = 25\%$.
\item Real transactions $T$ rose by $5\%$.
\end{itemize}

Predicted price level change:
\[
\frac{P_{2024}}{P_{2022}} = \frac{1.25}{1.05} \approx 1.1905 \quad \Rightarrow \quad \text{Predicted inflation} \approx \mathbf{19.05\%}.
\]

\textbf{Comparison:} The quantity theory predicts $\approx 19\%$ inflation, while the measured RPI gives $22.7\%$. The prediction is close but understates the actual rise by about $3.7$ percentage points.

\textbf{Possible reasons for the gap:}
\begin{itemize}
\item The basket covers only four food items; it omits services, transport, imported goods, etc., which may have risen faster.
\item Velocity $V$ may not be perfectly constant (e.g., increased use of mobile money could raise $V$).
\item Supply-side shocks: Mbankomo depends on the Yaound\'e--Mbalmayo road; fuel price rises or road disruptions increase transport costs, pushing food prices above the purely monetary prediction.
\item The RPI uses fixed 2022 weights (Laspeyres), which tends to overstate inflation if consumers substitute away from goods whose prices rise most.
\end{itemize}

\begin{aretenir}[Quantity Theory in Exams]
At GCE level, you must:
\begin{itemize}
\item State the equation $MV = PT$ and define each variable.
\item Explain the assumptions ($V$ constant, $T$ determined by real factors).
\item Derive $\% \Delta P \approx \% \Delta M - \% \Delta T$.
\item Apply it numerically and evaluate its limitations (e.g., $V$ not constant, measurement issues, short-run vs long-run).
\end{itemize}
\end{aretenir}
\end{exoresolu}

\begin{exoresolu}[Task 5: Credit Creation]
\textbf{Statement:} From the 5,000,000 FCFA primary deposit and a 10\% cash ratio, calculate (a) the credit multiplier; (b) total deposits created; (c) total new credit created. Show the first two rounds.

\tcblower
\textbf{Detailed Solution:}

\textbf{(a) Credit multiplier:}
\[
k = \frac{1}{\text{cash ratio}} = \frac{1}{0.10} = \mathbf{10}.
\]

\textbf{(b) Total deposits the banking system can create:}
\[
\Delta D = \text{Primary deposit} \times k = 5\,000\,000 \times 10 = \mathbf{50\,000\,000\ FCFA}.
\]

\textbf{(c) Total new credit (loans) created:}
\[
\Delta L = \Delta D - \text{Primary deposit} = 50\,000\,000 - 5\,000\,000 = \mathbf{45\,000\,000\ FCFA}.
\]

\textbf{First two rounds of the process:}

\begin{center}
\begin{tabular}{|c|c|c|c|c|}
\hline
\rowcolor{popblueL} \textbf{Round} & \textbf{New Deposit (FCFA)} & \textbf{Cash Reserves (10\%)} & \textbf{New Loans (90\%)} & \textbf{Cumulative Deposits} \\
\hline
Initial & 5\,000\,000 & 500\,000 & 4\,500\,000 & 5\,000\,000 \\
\hline
1 & 4\,500\,000 & 450\,000 & 4\,050\,000 & 9\,500\,000 \\
\hline
2 & 4\,050\,000 & 405\,000 & 3\,645\,000 & 13\,550\,000 \\
\hline
\end{tabular}
\end{center}

The process continues until total deposits approach 50 million FCFA, total reserves equal the initial 5 million FCFA, and total loans equal 45 million FCFA.

\begin{propriete}[Credit Multiplier Formula]
For a single bank: it can lend $(1-r) \times \text{deposit}$.
For the \textbf{banking system} (all loans re-deposited):
\[
\text{Total deposits} = \frac{1}{r} \times \text{Primary deposit}, \quad
\text{Total credit} = \frac{1-r}{r} \times \text{Primary deposit}.
\]
Assumptions: no cash drain, no excess reserves, all loans re-deposited, constant cash ratio $r$.
\end{propriete}
\end{exoresolu}

\begin{exoresolu}[Task 6: Balance Sheet of Avenir Bank S.A.]
\textbf{Statement:} Classify the twelve items into assets and liabilities, present the balance sheet in proper form, and verify it balances.

\tcblower
\textbf{Detailed Solution:}

\textbf{Classification:}
\begin{itemize}
\item \textbf{Liabilities (what the bank owes / sources of funds):} Customer deposits (50), Share capital (6), Reserves (2), Borrowing from central bank (5), Other liabilities (1).
\item \textbf{Assets (uses of funds, ranked by liquidity):} Cash in vault (2), Balances at central bank (3), Cheques in course of collection (1), Treasury bills (7), Investments in local firms (4), Loans and advances (41), Premises and equipment (6).
\end{itemize}

\textbf{Balance Sheet as at 31 December 2024 (millions of FCFA):}

\begin{center}
\begin{tabularx}{\linewidth}{|X|r|X|r|}
\hline
\rowcolor{popblueL} \textbf{Liabilities} & \textbf{Amount} & \textbf{Assets} & \textbf{Amount} \\
\hline
Share capital & 6 & Cash in vault & 2 \\
\hline
Reserves & 2 & Balances at central bank & 3 \\
\hline
Customer deposits & 50 & Cheques in course of collection & 1 \\
\hline
Borrowing from central bank & 5 & Treasury bills & 7 \\
\hline
Other liabilities & 1 & Investments in local firms & 4 \\
\hline
 &  & Loans and advances to customers & 41 \\
\hline
 &  & Premises and equipment & 6 \\
\hline
\textbf{Total Liabilities} & \textbf{64} & \textbf{Total Assets} & \textbf{64} \\
\hline
\end{tabularx}
\end{center}

The balance sheet balances at 64 million FCFA. Assets are arranged in order of decreasing liquidity (most liquid first); liabilities show the bank's obligations to depositors and creditors, plus shareholders' funds (capital and reserves).

\begin{attention}[Balance Sheet Convention]
In the GCE, always:
\begin{itemize}
\item Put liabilities on the left (or top) and assets on the right (or bottom).
\item List assets in order of liquidity: Cash $\to$ Central bank balances $\to$ Short-term securities (T-bills) $\to$ Investments $\to$ Loans $\to$ Fixed assets.
\item Ensure totals match exactly.
\item State the date and currency (millions of FCFA).
\end{itemize}
\end{attention}
\end{exoresolu}

\begin{exoresolu}[Task 7: Synthesis and Advice to MBATRACO]
\textbf{Statement:} In about ten lines, advise the cooperative on whether to keep savings in cash, deposit them, or borrow for the warehouse. Justify with figures.

\tcblower
\textbf{Model Advice (GCE Essay Style):}

Prices in Mbankomo rose by 22.7\% between 2022 and 2024 (RPI = 122.7), so the FCFA lost 18.5\% of its purchasing power. Keeping the cooperative's savings in cash guarantees a real loss of 18.5\% — a trader holding 260,000 FCFA in 2022 can now buy only 212,000 FCFA worth of goods at 2022 prices.

Depositing the 5,000,000 FCFA at Avenir Bank protects the nominal sum and earns interest, but if the deposit interest rate is below 22.7\% (the inflation rate), the real return remains negative. However, the banking system's credit multiplier of 10 means the cooperative's deposit supports up to 45,000,000 FCFA of new loans in the economy, making loanable funds abundant.

\textbf{Recommendation:} MBATRACO should \textbf{borrow now} to build the warehouse. The loan will be repaid in future FCFA of lower purchasing power (inflation benefits debtors), while the warehouse is a real asset whose value and storage revenues should rise with prices. The cooperative should keep only a minimal cash float for daily transactions, deposit the remainder to earn interest and support the loan application, and negotiate a fixed-rate loan whose real burden will shrink over time — provided the warehouse's expected return exceeds the interest cost.

\begin{aretenir}[What This Activity Consolidates]
\begin{itemize}
\item A \cle{Laspeyres price index} uses base-year quantities as weights: $\frac{\sum p_1 q_0}{\sum p_0 q_0} \times 100$.
\item \cle{Inflation} is the percentage rise in the index; the \cle{value of money} falls by $1 - \frac{100}{\text{Index}}$ — a smaller percentage.
\item The \cle{Quantity Theory} ($MV=PT$) predicts $\% \Delta P \approx \% \Delta M - \% \Delta T$; real-world deviations arise from supply shocks, velocity changes, and index limitations.
\item \cle{Credit creation} by the banking system: multiplier $k = 1/r$; total deposits $= k \times \text{primary deposit}$; total credit $= (k-1) \times \text{primary deposit}$.
\item A \cle{commercial bank balance sheet} must balance: Assets (liquid $\to$ illiquid) = Liabilities (deposits, borrowings, capital).
\item \cle{Financial decision-making} in inflation: avoid idle cash; use deposits for safety/interest; borrow for real assets if real interest rate $<$ expected return.
\end{itemize}
\end{aretenir}
\end{exoresolu}

\newpage

\coursec{Methods and worked examples}

\courssub{Skill 1: Identifying and Justifying the Functions of Money}

Before money can be studied as a macroeconomic variable (its demand, its supply, its value), the GCE candidate must be able to \emph{recognise} it in action. Examination questions frequently present a short scene --- a market transaction, a savings decision, a loan contract --- and ask which function of money is illustrated. The skill is to match each element of the scene to the correct function and to \emph{justify} the match with a quotation.

\begin{methode}[Analysing the functions of money in a given situation]
When a GCE question presents a transaction or a situation and asks you to identify which function of money is illustrated, proceed as follows:
\begin{enumerate}
\item \textbf{Recall the four functions precisely}: \cle{medium of exchange} (money used to settle transactions), \cle{unit of account / measure of value} (prices expressed in money), \cle{store of value} (purchasing power kept over time), \cle{standard of deferred payment} (debts and future payments expressed in money).
\item \textbf{Locate the key verb} in the statement: ``buys/sells'' $\rightarrow$ medium of exchange; ``is priced at / is valued at'' $\rightarrow$ unit of account; ``saves / keeps'' $\rightarrow$ store of value; ``borrows / repays later / pays in instalments'' $\rightarrow$ standard of deferred payment.
\item \textbf{Quote the relevant phrase} from the text to justify your answer: an identification without textual evidence earns only half the marks.
\item \textbf{Check the qualities of good money} if asked: general acceptability, durability, divisibility, portability, homogeneity (uniformity), scarcity, and stability of value.
\end{enumerate}
\end{methode}

\begin{exoresolu}[Functions of money in a market scene]
At Mokolo market in Yaound\'e, a trader displays a sign: ``Tomatoes: 500 FCFA per basket.'' Amina buys two baskets and pays 1\,000 FCFA in cash. She also keeps 20\,000 FCFA in a tin box at home for the school fees of her daughter, due in September. Finally, she has agreed to repay a 50\,000 FCFA loan from her \emph{tontine} in three monthly instalments.

Identify the function of money illustrated by each element of the scene.
\tcblower
\textbf{Solution.}

\textbf{Step 1: ``Tomatoes: 500 FCFA per basket.''} The price of tomatoes is \emph{expressed} in money; no exchange has yet taken place. Money is used as a \cle{unit of account} (measure of value): it allows the value of tomatoes to be stated and compared with that of other goods.

\textbf{Step 2: ``Amina buys two baskets and pays 1\,000 FCFA in cash.''} Money is physically handed over to settle the transaction immediately. Money functions as a \cle{medium of exchange}, eliminating the double coincidence of wants required under barter.

\textbf{Step 3: ``She keeps 20\,000 FCFA in a tin box \dots\ due in September.''} Purchasing power is held over time and will be spent later. Money acts as a \cle{store of value} (store of wealth). Note: this function is only well performed if the value of money is relatively stable, i.e.\ inflation is low.

\textbf{Step 4: ``Repay a 50\,000 FCFA loan \dots\ in three monthly instalments.''} A payment obligation is fixed in money terms but settled in the future. Money serves as a \cle{standard of deferred payment}, which makes borrowing, lending and hire-purchase contracts possible.

\textbf{Conclusion:} the four classical functions of money are all illustrated: unit of account, medium of exchange, store of value, and standard of deferred payment.
\end{exoresolu}

\courssub{Skill 2: Computing the Money Supply and Its Components}

The \cle{money supply} is the total stock of money held by the public in an economy at a given moment. Because modern money exists in several forms of differing liquidity, statisticians and central banks (such as the BEAC in the CEMAC zone) measure it through \emph{monetary aggregates}. The GCE candidate must know the two standard aggregates and, above all, what to \emph{include} and what to \emph{exclude}.

\begin{methode}[Calculating monetary aggregates (M1, M2)]
\begin{enumerate}
\item \textbf{Know the definitions}: $M1 =$ notes and coins in circulation with the public $+$ demand (sight) deposits. $M2 = M1 +$ quasi-money (savings and time deposits).
\item \textbf{Exclude} items held \emph{inside} banks (cash in bank vaults is not ``with the public'') and items that are not money at all (cheques and credit cards are payment \emph{instruments}, not money; treasury bills are financial assets).
\item \textbf{Add carefully}, showing each aggregate on a separate line.
\item \textbf{Interpret}: state which component dominates and what that reveals (degree of monetisation, banking habits of the population).
\end{enumerate}
\end{methode}

\begin{exoresolu}[Money supply in a hypothetical economy]
The following data (in billions of FCFA) relate to the economy of ``Kamerland'' in a given year: notes and coins held by the public: 800; cash in commercial bank vaults: 120; demand deposits at commercial banks: 1\,500; savings deposits: 900; time (fixed) deposits: 600; treasury bills held by the public: 300.

Calculate (a) $M1$; (b) quasi-money; (c) $M2$; (d) the percentage share of demand deposits in $M1$.
\tcblower
\textbf{Solution.}

\textbf{(a)} $M1 =$ notes and coins with the public $+$ demand deposits:
\[ M1 = 800 + 1\,500 = 2\,300 \text{ billion FCFA.} \]
Cash in bank vaults (120) is excluded: it is not ``with the public''.

\textbf{(b)} Quasi-money $=$ savings deposits $+$ time deposits:
\[ QM = 900 + 600 = 1\,500 \text{ billion FCFA.} \]

\textbf{(c)} $M2 = M1 + QM = 2\,300 + 1\,500 = 3\,800$ billion FCFA.

\textbf{(d)} Share of demand deposits in $M1$:
\[ \frac{1\,500}{2\,300}\times 100 \approx 65.2\%. \]

\textbf{Interpretation:} nearly two-thirds of narrow money consists of bank (scriptural) money, showing that the economy is well banked. Treasury bills (300) are excluded from $M2$ because they are financial assets, not money.
\end{exoresolu}

\courssub{Skill 3: Constructing and Interpreting a Retail Price Index}

To measure changes in the \emph{value of money} we first need a measure of the \emph{general price level}. The standard tool is the \cle{retail (consumer) price index}: the cost of a fixed ``basket'' of goods and services bought by a typical household, expressed as a percentage of its cost in a base year. This is the practical work most frequently set in the GCE on this chapter.

\begin{methode}[Calculating a Laspeyres consumer price index]
\begin{enumerate}
\item \textbf{Formula}: $I_t = \dfrac{\sum p_t q_0}{\sum p_0 q_0}\times 100$, where $p_0, q_0$ are base-year prices and quantities, $p_t$ current-year prices.
\item \textbf{Compute the cost of the fixed basket} in each year: multiply each price by the \emph{base-year} quantity, then sum.
\item \textbf{Divide and multiply by 100}. The base-year index must equal 100 --- use this as a check.
\item \textbf{Inflation rate} between two years: $\pi = \dfrac{I_t - I_{t-1}}{I_{t-1}}\times 100$.
\end{enumerate}
\textbf{Reflex:} if weights are given instead of quantities, use $I_t = \sum w_i \dfrac{p_{it}}{p_{i0}}$ with $\sum w_i = 100$ or 1.
\end{methode}

\begin{exoresolu}[Cost-of-living index for a household]
A household in Douala consumes three items with the following base-year (2020) quantities and prices, and 2023 prices:

\begin{center}
\begin{tabular}{|l|c|c|c|}
\hline
\rowcolor{popblueL} Item & $q_0$ (2020) & $p_0$ (FCFA) & $p_t$ (2023, FCFA) \\
\hline
Rice (kg) & 100 & 400 & 550 \\
\hline
Fish (kg) & 40 & 1\,500 & 2\,100 \\
\hline
Transport (trips) & 200 & 250 & 300 \\
\hline
\end{tabular}
\end{center}

(a) Calculate the Laspeyres price index for 2023 (base 2020 = 100). (b) Calculate the inflation rate between 2020 and 2023. (c) Comment.
\tcblower
\textbf{Solution.}

\textbf{(a)} Cost of the basket in 2020:
\[ \sum p_0 q_0 = (400\times 100) + (1\,500\times 40) + (250\times 200) = 40\,000 + 60\,000 + 50\,000 = 150\,000. \]
Cost of the same basket at 2023 prices:
\[ \sum p_t q_0 = (550\times 100) + (2\,100\times 40) + (300\times 200) = 55\,000 + 84\,000 + 60\,000 = 199\,000. \]
\[ I_{2023} = \frac{199\,000}{150\,000}\times 100 = 132.7 \quad (\text{to 1 d.p.}). \]

\textbf{(b)} Inflation rate over the period:
\[ \pi = \frac{132.7 - 100}{100}\times 100 = 32.7\%. \]

\textbf{(c)} The cost of the fixed basket rose by about 32.7\% between 2020 and 2023; the purchasing power of the FCFA for this household has therefore fallen. Fish contributed most to the rise (a 40\% price increase on a large weight). Note the limitation of the Laspeyres index: by fixing quantities at base-year levels, it ignores the fact that consumers substitute away from goods whose prices rise fastest, so it tends to \emph{overstate} the true rise in the cost of living.
\end{exoresolu}

\courssub{Skill 4: The Value of Money and Real versus Nominal Values}

The \cle{value of money} is what a unit of money can buy --- its purchasing power. It is measured by the inverse of the price level: when prices rise, the value of money falls. This inverse relationship is the key to distinguishing \emph{nominal} values (expressed in current FCFA) from \emph{real} values (expressed in constant purchasing power).

\begin{methode}[Converting nominal values into real values]
\begin{enumerate}
\item \textbf{Value of money} $= \dfrac{1}{P}$: it moves inversely to the price level. If the price index doubles, the value of money halves.
\item \textbf{Real value} $= \dfrac{\text{Nominal value}}{\text{Price index}}\times 100$.
\item \textbf{Purchasing power of 1 FCFA} in year $t$ (in base-year FCFA) $= \dfrac{100}{I_t}$.
\item Always state the \emph{direction}: rising prices $\Rightarrow$ falling value of money.
\end{enumerate}
\end{methode}

\begin{exoresolu}[Real income of a civil servant]
A civil servant in Bafoussam earned a nominal monthly salary of 150\,000 FCFA in 2020, when the consumer price index stood at 100. In 2023 his salary was raised to 180\,000 FCFA while the index reached 125.

(a) Calculate his real salary in 2023 at 2020 prices. (b) Has his standard of living improved? (c) Calculate the purchasing power of one 2023 FCFA in 2020 FCFA.
\tcblower
\textbf{Solution.}

\textbf{(a)} Real salary $= \dfrac{180\,000}{125}\times 100 = 144\,000$ FCFA (at 2020 prices).

\textbf{(b)} Although nominal salary rose by $\frac{180\,000-150\,000}{150\,000}\times 100 = 20\%$, prices rose by 25\%. Real salary fell from 150\,000 to 144\,000 FCFA, a loss of $\frac{6\,000}{150\,000}\times 100 = 4\%$. His standard of living has \emph{deteriorated} despite the pay rise: this is the effect of \cle{inflation} eroding money incomes.

\textbf{(c)} Purchasing power of 1 FCFA in 2023 $= \dfrac{100}{125} = 0.80$ FCFA of 2020. One FCFA in 2023 buys only what 0.80 FCFA bought in 2020; the value of money has fallen by 20\%.

\textbf{Exam tip:} never compare nominal incomes across years without deflating by the price index.
\end{exoresolu}

\courssub{Skill 5: Applying the Quantity Theory of Money}

Why does the price level --- and therefore the value of money --- change? The oldest and most famous answer is the \cle{quantity theory of money}, associated with the classical economists and formalised by Irving Fisher. It links the money supply directly to the price level and remains a favourite GCE topic, both for calculation and for essay evaluation.

\begin{methode}[Using Fisher's equation $MV = PT$ (or $MV = PY$)]
\begin{enumerate}
\item \textbf{Write the equation of exchange}: $MV = PT$, with $M$ = money supply, $V$ = velocity of circulation, $P$ = price level, $T$ = volume of transactions (or $Y$ = real output).
\item \textbf{Identify the unknown} and rearrange: $P = \dfrac{MV}{T}$, $M = \dfrac{PT}{V}$, $V = \dfrac{PT}{M}$.
\item \textbf{For proportional changes}, use the growth-rate form: $\%\Delta M + \%\Delta V = \%\Delta P + \%\Delta T$.
\item \textbf{Apply the classical assumptions} when asked for the theory's prediction: $V$ constant (institutional habits) and $T$ fixed at full-employment output $\Rightarrow$ $\%\Delta P = \%\Delta M$: money supply growth causes proportional inflation.
\item \textbf{Evaluate} (essay questions): Keynesians object that $V$ is unstable and $T$ can rise when there are idle resources, so money need not be neutral.
\end{enumerate}
\end{methode}

\begin{exoresolu}[Fisher's equation and inflation]
In the economy of ``Ndol\'eland'', the money supply is 500 billion FCFA, the velocity of circulation is 4, and the volume of transactions is 2\,000 billion units.

(a) Calculate the general price level. (b) The central bank increases the money supply by 20\% while $V$ and $T$ remain unchanged. Calculate the new price level and the rate of inflation. (c) State the two assumptions behind this result and one criticism of the theory.
\tcblower
\textbf{Solution.}

\textbf{(a)} From $MV = PT$:
\[ P = \frac{MV}{T} = \frac{500\times 4}{2\,000} = 1.0 \text{ FCFA per unit.} \]

\textbf{(b)} New money supply: $M' = 500\times 1.20 = 600$ billion FCFA.
\[ P' = \frac{M'V}{T} = \frac{600\times 4}{2\,000} = 1.20. \]
Inflation rate: $\pi = \dfrac{1.20 - 1.00}{1.00}\times 100 = 20\%$.

\textbf{(c)} Assumptions: (i) velocity $V$ is constant, determined by stable payment habits and institutions; (ii) transactions/output $T$ is fixed at the full-employment level in the long run. Under these assumptions the price level is \emph{directly proportional} to the money supply --- money is \cle{neutral}.

Criticism: Keynes argued that in a depressed economy with unemployed resources, extra money may raise output ($T$) rather than prices, and that $V$ varies with interest rates and confidence (liquidity preference). Hence the strict proportionality between $M$ and $P$ holds only under restrictive conditions.
\end{exoresolu}

\courssub{Skill 6: Credit Creation by Commercial Banks}

Commercial banks do not merely safe-keep deposits: by lending, they \emph{create} money. A single bank can lend only its excess reserves, but because loans are redeposited elsewhere in the system, the banking system as a whole multiplies an initial cash injection into a much larger volume of deposits. Mastering the \cle{credit (deposit) multiplier} is essential for the GCE.

\begin{methode}[Calculating the credit (deposit) multiplier]
\begin{enumerate}
\item \textbf{Key formula}: credit multiplier $k = \dfrac{1}{r}$, where $r$ is the cash (reserve) ratio expressed as a decimal.
\item \textbf{Total deposits created} $= \dfrac{\text{initial new deposit}}{r}$; \textbf{total credit created} $=$ total deposits $-$ initial deposit.
\item \textbf{Build the chain} if asked: each bank keeps $r$ of the deposit as cash and lends out $(1-r)$; the loan is re-deposited in the next bank, and so on.
\item \textbf{State the assumptions}: no cash leakage to the public, banks lend out all excess reserves, a single initial injection of new cash, borrowers redeposit in the banking system.
\item \textbf{Check the direction}: a \emph{lower} reserve ratio $\Rightarrow$ a \emph{larger} multiplier.
\end{enumerate}
\end{methode}

\begin{exoresolu}[Multiple credit creation]
A new deposit of 1\,000\,000 FCFA in cash is made at Bank A. All commercial banks keep a cash reserve ratio of 20\% and lend out the remainder; all loans are redeposited within the banking system.

(a) Show the first three rounds of the credit-creation process. (b) Calculate the total deposits and total credit created in the system. (c) What would total credit be if the reserve ratio fell to 10\%?
\tcblower
\textbf{Solution.}

\textbf{(a)} Round by round (FCFA):
\begin{center}
\begin{tabular}{|l|c|c|c|}
\hline
\rowcolor{popblueL} Bank & Deposit received & Cash kept (20\%) & Loan granted (80\%) \\
\hline
A & 1\,000\,000 & 200\,000 & 800\,000 \\
\hline
B & 800\,000 & 160\,000 & 640\,000 \\
\hline
C & 640\,000 & 128\,000 & 512\,000 \\
\hline
\end{tabular}
\end{center}

\textbf{(b)} Multiplier: $k = \dfrac{1}{0.20} = 5$.
\[ \text{Total deposits} = 5 \times 1\,000\,000 = 5\,000\,000 \text{ FCFA.} \]
\[ \text{Total credit created} = 5\,000\,000 - 1\,000\,000 = 4\,000\,000 \text{ FCFA.} \]

\textbf{(c)} With $r = 10\%$: $k = \dfrac{1}{0.10} = 10$; total deposits $= 10\,000\,000$ FCFA; total credit created $= 9\,000\,000$ FCFA.

\textbf{Comment:} a single bank can only lend its excess reserves (800\,000 FCFA for Bank A), but the banking \emph{system} multiplies deposits. In practice leakages (cash held by the public, banks' prudence, BEAC/COBAC regulations in the CEMAC zone) reduce the multiplier below its theoretical value.
\end{exoresolu}

\courssub{Skill 7: Constructing and Interpreting a Commercial Bank Balance Sheet}

A bank's \cle{balance sheet} is a snapshot, at a given date, of what the bank \emph{owes} (liabilities) and what it \emph{owns or is owed} (assets). Reading and constructing it is a classic GCE practical exercise, and it reveals the fundamental tension of banking: the most liquid assets earn the least, the most profitable assets are the least liquid.

\begin{methode}[Drawing up a bank's balance sheet]
\begin{enumerate}
\item \textbf{Remember the golden rule}: Total Assets $=$ Total Liabilities ($+$ capital). If they do not balance, an item is missing or misplaced.
\item \textbf{Classify each item}: \emph{Liabilities} = what the bank owes (share capital, reserves, demand deposits, savings/time deposits, borrowings). \emph{Assets} = what the bank owns or is owed, ranked from most liquid to least liquid: cash in hand and at the central bank, money at call, bills discounted, investments (securities), loans and advances, fixed assets (premises).
\item \textbf{Order assets by liquidity} (liquidity $\downarrow$, profitability $\uparrow$): cash earns nothing but is perfectly liquid; loans earn most but are least liquid. This is the \cle{liquidity--profitability conflict}.
\item \textbf{Verify} the balance: $\sum \text{Assets} = \sum \text{Liabilities}$.
\item \textbf{Interpret} ratios if asked: cash ratio $= \dfrac{\text{cash}}{\text{deposits}}$; liquidity ratio, etc.
\end{enumerate}
\end{methode}

\begin{exoresolu}[Balance sheet of the ``Wouri Commercial Bank'']
The following items (in millions of FCFA) relate to Wouri Commercial Bank at 31 December: share capital 5\,000; premises 3\,000; cash in hand and at the central bank 4\,000; demand deposits 18\,000; loans and advances to customers 14\,000; reserves 2\,000; investments in government securities 5\,000; savings and time deposits 7\,000; money at call and short notice 2\,000; bills discounted 2\,000; equipment 2\,000.

(a) Draw up the balance sheet. (b) Calculate the cash ratio and comment.
\tcblower
\textbf{Solution.}

\textbf{(a)} Classify and balance:

\begin{center}
\begin{tabular}{|l|r|}
\hline
\rowcolor{popblueL} \textbf{Liabilities} & \textbf{Millions FCFA} \\
\hline
Share capital & 5\,000 \\
\hline
Reserves & 2\,000 \\
\hline
Demand deposits & 18\,000 \\
\hline
Savings and time deposits & 7\,000 \\
\hline
\textbf{Total liabilities} & \textbf{32\,000} \\
\hline
\end{tabular}
\end{center}

\begin{center}
\begin{tabular}{|l|r|}
\hline
\rowcolor{popblueL} \textbf{Assets (most $\rightarrow$ least liquid)} & \textbf{Millions FCFA} \\
\hline
Cash in hand and at central bank & 4\,000 \\
\hline
Money at call and short notice & 2\,000 \\
\hline
Bills discounted & 2\,000 \\
\hline
Investments (government securities) & 5\,000 \\
\hline
Loans and advances & 14\,000 \\
\hline
Premises and equipment & 5\,000 \\
\hline
\textbf{Total assets} & \textbf{32\,000} \\
\hline
\end{tabular}
\end{center}

(Premises 3\,000 $+$ equipment 2\,000 $= 5\,000$.) Check: assets $=$ liabilities $= 32\,000$ million FCFA. \checkmark

\textbf{(b)} Cash ratio:
\[ \frac{\text{Cash}}{\text{Total deposits}} = \frac{4\,000}{18\,000 + 7\,000} = \frac{4\,000}{25\,000} = 16\%. \]
The bank keeps 16\% of its deposit liabilities in cash, enough to meet normal withdrawals while lending the rest. Loans and advances (14\,000 million, i.e.\ $\frac{14\,000}{32\,000}\times 100 = 43.75\%$ of assets) are the most profitable but least liquid asset --- illustrating the trade-off between \cle{liquidity}, \cle{profitability} and \cle{security} that governs portfolio management in commercial banking.
\end{exoresolu}

\begin{attention}[Common mistakes in this chapter]
\begin{itemize}
\item Confusing the \emph{value of money} with the \emph{price level}: they move in opposite directions.
\item Deflating nominal values by multiplying by the index instead of dividing.
\item Forgetting to subtract the initial deposit when asked for \emph{credit created} (as opposed to total deposits).
\item Placing deposits among assets in a balance sheet: deposits are \emph{liabilities} of the bank (it owes them to customers); loans are \emph{assets}.
\item Applying $MV=PT$ mechanically without stating the assumptions of constant $V$ and full-employment $T$.
\end{itemize}
\end{attention}

\newpage

\coursec{Exercises}

\courssub{I check my knowledge}

\begin{exercice}[Multiple choice questions]
For each question, choose the single correct answer and write it down.
\begin{enumerate}
\item Money is best defined as:
\begin{enumerate}
\item any commodity that is rare and durable;
\item anything that is generally acceptable as a means of payment for goods, services and the settlement of debts;
\item notes and coins issued by the central bank only;
\item gold and silver held by commercial banks.
\end{enumerate}
\item Which of the following is NOT a function of money?
\begin{enumerate}
\item Medium of exchange;
\item Store of value;
\item Standard of deferred payment;
\item Source of profit for commercial banks.
\end{enumerate}
\item In the equation of exchange $MV = PT$, the letter $V$ stands for:
\begin{enumerate}
\item the value of money;
\item the velocity of circulation of money;
\item the volume of bank deposits;
\item the value of transactions.
\end{enumerate}
\item If a commercial bank keeps a cash ratio of 10\%, the credit multiplier is equal to:
\begin{enumerate}
\item 0.1; \item 1; \item 10; \item 100.
\end{enumerate}
\end{enumerate}
\end{exercice}

\begin{exercice}[True or false — justify]
State whether each of the following statements is \textbf{true} or \textbf{false}, and justify your answer in one or two sentences.
\begin{enumerate}
\item Money eliminates the need for a double coincidence of wants that characterises barter.
\item Demand deposits (current accounts) are excluded from the narrow definition of the money supply, M1.
\item According to the quantity theory of money, an increase in the money supply raises the price level proportionally, whatever the state of the economy.
\item In a commercial bank's balance sheet, deposits of customers appear on the assets side because they are owned by the customers.
\end{enumerate}
\end{exercice}

\begin{exercice}[Definitions]
Define each of the following terms in one or two clear sentences, giving an example where appropriate:
\begin{enumerate}
\item legal tender;
\item near money (quasi-money);
\item the retail price index;
\item money at call and short notice;
\item liquidity (in relation to a bank's assets).
\end{enumerate}
\end{exercice}

\begin{exercice}[Direct calculations]
\begin{enumerate}
\item The retail price index of a country stood at 125 in 2022 and 140 in 2023. Calculate the rate of inflation between the two years.
\item Using the same figures, calculate the value of money (in index form) in 2023, taking 2022 as the base year. Comment on your result.
\item In an economy, the money supply is $M = 400$ billion FCFA, the velocity of circulation is $V = 5$ and the volume of transactions is $T = 2\,000$ billion units. Using the Fisher equation, calculate the general price level $P$.
\end{enumerate}
\end{exercice}

\courssub{I practise}

\begin{exercice}[Constructing a price index]
The table below shows the prices and quantities of the three goods consumed by a typical household in Bafoussam in 2023 (base year) and 2024.
\begin{center}
\begin{tabular}{|l|c|c|c|c|}
\hline
\rowcolor{popblueL} Good & $P_{2023}$ (FCFA) & $Q_{2023}$ & $P_{2024}$ (FCFA) & $Q_{2024}$ \\
\hline
Rice (kg) & 400 & 50 & 500 & 55 \\
\hline
Bread (loaf) & 150 & 200 & 180 & 210 \\
\hline
Palm oil (litre) & 800 & 30 & 1\,000 & 28 \\
\hline
\end{tabular}
\end{center}
\begin{enumerate}
\item Construct a Laspeyres price index for 2024, taking 2023 as the base year ($= 100$).
\item Calculate the rate of inflation between 2023 and 2024.
\item Calculate the value of money in 2024 and interpret your result.
\item State two reasons why this index may not perfectly measure the cost of living of the household.
\end{enumerate}
\end{exercice}

\begin{exercice}[The quantity theory of money]
In a hypothetical CEMAC economy, $M = 500$ billion FCFA, $V = 4$ and $T = 2\,000$ billion units of output.
\begin{enumerate}
\item Calculate the equilibrium general price level $P$.
\item Suppose the money supply increases by 20\%, while $V$ and $T$ remain unchanged. Calculate the new price level and the rate of inflation.
\item State the assumptions of the quantity theory of money that you used in (b), and explain briefly why each may not hold in practice.
\item If the economy were operating with substantial unemployed resources, what would you expect to happen instead? Explain.
\end{enumerate}
\end{exercice}

\begin{exercice}[Credit creation]
A commercial bank in Douala receives a new cash deposit of $2\,000\,000$ FCFA. All banks in the system keep a cash (reserve) ratio of 20\%, and there is no cash leakage to the public.
\begin{enumerate}
\item Calculate the credit (deposit) multiplier.
\item Calculate the total amount of deposits that the banking system can create from this initial deposit.
\item Calculate the total amount of loans (advances) created in the process.
\item State three limitations that may prevent the banking system from expanding credit to this theoretical maximum in Cameroon.
\end{enumerate}
\end{exercice}

\begin{exercice}[Building a bank's balance sheet]
The accounts of Fidelity Trust Bank Ltd as at 31 December 2024 contain the following items (in millions of FCFA), given in no particular order: cash in till 500; customers' deposits 8\,000; loans and advances 4\,500; premises 900; bills discounted 600; reserves at the central bank 1\,200; investments 1\,800; money at call and short notice 500; share capital 1\,500; undistributed profit 500.
\begin{enumerate}
\item Construct a properly ordered balance sheet, placing assets in reverse order of liquidity (from the least liquid to the most liquid) and liabilities appropriately.
\item Verify that the balance sheet balances.
\item Identify which items are the bank's most liquid assets and explain why a bank must hold them.
\item Explain why deposits are a liability to the bank while loans are an asset.
\end{enumerate}
\end{exercice}

\courssub{I prepare for the GCE}

\begin{exercice}[GCE structured question: money and its functions]
\textbf{(a)} Define money and briefly trace its evolution from barter to modern forms of money. \hfill (6 marks)

\textbf{(b)} Explain the four functions of money, illustrating each with an example from the Cameroonian economy. \hfill (8 marks)

\textbf{(c)} Mobile money services (such as MTN Mobile Money and Orange Money) are widely used in Cameroon. Discuss the extent to which mobile money performs each of the functions of money, and identify two problems that may limit its use as money. \hfill (6 marks)

\hfill \emph{(Total: 20 marks)}
\end{exercice}

\begin{exercice}[GCE essay: money supply and inflation]
\begin{quote}
\emph{``An increase in the money supply is the main cause of inflation in Cameroon.''}
\end{quote}
Discuss this statement with reference to the quantity theory of money. Your answer should:
\begin{enumerate}
\item state and explain the Fisher equation of exchange and its assumptions;
\item show, with the help of a numerical example, how an increase in the money supply can lead to inflation;
\item examine other possible causes of inflation in Cameroon (cost-push factors, imported inflation, structural factors);
\item reach a justified conclusion on the validity of the statement.
\end{enumerate}
\hfill \emph{(Total: 20 marks)}
\end{exercice}

\begin{exercice}[GCE structured question: commercial banks]
\textbf{(a)} Distinguish between M1 and M2 as measures of the money supply. \hfill (4 marks)

\textbf{(b)} Explain the three Keynesian motives for holding money (demand for money). \hfill (6 marks)

\textbf{(c)} With the aid of a clearly labelled diagram, show the effect of an increase in the money supply on the equilibrium rate of interest. \hfill (4 marks)

\textbf{(d)} Explain the functions of commercial banks and evaluate their contribution to the economic development of Cameroon. \hfill (6 marks)

\hfill \emph{(Total: 20 marks)}
\end{exercice}

\newpage

\coursec{Solutions to the exercises}

\begin{corrige}[Exercise 1 — Multiple choice questions]
\begin{enumerate}
\item \textbf{Answer: (b).} Money is defined by its \emph{general acceptability} as a means of payment and settlement of debts, not by its physical form. Rarity and durability (a) are desirable qualities of commodity money but do not define money; (c) is too narrow because bank deposits are also money; (d) confuses money with bank reserves.
\item \textbf{Answer: (d).} The four functions of money are: \cle{medium of exchange}, \cle{unit of account} (measure of value), \cle{store of value} and \cle{standard of deferred payment}. Being a source of profit for banks is not a function of money.
\item \textbf{Answer: (b).} In Fisher's equation of exchange $MV = PT$, $V$ is the \cle{velocity of circulation}, i.e.\ the average number of times a unit of money changes hands in a given period.
\item \textbf{Answer: (c).} The \cle{credit multiplier} is the reciprocal of the cash ratio: $k = \dfrac{1}{r} = \dfrac{1}{0.10} = 10$.
\end{enumerate}
\end{corrige}

\begin{corrige}[Exercise 2 — True or false — justify]
\begin{enumerate}
\item \textbf{True.} Barter requires each trader to find someone who wants what he offers and offers what he wants (\cle{double coincidence of wants}). As a generally accepted \cle{medium of exchange}, money allows a person to sell for money and buy with money separately, eliminating this requirement.
\item \textbf{False.} Demand deposits are \emph{included} in \cle{M1}. The narrow money supply $M1$ = notes and coins in circulation \emph{plus} demand (sight/current account) deposits, because both are immediately usable as means of payment.
\item \textbf{False.} The proportional link between $M$ and $P$ holds only under the \cle{quantity theory of money} assumptions: $V$ and $T$ constant, which requires full employment of resources. If there are unemployed resources, an increase in $M$ may raise output ($T$) rather than prices.
\item \textbf{False.} Customers' deposits appear on the \emph{liabilities} side precisely because they belong to customers: the bank owes this money to depositors and must repay it on demand or at maturity. What the bank owns (loans, investments, premises) appears on the assets side.
\end{enumerate}
\end{corrige}

\begin{corrige}[Exercise 3 — Definitions]
\begin{enumerate}
\item \textbf{\cle{Legal tender}:} money which, by law, must be accepted in settlement of a debt within a country, and which a creditor cannot legally refuse. In Cameroon, FCFA notes and coins issued by the BEAC are legal tender (coins only up to legally fixed amounts).
\item \textbf{\cle{Near money} (quasi-money):} assets that are not directly usable as a means of payment but can be converted into money quickly and with little or no loss of value, e.g.\ time (savings) deposits, treasury bills.
\item \textbf{\cle{Retail price index (RPI)}:} a weighted average measure of the prices of a representative ``basket'' of goods and services bought by a typical household, expressed as a percentage of prices in a base year; it is used to measure \cle{inflation} and changes in the cost of living.
\item \textbf{\cle{Money at call and short notice}:} short-term loans made by a commercial bank (often to the money market or discount houses) which must be repaid immediately on demand (``at call'') or within a very short period (``short notice''); it is a highly liquid asset.
\item \textbf{\cle{Liquidity} (of a bank's assets):} the ease and speed with which an asset can be converted into cash without loss of value. Cash is perfectly liquid; premises are highly illiquid.
\end{enumerate}
\end{corrige}

\begin{corrige}[Exercise 4 — Direct calculations]
\begin{enumerate}
\item Rate of \cle{inflation}:
\[ \pi = \frac{\text{RPI}_{2023} - \text{RPI}_{2022}}{\text{RPI}_{2022}} \times 100 = \frac{140 - 125}{125} \times 100 = \frac{15}{125} \times 100 = 12\%. \]
Prices rose by $12\%$ between 2022 and 2023.
\item The \cle{value of money} is the reciprocal of the price index (in index form, base year $= 100$):
\[ V_m = \frac{\text{RPI}_{2022}}{\text{RPI}_{2023}} \times 100 = \frac{125}{140} \times 100 \approx 89.3. \]
\textbf{Comment:} the value (purchasing power) of money has fallen by about $10.7\%$: one FCFA in 2023 buys only about $89\%$ of what it bought in 2022. Note that the fall in the value of money ($10.7\%$) is not numerically equal to the rise in prices ($12\%$), because the two are computed on different bases.
\item From $MV = PT$:
\[ P = \frac{MV}{T} = \frac{400 \times 5}{2\,000} = \frac{2\,000}{2\,000} = 1. \]
The general price level is $P = 1$ (FCFA per unit of transaction).
\end{enumerate}
\end{corrige}

\begin{corrige}[Exercise 5 — Constructing a price index]
\begin{enumerate}
\item The \cle{Laspeyres index} uses \emph{base-year quantities} as weights:
\[ L = \frac{\sum P_{2024} Q_{2023}}{\sum P_{2023} Q_{2023}} \times 100. \]
Cost of the base-year basket at 2023 prices:
\[ (400 \times 50) + (150 \times 200) + (800 \times 30) = 20\,000 + 30\,000 + 24\,000 = 74\,000 \text{ FCFA}. \]
Cost of the same basket at 2024 prices:
\[ (500 \times 50) + (180 \times 200) + (1\,000 \times 30) = 25\,000 + 36\,000 + 30\,000 = 91\,000 \text{ FCFA}. \]
\[ L = \frac{91\,000}{74\,000} \times 100 \approx 122.97 \approx 123.0. \]
\item Rate of inflation:
\[ \pi = \frac{123.0 - 100}{100} \times 100 \approx 23.0\%. \]
\item Value of money in 2024:
\[ V_m = \frac{100}{123.0} \times 100 \approx 81.3. \]
\textbf{Interpretation:} the purchasing power of the FCFA has fallen by about $18.7\%$; 100 FCFA in 2024 buys only what about 81 FCFA bought in 2023.
\item Two limitations (any two):
\begin{itemize}
\item The basket of only three goods is not representative of all household spending (housing, transport, school fees, health are omitted).
\item Quantities consumed change over time (here $Q_{2024} \neq Q_{2023}$): Laspeyres keeps base-year weights fixed and ignores substitution towards relatively cheaper goods, so it tends to overstate the rise in the cost of living (\cle{substitution bias}).
\item It ignores changes in the quality of goods.
\end{itemize}
\end{enumerate}
\end{corrige}

\begin{corrige}[Exercise 6 — The quantity theory of money]
\begin{enumerate}
\item Equilibrium price level:
\[ P = \frac{MV}{T} = \frac{500 \times 4}{2\,000} = \frac{2\,000}{2\,000} = 1. \]
\item New money supply: $M' = 500 \times 1.20 = 600$ billion FCFA. With $V$ and $T$ unchanged:
\[ P' = \frac{M'V}{T} = \frac{600 \times 4}{2\,000} = 1.2. \]
Rate of inflation:
\[ \pi = \frac{1.2 - 1}{1} \times 100 = 20\%. \]
The price level rises proportionally (by the same $20\%$) with the money supply, as the quantity theory predicts.
\item Assumptions used and their practical weaknesses:
\begin{itemize}
\item \textbf{$V$ is constant:} in practice velocity varies with payment habits, financial innovation (mobile money) and confidence; in a crisis people may hold more cash, reducing $V$.
\item \textbf{$T$ is constant (full employment):} output is only fixed at full employment; with idle resources, extra money can raise real output instead of prices.
\item \textbf{Causality runs from $M$ to $P$:} in reality prices can rise first (e.g.\ imported inflation) and the money supply may then accommodate the higher prices.
\item \textbf{Independence of the variables:} a change in $M$ may itself affect $V$ or $T$.
\end{itemize}
\item With substantial unemployed resources, the increase in the money supply would raise aggregate demand, encouraging firms to employ idle labour and capital. Output ($T$) would rise, so $P$ would rise \emph{less than proportionally} — or hardly at all. This is the Keynesian view: money is not neutral below full employment.
\end{enumerate}
\end{corrige}

\begin{corrige}[Exercise 7 — Credit creation]
\begin{enumerate}
\item \cle{Credit (deposit) multiplier}:
\[ k = \frac{1}{r} = \frac{1}{0.20} = 5. \]
\item Total deposits created:
\[ D = k \times \text{initial deposit} = 5 \times 2\,000\,000 = 10\,000\,000 \text{ FCFA}. \]
\item Total loans (advances) created = total new deposits minus the initial cash deposit:
\[ L = 10\,000\,000 - 2\,000\,000 = 8\,000\,000 \text{ FCFA}. \]
(Equivalently, each bank lends $80\%$ of what it receives: $1\,600\,000 + 1\,280\,000 + \dots = 8\,000\,000$.)
\item Three limitations in Cameroon (any three):
\begin{itemize}
\item \textbf{Cash leakages:} the public may hold part of loans as cash instead of redepositing it, draining reserves from the banking system.
\item \textbf{Insufficient demand for loans:} creditworthy borrowers may be lacking, especially in a largely informal economy where many agents have no collateral or bankable projects.
\item \textbf{Banks' prudence / excess reserves:} banks may voluntarily keep reserves above the required $20\%$ (liquidity preference, risk aversion).
\item \textbf{Central bank regulation:} the BEAC may impose reserve requirements, credit ceilings or refinancing conditions that restrict expansion.
\item \textbf{Lack of bank branches / financial exclusion:} much of the population is unbanked, so redeposits do not fully return to the system.
\end{itemize}
\end{enumerate}
\end{corrige}

\begin{corrige}[Exercise 8 — Building a bank's balance sheet]
\begin{enumerate}
\item Ordered balance sheet (millions of FCFA). Assets are ranked from the \emph{least} liquid (premises) to the \emph{most} liquid (cash):
\begin{center}
\begin{tabular}{|l|r|l|r|}
\hline
\rowcolor{popblueL} \textbf{Liabilities} & \textbf{Amount} & \textbf{Assets} & \textbf{Amount} \\
\hline
Share capital & 1\,500 & Premises & 900 \\
\hline
Undistributed profit & 500 & Investments & 1\,800 \\
\hline
Customers' deposits & 8\,000 & Loans and advances & 4\,500 \\
\hline
 & & Bills discounted & 600 \\
\hline
 & & Money at call and short notice & 500 \\
\hline
 & & Reserves at the central bank & 1\,200 \\
\hline
 & & Cash in till & 500 \\
\hline
\textbf{Total} & \textbf{10\,000} & \textbf{Total} & \textbf{10\,000} \\
\hline
\end{tabular}
\end{center}
\item Verification: Liabilities $= 1\,500 + 500 + 8\,000 = 10\,000$; Assets $= 900 + 1\,800 + 4\,500 + 600 + 500 + 1\,200 + 500 = 10\,000$. The balance sheet balances: $10\,000 = 10\,000$ million FCFA. $\blacksquare$
\item The most liquid assets are \textbf{cash in till} and \textbf{reserves at the central bank} (followed by money at call and short notice). A bank must hold them to meet customers' day-to-day withdrawals on demand, to comply with the reserve requirements of the central bank (BEAC), and to maintain public confidence. Because deposits are repayable on demand while most assets are tied up in illiquid loans, insufficient liquidity could cause a bank run and failure.
\item \textbf{Deposits are a liability} because the money belongs to the customers: the bank owes it to them and must repay it on demand (current accounts) or at an agreed date (time deposits). \textbf{Loans are an asset} because they represent claims of the bank on borrowers: borrowers owe the money to the bank, and the loan generates interest income. In short, deposits are what the bank \emph{owes}; loans are what is \emph{owed to} the bank.
\end{enumerate}
\end{corrige}

\begin{corrige}[Exercise 9 — GCE structured question: money and its functions]
\textbf{Indicative mark scheme and model content (Total: 20 marks).}

\textbf{(a) Definition and evolution (6 marks).}
\begin{itemize}
\item Definition (2 marks): money is anything that is \emph{generally acceptable} as a medium of exchange and in settlement of debts.
\item Evolution (4 marks, 1 mark per stage, any four): barter (direct exchange of goods, requiring double coincidence of wants) $\rightarrow$ commodity money (cowrie shells, salt, cattle — used in pre-colonial Cameroon) $\rightarrow$ metallic money (gold and silver coins) $\rightarrow$ paper money (convertible, then fiduciary notes) $\rightarrow$ bank money (chequeable deposits) $\rightarrow$ modern forms: plastic cards and electronic/mobile money (MTN Mobile Money, Orange Money).
\end{itemize}

\textbf{(b) Four functions (8 marks = 2 marks each: 1 for explanation, 1 for a Cameroonian illustration).}
\begin{enumerate}
\item \textbf{Medium of exchange:} money is accepted in payment for goods and services, removing barter's double coincidence of wants. \emph{Example:} a trader in Mokolo market sells plantains for FCFA and uses it to buy clothes.
\item \textbf{Unit of account (measure of value):} prices and accounts are expressed in money, allowing comparison of values. \emph{Example:} cocoa is quoted in FCFA per kilogram, so farmers can compare its value with fertiliser costs.
\item \textbf{Store of value:} money keeps purchasing power over time (provided inflation is low), enabling saving. \emph{Example:} a civil servant keeps part of her salary in a savings account at a microfinance institution.
\item \textbf{Standard of deferred payment:} money allows borrowing and lending with repayment fixed in money terms. \emph{Example:} a trader borrows 500\,000 FCFA from a credit union (tontine/njangui or MC$^2$) to be repaid in monthly instalments.
\end{enumerate}

\textbf{(c) Mobile money (6 marks).}
\begin{itemize}
\item \emph{Functions performed (up to 4 marks):} medium of exchange — transfers, airtime purchase, bill payment, payment in many shops; unit of account — balances are denominated in FCFA; store of value — e-wallets hold value, though without interest; standard of deferred payment — limited, though nano-credit services are emerging.
\item \emph{Two problems (2 marks, any two):} network failures and dependence on electricity/telecom coverage; fraud and SIM-swap scams; agent illiquidity (agents running out of cash); transaction fees; exclusion of those without phones or IDs; balances not universally accepted for large transactions.
\item \emph{Examiner expectation:} a balanced judgement — mobile money performs the medium-of-exchange function remarkably well in Cameroon but is still a complement to, rather than a full substitute for, cash and bank money.
\end{itemize}
\end{corrige}

\begin{corrige}[Exercise 10 — GCE essay: money supply and inflation]
\textbf{Indicative mark scheme (Total: 20 marks).}

\textbf{Introduction (2 marks).} Define inflation (a sustained rise in the general price level) and the money supply; announce the plan: the monetarist view via the quantity theory, then alternative explanations, then a judgement for Cameroon.

\textbf{1. The Fisher equation and its assumptions (6 marks).}
\begin{itemize}
\item Statement: $MV = PT$, where $M$ = money supply, $V$ = velocity of circulation, $P$ = general price level, $T$ = volume of transactions (output). The equation is an identity: total spending equals total receipts.
\item The quantity theory turns it into a theory of prices by assuming $V$ constant (stable payment habits) and $T$ fixed at the full-employment level; then $P = \dfrac{MV}{T}$ and any change in $M$ causes a \emph{proportional} change in $P$.
\item Assumptions clearly listed: constant $V$; constant/full-employment $T$; causality from $M$ to $P$; independence of variables.
\end{itemize}

\textbf{2. Numerical illustration (4 marks).} Example: $M = 500$, $V = 4$, $T = 2\,000 \Rightarrow P = \dfrac{500 \times 4}{2\,000} = 1$. If $M$ rises by $20\%$ to $600$ with $V$, $T$ unchanged: $P' = \dfrac{600 \times 4}{2\,000} = 1.2$, i.e.\ $20\%$ inflation. Money growth translates fully into price growth.

\textbf{3. Other causes of inflation in Cameroon (6 marks, any three developed).}
\begin{itemize}
\item \textbf{Cost-push inflation:} rising costs of fuel, electricity, imported inputs and wages push supply curves leftward; e.g.\ increases in pump prices raise transport and food costs.
\item \textbf{Imported inflation:} Cameroon imports fuel, wheat, rice, machinery; a rise in world prices or a rise in import costs transmits directly into domestic prices.
\item \textbf{Structural factors:} poor roads and storage (high distribution costs), supply bottlenecks in agriculture, insecurity in some regions disrupting food supply, market imperfections and hoarding.
\item \textbf{Demand-pull from non-monetary sources:} government deficit spending, export booms.
\item \textbf{Expectations:} once expected, inflation feeds wage and price demands.
\end{itemize}

\textbf{4. Evaluation and conclusion (2 marks).} The statement is only \emph{partly} valid. In the long run and near full employment, sustained money growth is indeed associated with inflation, and in the CEMAC zone monetary expansion is constrained by the BEAC and the CFA franc's peg to the euro, which has historically limited monetary financing. But in Cameroon, with unemployed resources and large import dependence, much observed inflation is cost-push, imported and structural rather than purely monetary. Hence money supply growth is \emph{a} cause, and a necessary permissive condition for sustained inflation, but not always the \emph{main} cause.

\emph{Examiner expectations:} correct statement and use of $MV = PT$; a worked numerical example; at least two non-monetary causes applied to Cameroon; a reasoned, two-sided conclusion. Award full marks only where the answer evaluates rather than merely describes.
\end{corrige}

\begin{corrige}[Exercise 11 — GCE structured question: commercial banks]
\textbf{Indicative mark scheme and model content (Total: 20 marks).}

\textbf{(a) M1 vs M2 (4 marks).}
\begin{itemize}
\item \cle{M1} (narrow money) = notes and coins in circulation \emph{plus} demand (sight/current account) deposits — assets immediately usable for payments. (2 marks)
\item \cle{M2} (broad money) = $M1$ \emph{plus} near money: time and savings deposits and similar quasi-liquid assets that can be converted into means of payment with some delay. (2 marks)
\end{itemize}

\textbf{(b) Keynesian motives for holding money (6 marks = 2 each).}
\begin{enumerate}
\item \textbf{\cle{Transactions motive}:} money is held to finance everyday purchases because receipts and payments are not synchronised; demand rises with income.
\item \textbf{\cle{Precautionary motive}:} money is held as a reserve against unforeseen expenses (illness, repairs); it also depends mainly on income.
\item \textbf{\cle{Speculative motive}:} money is held to take advantage of future changes in bond/asset prices; it varies \emph{inversely} with the rate of interest — at high interest rates people prefer bonds, at low rates they prefer liquidity (\cle{liquidity trap} at very low rates).
\end{enumerate}

\textbf{(c) Diagram (4 marks).} Money market equilibrium: the \cle{liquidity preference} curve $L$ (or $M_d$) slopes downward; money supply $M_s$ is vertical (fixed by the central bank). An increase in the money supply shifts $M_s$ rightwards (from $M_{s1}$ to $M_{s2}$); at the old interest rate there is excess supply of money, agents buy bonds, bond prices rise and the equilibrium interest rate falls from $r_1$ to $r_2$.
\begin{popfigure}
\begin{tikzpicture}[scale=1.1]
\draw[->,thick] (0,0) -- (6.2,0) node[below left,align=center,text width=2.2cm]{Quantity of money};
\draw[->,thick] (0,0) -- (0,4.2) node[below left,align=center,text width=1.8cm]{Rate of interest $r$};
\draw[very thick,popblue,domain=0.9:5.4] plot (\x,{2.6/(\x-0.2)+0.15}) node[right]{$L$};
\draw[very thick,poporange] (2.2,0) -- (2.2,3.6) node[above]{$M_{s1}$};
\draw[very thick,popgreen] (3.6,0) -- (3.6,3.6) node[above]{$M_{s2}$};
\draw[dashed,gray] (0,1.45) node[left]{$r_1$} -- (2.2,1.45);
\draw[dashed,gray] (0,0.91) node[left]{$r_2$} -- (3.6,0.91);
\fill[popdark] (2.2,1.45) circle (1.5pt);
\fill[popdark] (3.6,0.91) circle (1.5pt);
\end{tikzpicture}
\legende{An increase in the money supply lowers the equilibrium rate of interest.}
\end{popfigure}
Marks: axes labelled (1); downward-sloping $L$ and vertical $M_s$ (1); rightward shift of $M_s$ (1); fall in $r$ shown and explained (1).

\textbf{(d) Functions of commercial banks and their contribution to Cameroon (6 marks).}
\begin{itemize}
\item \emph{Functions (3 marks, any three):} accepting deposits (current, savings, time accounts); granting loans and advances and discounting bills; creating credit through the deposit multiplier; providing means of payment (cheques, cards, transfers); agency services (standing orders, foreign exchange, safe custody); financial advice.
\item \emph{Contribution to development (2 marks):} mobilising scattered savings and channelling them into productive investment; financing trade (cocoa, imports/exports through the port of Douala); supporting SMEs and agriculture; facilitating payments and financial inclusion.
\item \emph{Evaluation (1 mark):} limitations — lending concentrated in urban areas and large firms, high collateral requirements excluding small farmers and informal businesses, high interest spreads, competition from mobile money and informal tontines; hence their contribution, though real, remains below potential.
\end{itemize}
\end{corrige}

\newpage

\coursec{Summary sheet}

\begin{popfigure}
\begin{tikzpicture}[mindmap, grow cyclic,
  every node/.style={concept, execute at begin node=\hspace{0pt}},
  root concept/.append style={concept color=popdark, font=\small\bfseries, text=white, minimum size=3.2cm},
  level 1 concept/.append style={level distance=4.4cm, sibling angle=72, font=\scriptsize, minimum size=2.6cm},
  level 2 concept/.append style={level distance=2.6cm, font=\tiny, minimum size=1.9cm}]
\node[root concept]{Money and Banking}
  child[concept color=popblue]{ node{Money: nature \& functions}
    child[concept color=popblue!60]{ node{Barter, evolution}}
    child[concept color=popblue!60]{ node{4 functions, qualities}}
  }
  child[concept color=poporange]{ node{Demand \& supply of money}
    child[concept color=poporange!60]{ node{Keynes: 3 motives}}
    child[concept color=poporange!60]{ node{$M_1, M_2, M_3$}}
  }
  child[concept color=popgreen]{ node{Value of money}
    child[concept color=popgreen!60]{ node{RPI, inflation}}
    child[concept color=popgreen!60]{ node{$MV=PT$, $MV=PY$}}
  }
  child[concept color=poppurple]{ node{Commercial banks}
    child[concept color=poppurple!60]{ node{Functions, types}}
    child[concept color=poppurple!60]{ node{Credit creation}}
  }
  child[concept color=popgold]{ node{Bank balance sheet}
    child[concept color=popgold!60]{ node{Assets vs liabilities}}
    child[concept color=popgold!60]{ node{Liquidity vs profit}}
  };
\end{tikzpicture}
\legende{Mind map of the chapter ``Money and Banking''.}
\end{popfigure}

\begin{bilan}
\textbf{1. Key definitions}\par
\begin{itemize}
\item \cle{Money}: anything generally acceptable as a medium of exchange and in settlement of debts (notes, coins, bank deposits).
\item \cle{Barter}: direct exchange of goods for goods; limited by the \emph{double coincidence of wants}.
\item \cle{Functions of money}: medium of exchange; unit of account (measure of value); store of value; standard of deferred payment.
\item \cle{Qualities of good money}: acceptability, durability, portability, divisibility, homogeneity, scarcity, recognisability.
\item \cle{Money supply}: total stock of money in an economy at a time. $M_1$ = notes and coins + demand (sight) deposits; $M_2 = M_1$ + savings and time deposits; $M_3 = M_2$ + deposits at non-bank financial institutions (broad money).
\item \cle{Demand for money} (Keynes): desire to hold wealth in liquid form, arising from three motives --- \emph{transactions} and \emph{precautionary} motives (depend on income) and the \emph{speculative} motive (depends inversely on the interest rate).
\item \cle{Value of money}: the purchasing power of a unit of money; varies inversely with the general price level.
\item \cle{Retail Price Index (RPI)}: weighted average index measuring changes in the prices of a basket of goods and services consumed by a typical household.
\item \cle{Inflation}: persistent rise in the general price level (fall in the value of money); \cle{deflation} is the opposite.
\item \cle{Commercial bank}: a financial institution that accepts deposits, grants loans and creates credit for profit (e.g. Afriland, BICEC, SGC in Cameroon).
\end{itemize}

\textbf{2. Laws, formulas and theories to know}\par
\begin{itemize}
\item \cle{Quantity Theory of Money} (Fisher): $MV = PT$ (transactions version) or $MV = PY$ (income version), where $M$ = money supply, $V$ = velocity of circulation, $P$ = price level, $T$/$Y$ = volume of transactions/real output. With $V$ and $T$ constant, $P$ varies proportionally with $M$: $\%\Delta M \approx \%\Delta P$.
\item \cle{Price index}: $\mathrm{RPI} = \dfrac{\sum p_1 q_0}{\sum p_0 q_0}\times 100$ (base-weighted / Laspeyres form).
\item \cle{Inflation rate}: $\dfrac{\mathrm{RPI}_1 - \mathrm{RPI}_0}{\mathrm{RPI}_0}\times 100$.
\item \cle{Real value of money}: $\text{Real income} = \dfrac{\text{Money income}}{\mathrm{RPI}}\times 100$.
\item \cle{Credit creation multiplier}: $\text{Total deposits} = \dfrac{\text{Initial deposit}}{\text{cash reserve ratio}}$; e.g. with a 10\% ratio, FCFA 100\,000 creates up to FCFA 1\,000\,000 of deposits in the banking system.
\item \cle{Balance sheet rule}: Total assets = Total liabilities (capital + deposits + other liabilities).
\end{itemize}

\textbf{3. Methods}\par
\begin{itemize}
\item To compute an RPI: (i) collect base-year and current-year prices; (ii) weight each item by its base-year quantity or expenditure share; (iii) sum weighted prices for both years; (iv) divide and multiply by 100.
\item To apply the quantity theory: identify $M$, $V$, $P$, $T$; state the assumptions ($V$, $T$ constant); solve for the unknown; interpret in words.
\item To compute credit creation: find the reserve ratio $r$; multiplier $= 1/r$; multiply by the initial new deposit; remember a \emph{single} bank can lend only its excess reserves, but the \emph{system} multiplies.
\item To draw up a balance sheet: liabilities on one side (capital, reserves, deposits, borrowings), assets on the other (cash, balances at central bank, bills, investments, loans and advances, fixed assets); order assets from most liquid to least liquid.
\end{itemize}

\textbf{4. Common mistakes}\par
\begin{itemize}
\item Confusing the \emph{value} of money with its \emph{face value}; value = purchasing power, which falls when prices rise.
\item Forgetting to multiply the index ratio by 100, or using current-year weights when the question specifies base-year weights.
\item Believing one bank alone can create the full multiplied amount of credit --- only the banking system as a whole can.
\item Mixing up assets and liabilities: deposits are a \emph{liability} of the bank (owed to customers); loans are an \emph{asset}.
\item Assuming $MV=PT$ proves money is the only cause of inflation --- remember the assumptions (constant $V$ and $T$) and cost-push causes.
\item Confusing $M_1$ (narrow) with $M_3$ (broad) money.
\end{itemize}

\textbf{5. GCE tips}\par
\begin{itemize}
\item In essays, always \emph{define} money and state its functions before analysing; illustrate with Cameroonian examples (CFA franc, BEAC, mobile money such as MTN MoMo and Orange Money).
\item In calculations, show every step, quote the formula first, and end with a sentence of interpretation --- method marks are awarded.
\item For evaluation marks, discuss limitations: quantity theory ignores cost-push inflation; credit creation is limited by cash leakages, the central bank's reserve requirements and borrowers' demand.
\item Distinguish clearly the aims of commercial banks: \emph{liquidity}, \emph{profitability} and \emph{security} --- and the conflict between them when arranging assets.
\end{itemize}
\end{bilan}

\findecours

\end{document}
